\chapter{Integración de campos vectoriales}\label{chap:stokes}

Los campos vectoriales se presentaron con algunas generalidades en la sección \secref{funcionesenelespacio}. \label{integracioncamposvectoriales}
Ahora nos ocuparemos de algunos resultados relativos a la integración de estas funciones.

\section{Integrales de línea en campos vectoriales}

Algunos libros como \cite{stewart_calculo_2018} o \cite{thomas_calculo_2015} motivan la definición de integrales sobre una curva mediante conceptos físicos como el trabajo. En tal sentido, una fuerza cons\-tante ${\bf F}$ mueve una partícula sobre el segmento que une al punto $A$ con el punto $B$. Si $W$ es el trabajo realizado, entonces $W={\bf F}\cdot \overrightarrow{AB}$.

Abordaremos inicialmente el caso bidimensional.
Si una fuerza no constante está determinada por ${\bf F}(x,y)= \langle P(x,y),Q(x,y)\rangle $
y el movimiento no es rectilíneo, sino que se hace a través de una trayectoria la cual está dada
${\bf r}(t) =\left\langle f(t),g(t) \right\rangle $, con $t\in [a,b]$, siendo $f$ y $g$ funciones continuas en $[a,b]$.
En cualquier punto de la curva $C$ la fuerza está dada por ${\bf F}({\bf r}(t))$.

\begin{figure}[H]
  \centering
  \caption{Curva en $\mathbb{R}^2$}
  \includegraphics[width=0.8\textwidth]{Fig/170.pdf}
  \fuentepropia
\end{figure}

Consideremos una partición del intervalo $[a,b]$, esto es, un conjunto finito de puntos
$\{t_0, t_1, t_2, \dots, t_{n-1}, t_n\}$, con $a = t_0 < t_1 < t_2 < \cdots < t_{n-1} < t_n = b$.

Entonces, el vector $\overrightarrow{A_{i-1}A_i}$ está dado por
\[
  \overrightarrow{A_{i-1}A_i} = \mathbf{r}(t_i) - \mathbf{r}(t_{i-1}) =
  \left\langle f(t_i) - f(t_{i-1}),\; g(t_i) - g(t_{i-1}) \right\rangle\!.
\]

Por el teorema del valor medio, existen $t_i^{\ast} \in [t_{i-1}, t_i]$ y
$t_i^{\ast\ast} \in [t_{i-1}, t_i]$, tales que
\[
  \overrightarrow{A_{i-1}A_i} =
  \left\langle f'(t_i^{\ast}),\; g'(t_i^{\ast\ast}) \right\rangle \, \Delta t_i.
\]

Para cada $i$, consideremos $\mathbf{F}_i$ como una aproximación al vector fuerza
$\mathbf{F}(\mathbf{r}(t))$ sobre la curva $C$ desde el punto $A_{i-1}$ hasta $A_i$. Por lo tanto, el trabajo realizado por $\mathbf{F}$ para trasladar una partícula desde $A_{i-1}$ hasta $A_i$ está dado por:
{\small
  \begin{multline*}
    \Delta W_i =
    \left\langle
    P\left( \left\langle f(t_i^{\ast}),\; g(t_i^{\ast}) \right\rangle \right),
    Q\left( \left\langle f(t_i^{\ast\ast}),\; g(t_i^{\ast\ast}) \right\rangle \right)
    \right\rangle \cdot
    \left\langle f'(t_i^{\ast}),\; g'(t_i^{\ast\ast}) \right\rangle \, \Delta t_i
  \end{multline*}
}

Una aproximación al trabajo total para desplazar una partícula sobre la curva $C$ desde $\mathbf{r}(a)$ hasta $\mathbf{r}(b)$ está dada por:
\[
  W \approx \sum_{i=1}^{n} \Delta W_i.
\]

Si $n \to \infty$, entonces $\Delta t_i \to 0$, por lo tanto:
\begin{eqnarray}
  W &=& \lim_{n\rightarrow\infty}\sum_{i=1}^{n}\Delta W_{i} \nonumber \\
  &=& \int_{a}^{b}\left\langle P\left( \left\langle f(t) ,g\left(t\right) \right\rangle \right) ,Q\left( \left\langle f(t)
  ,g(t) \right\rangle \right) \right\rangle \cdot \left\langle f^{\prime }(t) ,g^{\prime }(t) \right\rangle dt \label{integraldelinea2}
\end{eqnarray}
Simplificando la ecuación \ref{integraldelinea2} se precisa la expresión que determina el siguiente concepto.

\begin{defi}
  Si $C$ es una curva suave en el plano parametrizada por la función ${\bf r}(t)=\langle f(t),g(t)\rangle$, con $t \in [a,b]$
  y ${\bf F}(x,y)= \langle P(x,y),Q(x,y)\rangle $ es un campo vectorial con funciones componentes continuas, se define la integral del campo vectorial ${\bf F}$ sobre la curva $C$ como
  \begin{equation} \label{integraldelinea3}
    \int_C {\bf F} \cdot d {\bf r} =\int_{a}^{b}{\bf F}({\bf r}(t)) \cdot {\bf r}^{\prime }(t) dt
  \end{equation} \index{Integrales de línea}
\end{defi}

Este concepto se puede extender a campos vectoriales tridimensionales. Para una función vectorial
{\small ${\bf r}(t)=\langle f(t),g(t),h(t)\rangle$,} con {\small $t \in [a,b]$;}
que parametriza a una curva suave $C$ y {\footnotesize ${\bf F}(x,y,z)= \langle P(x,y,z),Q(x,y,z), R(x,y,z)\rangle$} es un campo vectorial con funciones componentes continuas, entonces la integral del campo vectorial ${\bf F}$ sobre la curva $C$ está dada por
\begin{equation}\label{integraldelinea7}
  \scalebox{0.82}{%
    $\displaystyle \int_C {\bf F} \cdot d {\bf r} =\int_{a}^{b}{\bf F}(\langle f(t),g(t),h(t)\rangle) \cdot \langle f'(t),g'(t),h'(t)\rangle dt =\int_{a}^{b}{\bf F}({\bf r}(t)) \cdot {\bf r}^{\prime }(t)$
  }
\end{equation}

\begin{ejemplo} \label{integraldelinea}
  Evaluar la integral de línea de ${\bf F}\left( x,y,z\right) =\left\langle xy,yz,-x^{2}\right\rangle $ sobre la curva que une los puntos $\left(-1,1,3\right) $ con
  $\left( 1,1,1\right)$, considerando las siguientes parametrizaciones:

  \begin{itemize}
    \item ${\bf r}(t) =\left\langle t^{2},t,2-t\right\rangle $, con $t\in[-1,1]$. En este caso,\\ ${\bf F}\left({\bf r}(t) \right) \cdot {\bf r}^{\prime }(t)= \left\langle t^{3},t\left( 2-t\right) ,-t^{4}\right\rangle \cdot \left\langle 2t,1,-1\right\rangle = 3t^{4}-t^{2}+2t$.
      \begin{eqnarray*}
        \int_{C}{\bf F}\cdot d{\bf r} &=&\int_{-1}^{1}{\bf F}\left({\bf r}(t) \right) \cdot {\bf r}^{\prime }(t)
        = \int_{-1}^{1}(3t^{4}-t^{2}+2t) dt= \frac{8}{15}
      \end{eqnarray*}

    \item ${\bf r}\left( r\right) =\left\langle 2t-1,1,3-2t\right\rangle $, con $t\in [0,1];$ Para este caso ${\bf F}\left({\bf r}(t) \right) \cdot {\bf r}^{\prime }(t)= \left\langle 2t-1,3-2t,-\left( 2t-1\right) ^{2}\right\rangle \cdot \left\langle 2,0,-2\right\rangle = 8t^{2}-4t$
      \begin{eqnarray*}
        \int_{C}{\bf F}\cdot d{\bf r}&=& \int_{0}^{1}{\bf F}\left({\bf r}(t) \right) \cdot {\bf r}^{\prime }(t) dt = \int_{0}^{1} (8t^{2}-4t) dt = \frac{2}{3}
      \end{eqnarray*}
  \end{itemize}
  Nótese que, para el mismo campo vectorial, la curva une los mismos puntos y el valor de la integral es distinto.
\end{ejemplo}

Los cálculos anteriores se evalúan en \verb"Mathematica" ejecutando los siguien\-tes comandos:

\begin{tcolorbox}[title=\bf Instrucción \verb"Mathematica"]\begin{Verbatim}[formatcom=\letraverbatim]
r={x,y,z}={t^2,t,2-t}
Integrate[Dot[{x*y,y*z,-x^2},D[r,t]],{t,-1,1}]
r={x,y,z}={2*t-1,1,3-2*t}
Integrate[Dot[{x*y,y*z,-x^2},D[r,t]],{t,0,1}]
\end{Verbatim}
\end{tcolorbox}

\subsection{Teorema fundamental para integrales de línea}
De acuerdo con \cite{acosta_gempeler_doce_2020},
para un campo escalar $\phi$ con dominio $D \subseteq \mathbb{R}^3$ y $C$ una curva en
$D$ que se parametriza mediante una función vectorial continua ${\bf r} : [a, b]\to D$. Para calcular la circulación del campo vectorial ${\bf F}=\nabla \phi$ a lo largo de $C$, se sigue que
\begin{eqnarray*}
\int_{C}F\cdot d{\bf r}&=&\int_{C}\nabla\phi\cdot dr=\int_{a}^{b}\nabla \phi\left({\bf r}(t)\right)\cdot {\bf r}^{\prime}(t) dt\\
&=& \int_{a}^{b} \frac{d}{dt} \left(\phi \left({\bf r}(t)\right) \right) dt= \phi \left({\bf r}(b)\right) -\phi\left({\bf r}\left( a\right) \right)\!.
\end{eqnarray*}
Por lo tanto, si ${\bf F}$ es un campo gradiente, entonces la circulación de ${\bf F}$ a lo largo de toda curva dentro del dominio de ${\bf F}$ solo depende de los puntos extremos de dicha curva.

\begin{ejemplo}
Evaluar la integral de línea de {\small ${\bf F}\left( x,y,z\right) =\left\langle y^{2}-2xz,2xy,-x^{2}\right\rangle $}
sobre la curva que une los puntos $\left( -1,1,3\right) $ con $\left( 1,1,1\right)$,
considerando las siguien\-tes parametrizaciones.

\begin{itemize}
  \item ${\bf r}(t) =\left\langle t,t^{2},2-t\right\rangle $, con $t\in [-1,1]$
    \begin{eqnarray*}
      \int_{C}{\bf F}\cdot dr &=& \int_{-1}^{1}{\bf F}\left({\bf r}(t) \right) \cdot {\bf r}^{\prime }(t) dt\\
      &=& \int_{-1}^{1}\left\langle t^{4}-2t\left( 2-t\right) ,2t^{3},-t^{2}\right\rangle \cdot \left\langle 1,2t,-1\right\rangle dt \\
      &=&\int_{-1}^{1}\left(5t^{4}+3t^{2}-4t\right) dt= 4
    \end{eqnarray*}

  \item ${\bf r}(t) =\left\langle 2t-1,1,3-2t\right\rangle $, con $t\in [0,1]$
    \begin{eqnarray*}
      \int_{C}{\bf F}\cdot d{\bf r}&=&\int_{0}^{1}{\bf F}\left({\bf r}(t) \right) \cdot{\bf r}^{\prime }(t) dt \\
      %&=& \int_{0}^{1}\left\langle 1-2\left( 2t-1\right) \left( 3-2t\right) ,2\left(2t-1\right) ,-\left( 2t-1\right) ^{2}\right\rangle \cdot \left\langle 2,0,-2\right\rangle dt \\
      &=& \int_{0}^{1}\left( 24t^{2}-40t+16\right)dt= 4.
    \end{eqnarray*}
\end{itemize}

Nótese que ${\bf F}$ es conservativo con función potencial $\phi(x,y,z)=xy^2-x^2z$, con lo cual
$\int_{C}{\bf F}\cdot d{\bf r}= \phi(1,1,1)-\phi(-1,1,3)=4$.
Compárese con los resultados del ejemplo (\ref{integraldelinea}).
\end{ejemplo}

En \verb"Mathematica" se pueden evaluar los cálculos anteriores ejecutando los siguientes comandos:

\begin{tcolorbox}[title=\bf Instrucción \verb"Mathematica"]\begin{Verbatim}[formatcom=\letraverbatim]
r={x,y,z}={t,t^2,2-t}
Integrate[Dot[{y^2-2*x*z,2*x*y,-x^2},D[r,t]],{t,-1,1}]
r={x,y,z}={2*t-1,1,3-2*t}
Integrate[Dot[{y^2-2*x*z,2*x*y,-x^2},D[r,t]],{t,0,1}]
\end{Verbatim}
\end{tcolorbox}

La función potencial se obtiene en \verb"Mathematica" con el siguiente proceso.

\begin{tcolorbox}[title=\bf Instrucción \verb"Mathematica"]\begin{Verbatim}[formatcom=\letraverbatim]
{P,Q,R}={y^2-2*x*z,2*x*y,-x^2};
DSolve[{D[u[x,y,z],x]==P,D[u[x,y,z],y]==Q,
D[u[x,y,z],z]==R},u[x,y,z],x,y,z]
\end{Verbatim}
\end{tcolorbox}

\begin{ejemplo}
Evaluar la integral de línea de {\small ${\bf F}\left( x,y,z\right) =\left\langle 4 x^3 y,x^4+3 y^2 z,y^3 \right\rangle$}
sobre la curva $\alpha$ parametrizada por ${\bf r}(t)=\langle 2 t-1,t+2,9 t^2 \rangle$ con $t \in [-1,2]$.

El campo es conservativo, la función potencial es $\phi(x,y,z)=x^4 y+y^3z+C$, con lo cual
\[\int_\alpha{\bf F} \cdot d{\bf r}=\phi({\bf r}(2))-\phi({\bf r}(-1))=\phi(3,4,36)-\phi(-3,1,9)=2538.\]
\end{ejemplo}

Con las siguientes instrucciones se evalúa la integral de línea.

\begin{tcolorbox}[title=\bf Instrucción \verb"Mathematica"]\begin{Verbatim}[formatcom=\letraverbatim]
F={4x^3y,x^4+3y^2z,y^3};
r[t_]:={2t-1,2+t,3*3*t^2}
{x,y,z}=r[t];
Integrate[Dot[F,D[r[t],t]],{t,-1,2}]
\end{Verbatim}
\end{tcolorbox}

Con los siguientes comandos se halla la función potencial y se utiliza el teorema fundamental para integrales de línea.

\begin{tcolorbox}[title=\bf Instrucción \verb"Mathematica"]\begin{Verbatim}[formatcom=\letraverbatim]
Clear[x,y,z,r,t]
{P,Q,R}={4x^3y,x^4+3y^2z,y^3};
\[Phi][x_,y_,z_]=u[x,y,z]/.First@DSolve[{D[u[x,y,z],
x]==P,D[u[x,y,z],y]==Q,D[u[x,y,z],z]==R},
u[x,y,z],x,y,z]
\[Phi][3,4,36]-\[Phi][-3,1,9\]
\end{Verbatim}
\end{tcolorbox}

En la ecuación \ref{integraldelinea3} al considerar el vector tangente unitario y $s$ es la longitud de la curva, haciendo el cambio ${\bf T}(t)= \frac{{\bf r}^\prime(t)}{\Vert {\bf r}^\prime(t) \Vert} = \frac{d{\bf r}}{ds}ds$,
se obtiene una expresión equivalente a la fórmula \ref{integraldelinea7} en términos del vector tangente unitario.

\begin{figure}[H]
\centering
\caption{Vectores ${\bf F}(x,y,z)$, ${\bf T}$ y ${\bf n}$}
\includegraphics[width=0.8\textwidth]{Fig/171.png}
\fuentepropia
\end{figure}

{\small
\begin{equation} \label{integraldelinea4}
\int_C {\bf F} \cdot d{\bf r}=\int_C{\bf F}({\bf r}(t))\cdot\frac{{\bf r}^\prime(t)}{\Vert{\bf r}^\prime(t)\Vert}ds
=\int_C{\bf F}({\bf r}(t)) \cdot {\bf T}(t)ds=\int_C{\bf F}\cdot {\bf T}ds
\end{equation}
}

\subsection{Integrales de flujo y de circulación para campos de velocidad}

En la ecuación \ref{integraldelinea4}, según \cite{thomas_calculo_2015},
si ${\bf F}$ representa el campo de velocidad de un fluido que fluye a través de una región en el espacio, la integral de ${\bf F}\cdot {\bf T}$ a lo largo de una curva $C$ nos proporciona el \emph{flujo} del fluido a lo largo de la curva, o la \emph{circulación} alrededor de $C$.

\begin{defi}
Si $C$ es una curva suave parametrizada por ${\bf r}(t)$ y sea ${\bf F}(x,y)$ un campo vectorial cuyas funciones componentes son continuas. La circulación de ${\bf F}$ a lo largo de la curva $C$
es \[\int_C{\bf F} \cdot {\bf T} ds.\] \index{Circulación}
\end{defi}

En el plano $xy$ consideremos una curva cerrada simple $C$, la razón con que un fluido entra o sale de una región acotada por $C$ se determina de la siguiente manera.

\begin{defi}
Si $C$ es una curva cerrada, simple y suave en el dominio de un campo vectorial ${\bf F}(x,y) = P(x, y)\vi + Q(x, y)\vj$, cuyas funciones componentes son continuas, el flujo de {\bf F} a través de $C$ es \[\text{Flujo de ${\bf F}$ a través de $C$:}\int_C{\bf F} \cdot {\bf n} ds\]
En donde {\bf n} denota el vector normal unitario a la curva $C$ que apunta hacia afuera.
\end{defi}

Una exposición más detallada se encuentra en \cite{thomas_calculo_2015}.

Al término ${\bf F}(x(t),y(t)) \cdot {\bf T}(t)$ se le denomina la \emph{densidad de circulación} del campo ${\bf F}$ en el punto $(x(t),y(t))$ de la curva $C$. La \emph{densidad de flujo}
del campo ${\bf F}$, en el punto $(x(t),y(t))$ de la curva $C$, se define como la componente normal de ${\bf F}$, es decir,
${\bf F}(x(t),y(t)) \cdot {\bf n}(t)$. \index{Densidad!de circulación} \index{Densidad!de flujo}
De modo que integrando la densidad de circulación sobre la curva se obtiene la circulación e integrando la densidad de flujo se obtiene el flujo.

Si ${\bf T}$ y ${\bf n}$ denotan los vectores tangente y normal unitario a la curva $\alpha$
en el punto $(x(t),y(t))$ entonces,
${\bf T}(x,y)= \frac{1}{\sqrt{(x^\prime(t))^2+y^\prime(t))^2}}\la x^\prime(t), y^\prime(t)\ra$ y
${\bf n}(x,y)= \frac{1}{\sqrt{(x^\prime(t))^2+y^\prime(t))^2}}\la -y^\prime(t), x^\prime(t)\ra$.

\subsection{Densidad rotacional y densidad de expansión}
En \cite{acosta_gempeler_doce_2020}se define la densidad rotacional como la circulación por unidad de área, y la densidad de expansión es el flujo por unidad de área. Analizaremos estos conceptos en campos vectoriales bidimensionales. Posteriormente, se extienden a campos tridimensionales.

Consideremos ${\bf F} (x,y)=P(x,y)\vi+Q(x,y)\vj$ un campo de velocidad del flujo de un fluido en el plano tal que las primeras derivadas parciales de $P(x,y) $ y $Q(x,y) $ son continuas en cada punto del rectángulo $R$.

\begin{figure}[H]
\centering
\caption{La razón del flujo a través de cada lado en la dirección indicada}
\includegraphics[width=0.8\textwidth]{Fig/172.pdf}
\fuentepropia
\end{figure}

Según \cite{thomas_calculo_2015}, se desea medir la tasa en un punto del giro de una rueda con paletas, cuyo eje es perpendicular al plano, en un fluido que se mueve en una región plana. Esto nos da una idea de la forma en que circula un fluido con respecto a los ejes, localizados en diferentes puntos, perpendiculares a la región. En ocasiones, los físicos se refieren a esto como la densidad de circulación de un campo vectorial {\bf F}  en un punto.

Supongamos un rectángulo $R$ con lados paralelos a los ejes coordenados y con longitudes $\Delta x$ y $\Delta y$. Si las funciones componentes
$P(x,y) $ y $Q(x,y) $ preservan el signo en una pequeña región del plano que contiene al rectángulo $R$.
Para este procedimiento supondremos que los componentes de ${\bf F} $ son positivas.
La razón de circulación de ${\bf F} $ alrededor de la frontera de $R$ es la suma de las razones de flujo a lo largo de los lados en dirección tangencial. Para la orilla inferior, la razón de flujo es aproximadamente el producto de la densidad de circulación de ${\bf F}$ en el punto medio para la longitud del segmento, es decir,
${\bf F} (x,y)\cdot\vi\Delta x=P(x,y)\Delta x$.

Las razones de flujo son positivas o negativas, dependiendo de las componentes de
${\bf F} $. Aproximamos la razón de circulación neta alrededor de la frontera rectangular de $R$ sumando las razones de flujo a lo largo de los cuatro lados, como lo indican los siguientes productos punto.

\begin{tabular}{rl}
Arriba &  ${\bf F} \left( x,y+\Delta y\right) \cdot (-\vi) \Delta x=-P\left( x,y+\Delta y\right) \Delta x$ \\
Abajo & ${\bf F} (x,y) \cdot \vi\Delta x=P(x,y) \Delta x$ \\
Lado derecho & ${\bf F} ( x+\Delta x,y) \cdot \vj\Delta y=Q\left(x+\Delta x,y\right) \Delta y$ \\
Lado izquierdo & ${\bf F} (x,y) \cdot \left( -\vj\right) \Delta y=-Q(x,y) \Delta y$.
\end{tabular}

Al sumar estos elementos y aplicar el teorema del valor medio, se sigue que

\begin{multline}
Q( x+\Delta x,y) \Delta y-Q(x,y) \Delta y-P\left(x,y+\Delta y\right) \Delta x+P(x,y) \Delta x \\
=(Q( x+\Delta x,y) -Q(x,y)) \Delta y-\left(P\left( x,y+\Delta y\right) -P(x,y) \right) \Delta x\\
\approx \left( \frac{\partial Q}{\partial x}(x^{\ast },y)\Delta x\right) \Delta y-\left( \frac{\partial P}{\partial y}\left(x,y^{\ast }\right) \Delta y\right) \Delta x \
\end{multline}
Para algún $x^* \in (x,x+\Delta x)$ e $y^* \in (y,y+\Delta y)$.
Es por esto que la razón de circulación alrededor del rectángulo $R$ es

\[\left( \frac{\partial Q}{\partial x}(x^{\ast },y) \Delta x\right) \Delta y-\left( \frac{\partial P}{\partial y}\left( x,y^{\ast }\right) \Delta y\right) \Delta x\]

Al dividir  el área de $R$ se obtiene una estimación de la razón de circulación por unidad de área o densidad rotacional para el rectángulo

\[\frac{\left( \frac{\partial Q}{\partial x}(x^{\ast },y) \Delta x\right) \Delta y-\left( \frac{\partial P}{\partial y}\left( x,y^{\ast}\right) \Delta y\right) \Delta x}{\Delta x\Delta y}=\frac{\partial Q}{\partial x}(x^{\ast },y) -\frac{\partial P}{\partial y}\left(x,y^{\ast }\right)\!. \]

Lo anterior motiva la siguiente definición.

\begin{defi}
Sea ${\bf F} (x,y)=P(x,y) \vi+Q(x,y) \vj$ un campo vectorial cuyas funciones componentes poseen primeras derivadas parciales continuas.
{\bf La densidad rotacional} del campo vectorial ${\bf F} (x,y)$ en el punto $(x,y)$ está dada por
\begin{equation} \label{integraldelinea5}
\frac{\partial Q}{\partial x}(x,y) -\frac{\partial P}{\partial y}(x,y)
\end{equation}
Este término se denomina la componente $\vk$ del rotacional y se denota
$({\rm rot}{\bf F} )\cdot \vk$.
\index{Densidad!rotacional}
\end{defi}

De manera similar al caso anterior, supongamos que las componentes del campo preservan el signo en toda la pequeña del rectángulo $R$; se desea determinar la razón con la que el fluido sale de $R$.

\begin{figure}[H]
\centering
\caption{La razón del flujo a través de cada lado en la dirección indicada}
\includegraphics[width=0.8\textwidth]{Fig/173.pdf}
\fuentepropia
\end{figure}

De manera aproximada, la razón del flujo neto a través de la frontera del rectángulo $R$ corresponde a la suma de las razones de flujo a través de los cuatro lados de $R$. En el lado inferior, el flujo es aproximadamente el producto de la densidad de flujo de ${\bf F}$ en el punto medio por la longitud del segmento, es decir
${\bf F} (x,y) \cdot \left( -\vj\right) \Delta x$.
Por lo anterior, se consideran los siguientes productos punto para el flujo de salida

\begin{tabular}{rl}
Arriba & ${\bf F} \left( x,y+\Delta y\right) \cdot \mathbf{j}\Delta x=Q\left(x,y+\Delta y\right) \Delta x$ \\
Abajo & ${\bf F} (x,y) \cdot \left( -\vj\right) \Delta x=-Q\left(x,y\right) \Delta x$ \\
Lado derecho & ${\bf F} ( x+\Delta x,y) \cdot \vi\Delta y=P\left(x+\Delta x,y\right) \Delta y$ \\
Lado izquierdo & ${\bf F} (x,y) \cdot \left( -\vi\right) \Delta y=-P(x,y) \Delta y$
\end{tabular}

Al sumar estos elementos y aplicar el teorema del valor medio, se sigue que
\begin{multline}
P( x+\Delta x,y) \Delta y-P(x,y) \Delta y+Q\left(x,y+\Delta y\right) \Delta x-Q(x,y) \Delta x \\
= \left( P( x+\Delta x,y) -P(x,y) \right)\Delta y+\left( Q\left( x,y+\Delta y\right) -Q(x,y) \right)\Delta x \\
\approx \left( \frac{\partial P}{\partial x}(x^{\ast },y)\Delta x\right) \Delta y+\left( \frac{\partial Q}{\partial y}\left(x,y^{\ast }\right) \Delta y\right) \Delta x
\end{multline}
Para algún $x^* \in (x,x+\Delta x)$ y $y^* \in (y,y+\Delta y)$.
Por lo tanto, el flujo a través de la frontera del rectángulo $R$ es aproximadamente
\[\frac{\partial P}{\partial x}(x^*,y) \Delta x\Delta y+\frac{\partial Q}{\partial y}(x,y^*) \Delta y\Delta x.\]
Al dividir  el área del rectángulo $R$, se establece una estimación del flujo total por unidad de área o densidad de flujo para el rectángulo $R$. Lo anterior motiva la siguiente definición.

\begin{defi}
Si ${\bf F} (x,y)=P(x,y) \vi+Q(x,y) \vj$ es un campo vectorial cuyas funciones componentes poseen primeras derivadas parciales continuas.
La divergencia o \textit{densidad de expansión} del campo vectorial ${\bf F}$ en el punto $(x,y)$.
\begin{equation} \label{integraldelinea6}
{\rm div}{\bf F} =\frac{\partial P}{\partial x}(x,y) +\frac{\partial Q}{\partial y}(x,y) .
\end{equation}
\index{Densidad!de expansión} \index{Divergencia}
\end{defi}

La densidad rotacional y la densidad de expansión para campos vectoriales tridimensionales ${\bf F}(x,y,z) =P(x,y,z)\vi +Q(x,y,z)\vj+R(x,y,z)\vk$,
se pueden generalizar a partir de las ecuaciones (\ref{integraldelinea5})  y (\ref{integraldelinea6});
estos elementos corresponden al rotacional y la divergencia, respectivamente.
Estos elementos se presentaron en la sección \secref{divergenciayrotacional}.
\index{Rotacional} \index{Divergencia}

\section{Teorema de Green} \label{teoremadegreen}
En \cite{thomas_calculo_2015}, se presentan los siguientes resultados en dos formas; recordemos que ${\bf T}$ y ${\bf n}$
denotan los vectores tangente unitario y normal unitario a la curva $C$.
Una exposición completa se puede consultar en \cite{acosta_gempeler_doce_2020}.
\begin{teorema}[\bfseries Teorema de Green $-$ circulación rotacional o forma tangencial]
Sea $C$ una curva suave por tramos, cerrada y simple que encierra una región $R$ en el plano. Sea ${\bf F} (x,y)=P(x,y)\vi+Q(x,y)\vj$ un campo vectorial donde $P(x,y)$ y $Q(x,y)$ tienen primeras derivadas parciales continuas en una región abierta que contiene a $R$. Entonces, la circulación en contra de las manecillas del reloj de ${\bf F}$ alrededor de $C$ es la doble integral de $({\rm rot}{\bf F})\cdot \vk$
sobre la región $R$. \index{Región abierta} \index{Teorema!de Green}
\begin{equation}
\int_C {\bf F}\cdot{\bf T}ds=\int_C Pdx+Qdy=\iint_R \left(\frac{\partial Q}{\partial x}-\frac{\partial P}{\partial y}\right)dxdy
\end{equation}
\end{teorema}

\begin{teorema}[\bfseries Teorema de Green $-$ divergencia de flujo o forma normal]
Sea $C$ una curva suave por tramos, cerrada y simple que encierra una región R en el plano.
Sea ${\bf F} (x,y)=P(x,y)\vi+Q(x,y)\vj$ un campo vectorial donde $P(x,y)$ y $Q(x,y)$ tienen primeras derivadas parciales continuas en una región abierta que contiene a $R$. Entonces, el flujo hacia fuera de ${\bf F}$ a través de $C$ es igual a la doble integral de ${\rm div}{\bf F}$ sobre la región $R$ encerrada por $C$. \index{Región abierta}
\begin{equation}
\int_C{\bf F}\cdot{\bf n}ds=\int_C Pdy-Qdx=\iint_R \left(\frac{\partial P}{\partial x}+\frac{\partial Q}{\partial y}\right)dxdy
\end{equation}
\end{teorema}

Ilustremos lo anterior.

\begin{ejemplo}
\textls[-10]{Calcular la circulación en sentido antihorario  y el flujo hacia afuera del campo
${\bf F}(x,y)=(x^{2}y+3x)\vi+(2x^{3}-y)\vj$ y la curva $C:$ $x^{2}+2y^{2}=2$.}
\end{ejemplo}

\begin{parrafoitemize}
\item Cálculo de la circulación a través del teorema de Green con la integral doble
\[\frac{\partial Q}{\partial x}-\frac{\partial P}{\partial y}=\frac{\partial }{\partial x}\left( 2x^{3}-y\right)-
\frac{\partial }{\partial y}\left( x^{2}y+3x\right) = 5x^{2}\]
por lo tanto
\[\iint\limits_{R}\left( \frac{\partial Q}{\partial x}-\frac{\partial P}{\partial y}\right) dA=\int_{-1}^{1}\int_{-\sqrt{2-2y^{2}} }^{\sqrt{2-2y^{2}} }5x^{2}dxdy= \frac{5}{2}\sqrt{2}\pi\]
\item Cálculo de la circulación a través de la integral de línea. En este caso, la curva se parametriza como
${\bf r}(t) =\left\langle \sqrt{2}\cos t,\sin t\right\rangle $, $t\in \lbrack 0,2\pi ]$.
Por lo tanto ${\bf F}({\bf r}(t))=\left\langle \left( \sin 2t+3\sqrt{2}\right) \cos t,4\sqrt{2}\cos ^{3}t-\sin t\right\rangle$, además
\[{\bf F}({\bf r}(t))\cdot{\bf r}^{\prime}(t)=2\sqrt{2}\cos(2t)-\frac{7}{2}\sin(2t)+\frac{3}{4}\sqrt{2}\cos(4t)+\frac{5}{4}\sqrt{2}.\] Entonces la circulación está dada por
\begin{multline*} \int_{0}^{2\pi}{\bf F}\left({\bf r}(t)\right)\cdot{\bf r}^{\prime}(t) dt = \\ \int_{0}^{2\pi}\left( 2\sqrt{2}\cos 2t-\frac{7}{2}\sin 2t+\frac{3}{4}\sqrt{2}\cos 4t+\frac{5}{4}\sqrt{2}\right) dt
= \frac{5}{2}\sqrt{2}\pi
\end{multline*}
Este cálculo se efectúa en \verb"Mathematica" con los siguientes comandos:

\begin{tcolorbox}[title=\bf Instrucción \verb"Mathematica"]\begin{Verbatim}[formatcom=\letraverbatim]
r={x,y}={Sqrt[2]*Cos[t],Sin[t]};
F={x^2*y+3*x,2*x^3-y};
Integrate[Dot[F,D[r,t]],{t,0,2*Pi}]
\end{Verbatim}
\end{tcolorbox}

\item Para el flujo hacia afuera a través de la integral doble
\[\frac{\partial P}{\partial x}+\frac{\partial Q}{\partial y}=\frac{\partial }{\partial x}\left( x^{2}y+3x\right) +\frac{\partial }{\partial y}\left(2x^{3}-y\right) = 2xy+2.\]
entonces el flujo es
\[\iint\limits_{R}\left(\frac{\partial P}{\partial x}+\frac{\partial Q}{\partial y}\right)dA = \int_{-1}^{1}\int_{-\sqrt{2-2y^{2}}}^{\sqrt{2-2y^{2}} }\left( 2xy+2\right) dxdy= 2\sqrt{2}\pi.\]

\item Flujo hacia afuera a través de la integral de línea
$Pdy-Qdx=$ $\left( 10\sin t\cos ^{3}t+4\sqrt{2}\cos ^{2}t-\sqrt{2}\right) dt;$
por lo tanto,
\[\int_{C}Pdx-Qdy=\int_{0}^{2\pi }\left( 10\sin t\cos ^{3}t+4\sqrt{2}\cos ^{2}t-\sqrt{2}\right) dt= 2\sqrt{2}\pi\]
\end{parrafoitemize}

\begin{ejemplo}
\textls[-10]{Considere el campo vectorial
\[{\bf F}(x,y)=\left\langle x^{2}-\frac{x}{1+y^{2}},e^{x}+\arctan y\right\rangle\!,\]
calcule el flujo hacia afuera de ${\bf F}$ a través de la cardioide $r=a(1+\cos \theta )$, $a>0$.}
% La curva se parametriza por medio de la función ${\bf r}\left( \theta \right) =\left\langle a(1+\cos\theta )\cos \theta ,a(1+\cos \theta )\sin \theta \right\rangle\!,$ $\theta \in [0,2\pi].$
En este caso,
\[\frac{\partial P}{\partial x}+\frac{\partial Q}{\partial y}=\frac{\partial }{\partial x}\left( x^{2}-\frac{x}{1+y^{2}}\right) +\frac{\partial }{\partial y}\left( e^{x}+\arctan y\right) = 2x;\]
por lo tanto, el flujo está dado por
\[\iint\limits_{R}\left( \frac{\partial P}{\partial x}+\frac{\partial Q}{\partial y}\right) dA=\iint\limits_{R}2xdA=\int_{0}^{2\pi}\int_{0}^{a(1+\cos \theta )} 2r\cos \theta  rdrd\theta= \frac{5}{2}\pi a^{3}.\]

\end{ejemplo}
Este cálculo se efectúa en \verb"Mathematica" con los siguientes comandos:

\begin{tcolorbox}[title=\bf Instrucción \verb"Mathematica"]\begin{Verbatim}[formatcom=\letraverbatim]
{P,Q}={x^2-x/(y^2+1),Exp[x]+ArcTan[y]};
A=D[P,x]+D[Q,y];
{x,y}={r*Cos[\[Theta]],r*Sin[\[Theta]]};
Integrate[A*r,{\[Theta],0,2*Pi},
{r,0,a*(1+Cos[\[Theta]])}]
\end{Verbatim}
\end{tcolorbox}

\section{Superficies orientadas}

Según \textcite{marsden_calculo_2013}, una superficie parametrizada es una función ${\bf r}:D\subset \mathbb{R}^{2}\rightarrow \mathbb{R}^{3}$, en donde $D$ es algún dominio en $\mathbb{R}^{2};$ la superficie $\sigma $ correspondiente a la función ${\bf r}$ es su imagen $\sigma ={\bf r}\left( D\right) $, que se escribe como
\[{\bf r}\left( u,v\right) =\left\langle x\left( u,v\right) ,y\left( u,v\right),z\left( u,v\right) \right\rangle\!.\]
Si las funciones $x(u,v) $, $y(u,v) $ y $z(u,v)$ son diferenciables, se dice que ${\bf r}$ es diferenciable. Si ${\bf r}$ es diferenciable en $\left( u_{0},v_{0}\right) $, al fijar $u$, se obtiene que la función ${\bf r}\left( u_{0},v\right) $ posee como imagen una curva en $\mathbb{R}^{3}$.
Los vectores ${\bf r}_{u}$ y ${\bf r}_{v}$ son tangentes a esta curva en $\left(u_{0},v_{0}\right) ;$
de modo que ${\bf r}_{u}\times {\bf r}_{v}$ es normal a la superficie en $\left( u_{0},v_{0}\right) $.
Se dice que la superficie es \emph{suave} \index{Superficie!suave} si ${\bf r}_{u}\times {\bf r}_{v}\neq {\bf 0}$, en todo punto de $S$ (esto puede depender de la parametrización elegida).

La superficie $S$ se dice \emph{orientada} \index{Superficie!orientada} si consta de dos lados (uno exterior o positivo y el otro interior o negativo) y en todo punto, hay dos vectores normales unitarios ${\bf N}_{1}$ y ${\bf N}_{2}$, con ${\bf N}_{2}=-{\bf N}_{1}$.

Ilustraremos lo anterior con unos ejemplos. En estos casos se presentan algunas superficies con una parametrización determinada (${\bf r}(u,v)$) y unos vectores del campo vectorial ${\bf r}_{u}\times {\bf r}_{v}$
(no necesariamente normalizados), se muestran pocos de ellos para diferenciarlos entre sí.

\begin{parrafoitemize}

\item La función ${\bf r}(u,t)= \langle \cos(t)\sin(u),\sin(t)\sin(u),\cos(u) \rangle $, para $t \in [0,2\pi]$ y $u \in [0,\pi]$, parametriza la esfera unitaria, que es una superficie suave y orientada.
%

\begin{figure}[H]
\centering
\caption{Campo vectorial ${\bf r}_u \times {\bf r}_v$ sobre la esfera}
\includegraphics[width=0.8\textwidth]{Fig/174.pdf}
\fuentepropia
\end{figure}

\item En las siguientes figuras se ilustra la superficie $z=xy/4$ y el campo vectorial $\langle y/4,x/4,1 \rangle$
en algunos puntos.

\begin{figure}[H]
\centering
\caption{Campo vectorial ${\bf r}_u \times {\bf r}_v$ sobre la superficie $4z=xy$}
\includegraphics[height=4.5cm]{Fig/175.pdf}
\fuentepropia
\end{figure}

Usando la parametrización ${\bf r}(x,y)=\langle x,y,xy/4\rangle$, $x \in [-2,2]$,
$y \in [-2,2]$.

\begin{figure}[H]
\centering
\caption{Campo vectorial  ${\bf r}_u \times {\bf r}_v$}
\includegraphics[width=0.8\textwidth]{Fig/176.pdf}
\fuentepropia
\end{figure}

Es sencillo ver que ${\bf r}_x \times {\bf r}_y \neq \langle 0,0,0\rangle$, para todo $x,y$,
es decir que dicha superficie es suave, también puede probarse que esta superficie es orientada.

\item La \emph{cinta de Mobius} \index{Cinta de Mobius} es un ejemplo de una superficie no orientada.
\index{Superficie!no orientada}
Para $t \in [0,2\pi]$ y $u\in[-1,1]$, esta región se puede parametrizar en la forma
{\small
\[{\bf r}(u,t)= \langle \cos(t)(2+u\cos(t/2)/2),\sin(t)(2+u\cos(t/2)),u\sin(t/2)/2 \rangle.\]
}

\begin{figure}[H]
\centering
\caption{Cinta de Mobius}
\includegraphics[width=0.8\textwidth]{Fig/177.pdf}
\fuentepropia
\end{figure}

\item El hiperboloide de un manto con ecuación $x^2-y^2=1+z^2$ se parametriza en la forma
${\bf r}(u,v) = \langle \cosh(u)\cos(t), \cosh(u)\sin(t),\sinh(u) \rangle$, con $t \in [0,2\pi]$, $u \in \mathbb{R}$.
Puede comprobarse que esta superficie es orientada y suave. Se muestran algunos vectores del campo vectorial ${\bf r}_t \times {\bf r}_u$ que son ortogonales a la superficie. \index{Hiperboloide}

\begin{figure}[H]
\centering
\caption{Campo vectorial ${\bf r}_x \times {\bf r}_y$ sobre el hiperboloide}
\includegraphics[width=0.8\textwidth]{Fig/178.pdf}
\fuentepropia
\end{figure}

\item La ecuación $z= \frac{-1}{10}(x^2+y^2) \cos(2 \sqrt{x^2+y^2})$, representa una superficie que se parametriza por $x=u\cos(t)$, $y=u\sin(t)$ y $z=-\frac{u^2}{10}\cos(2u)$, en donde $t\in [0,2\pi]$ y
$u\in [0,\infty)$. Esta superficie está orientada y es suave para $(x,y,z) \neq (0,0,0)$.
%\clearpage

\begin{figure}[H]
\centering
\caption{Campo vectorial ${\bf r}_u \times {\bf r}_v$}
\includegraphics[width=0.8\textwidth]{Fig/179.pdf}
\fuentepropia
\end{figure}

\item \emph{La botella de Klein}, \index{Botella de Klein}: esta región es un ejemplo de superficie no orientada. \index{Superficie!no orientada}
Para $u\in [0, 2\pi]$ y $v\in [0, \pi]$. La siguiente función vectorial determina dicha región:
\begin{multline*}
\overrightarrow{r} (u,v)= \big \langle 6\cos(v)(1+\sin(v))-4(1-\cos(v)/2)\cos(u),\\
16\sin(v), 4(1-\cos(v)/2)\sin(u) \big \rangle.
\end{multline*}

\begin{figure}[H]
\centering
\caption{La botella de Klein}
\includegraphics[width=0.5\textwidth]{Fig/180.pdf}
\fuentepropia
\end{figure}

\item La porción del paraboloide con ecuación $z=1-\frac{x^2+y^2}{4} $, $z\geq 0$
se parametriza como  ${\bf r}(t,u) = \langle u\cos(t),u\sin(t),1-\frac{u^2}{4}\rangle $, $t\in [0,2\pi]$ y $u \in [0,2]$.

\begin{figure}[H]
\centering
\caption{Campo vectorial ${\bf r}_u \times {\bf r}_v$ sobre el paraboloide}
\includegraphics[width=0.8\textwidth]{Fig/181.pdf}
\fuentepropia
\end{figure}

\item La función ${\bf r}(u,t)=\langle (3+\sin(t))\cos(u),(3+\sin(t))\sin(u),\cos(t) \rangle $, con $u,t \in [0,2\pi]$,
parametriza un toro; puede probarse que esta es una superficie suave y orientada.

\begin{figure}[H]
\centering
\caption{Campo vectorial ${\bf r}_u \times {\bf r}_v$ sobre el toro}
\includegraphics[width=0.8\textwidth]{Fig/182.pdf}
\fuentepropia
\end{figure}

\end{parrafoitemize}

\section{Dos teoremas fundamentales}
Los siguientes dos resultados son importantes para la realización de los ejercicios presentados en este capítulo; el enunciado es tomado de \cite{thomas_calculo_2015}.
\begin{teorema}[\bfseries Teorema de Stokes] \index{Teorema!de Stokes}
Sea $S$ una superficie orientada suave por partes que tiene como frontera una curva suave por partes $C$. Sea ${\bf F}(x,y,z)= P(x,y,z)\vi+ Q(x,y,z)\vj + R(x,y,z)\vk$
un campo vectorial cuyos componentes tienen primeras derivadas parciales continuas sobre una región abierta que contiene a $S$. Entonces, la circulación de ${\bf F}$ alrededor de $C$ en la dirección contraria a las manecillas del reloj con respecto al vector unitario ${\bf n}$ normal a la superficie es igual a la integral de \[\int_C {\bf F} \cdot d{\bf r}=\iint_S \rot{F}\cdot {\bf n} d\sigma\]
\index{Región abierta}
\end{teorema}

La prueba se puede conseguir en \cite{apostol_calculus_1999}.
\begin{teorema}[\bfseries Teorema de la divergencia] \index{Teorema!de la divergencia}
Sea ${\bf F}(x,y,z)= P(x,y,z)\vi+ Q(x,y,z)\vj + R(x,y,z)\vk$
un campo vectorial cuyos componentes tienen primeras derivadas parciales, y sea $S$ una superficie cerrada, orientada, suave por partes. El flujo de ${\bf F}$ a través de $S$ en la dirección del campo unitario normal exterior ${\bf n}$ es igual a la integral de ${\rm div\,{\bf F}}$ sobre la región $D$ encerrada por la superficie: \index{Flujo}
\[\iint_S {\bf F}\cdot {\bf n} d\sigma = \iiint_D {\rm div\,{\bf F}} dV.\]
\end{teorema}
La prueba se puede conseguir en \cite{apostol_calculus_1999} o \cite{marsden_calculo_2013}.

\section{Ejercicios que usan el teorema de Stokes}
Ilustraremos el uso de estos resultados en algunas situaciones particulares; en ellas se comprueba el cálculo señalado de más de una forma.
\subsection{Superficie sobre un paraboloide}

Para el campo vectorial ${\bf F}(x,y,z)=\la y^3,z^3,x^3 \ra$,
evaluar {\footnotesize $\ds \iint_S {\rm rot\,}{\bf F}\cdot d{\bf S}$} sobre la superficie del paraboloide
$x=3-\left( y-2\right) ^{2}-\left(z-1\right) ^{2}$
acotada por el cilindro $\left( y-2\right) ^{2}+\left(
z-2\right) ^{2}=4$. Esta región se consideró en la página \pageref{curvas2}.

\begin{figure}[H]
\centering
\caption{Superficie sobre el paraboloide interior al cilindro}
\includegraphics[width=0.4\textwidth]{Img/70b.pdf} \hspace{5mm}
\includegraphics[width=0.4\textwidth]{Img/T_22a.pdf}

\includegraphics[width=0.4\textwidth]{Img/T_22b.pdf}
\fuentepropia
\end{figure}

La figura se obtiene con las siguientes instrucciones:

\begin{tcolorbox}[title=\bf Instrucción \verb"Mathematica"]\begin{Verbatim}[formatcom=\letraverbatim]
{y,z,x}={2+2*u*Cos[v],2+2*u*Sin[v],(y-2)^2+(z-1)^2-3};
a=ParametricPlot3D[{x,y,z},{u,0,1},{v,0,2Pi}];
b=SliceVectorPlot3D[{1,-4+2y,-2+2z},x==-3+(y-2)^2+
(z-1)^2,{x,-3,7},{y,0,4},{z,0,4.2},VectorScale->{0.08,
Scaled[0.4]},RegionFunction->Function[{x,y,z},(y-2)^2+
(z-2)^2<4],VectorStyle->Black,VectorPoints->18,
PlotStyle->Opacity[0]];
c=ParametricPlot3D[{2+4*Sin[t],2+2*Cos[t],2+2*Sin[t]},
{t,0,2*Pi},PlotStyle->{Blend[{Red,Black},1/2],
Thickness->0.01},PlotPoints->80];
Show[b,a,c]
\end{Verbatim}
\end{tcolorbox}

Con el campo vectorial ${\bf F}$ se verifica que
${\rm rot\,}{\bf F}(x,y,z)=\la -3z^2,-3 x^2,-3 y^2\ra$.
La proyección de la superficie en el plano $yz$ se parametriza mediante las expresiones
$y = 2+2u \cos (v)$, $z = 2+2u \sin (v)$. Reemplazando en la expresión que define el paraboloide se sigue que $x = -3+ (y-2)^2+(z-1)^2=-2+4u^2+4u\sin(v)$.
En resumen, para $u\in [0,1],\, v\in [0,2\pi]$, la siguiente función parametriza la superficie sobre el paraboloide con la restricción mencionada.
\[{\bf r}(u,v)=\la -2+4u^2+4u\sin(v),2+2u\cos(v),2+2u\sin(v)\ra.\]

De esta función se obtiene
${\bf r}_u \times {\bf r}_v=\la u,-16u^2\cos(v),-8u(1+2u\sin(v)) \ra$
\begin{multline*}
{\rm rot\,}{\bf F}({\bf r}(u,v)) \cdot ({\bf r}_u \times {\bf r}_v) = 48u\bigg(2(1+u\cos(v))^2(1+2u\sin(v))\\ -(1+u\sin(v))^2
+u\cos(v)(4u^2-2+4u\sin(v))^2 \bigg ).
\end{multline*}
Con lo cual,
\begin{equation} \label{stokes1}
\iint_S {\rm rot\,}{\bf F}\cdot d{\bf S}=\int_0^1\int_0^{2\pi}{\rm rot\,}{\bf F}({\bf r}(u,v))\cdot({\bf r}_u\times{\bf r}_v) dvdu = 60 \pi
\end{equation}
Este cálculo puede efectuarse de manera simplificada mediante la ejecución de las siguientes instrucciones:

\begin{tcolorbox}[title=\bf Instrucción \verb"Mathematica"]\begin{Verbatim}[formatcom=\letraverbatim]
G=Curl[{y^3,z^3,x^3},{x,y,z}];
{y,z,x}={2+2*u*Cos[v],2+2*u*Sin[v],y^2-4*y+z^2-2*z+2};
r={x,y,z};
Integrate[Dot[G,Cross[D[r,u],D[r,v]]],
{v,0,2*Pi},{u,0,1}]
\end{Verbatim}
\end{tcolorbox}

La curva de intersección del paraboloide y el cilindro se determina reduciendo las ecuaciones originales a
$x = 2z - 2$. Con las expresiones $y = 2 + 2\cos(t)$, $z = 2 + 2\sin(t)$, se sigue que la función vectorial que describe dicha curva es

{\small
\[{\bf r}(t) = \left\langle 2 + 4\sin(t),\, 2 + 2\cos(t),\, 2 + 2\sin(t) \right\rangle\!,\; t \in [0, 2\pi].\]}

Entonces,

{\small
\[
{\bf F}({\bf r}(t)) =
\left\langle
(2 + 2\cos(t))^3,\;
(2 + 2\sin(t))^3,\;
(2 + 4\sin(t))^3
\right\rangle
\]
}
y también
{\small
\[
{\bf F}({\bf r}(t)) \cdot {\bf r}'(t) =
16\cos(t)\left(2(1 + \cos(t))^3 + (1 + 2\sin(t))^3\right)
+ 16\sin(t)(1 + \sin(t))^3.
\]
}
Por lo tanto
\begin{equation} \label{stokes2}
\int_0 ^{2\pi}{\bf F}({\bf r}(t)) \cdot{\bf r}'(t) dt = 60 \pi
\end{equation}
De acuerdo con el teorema de Stokes, los resultados obtenidos en \ref{stokes1} y \ref{stokes2} son iguales.
Este valor se obtiene con la ayuda de las siguientes instrucciones:

\begin{tcolorbox}[title=\bf Instrucción \verb"Mathematica"]\begin{Verbatim}[formatcom=\letraverbatim]
r={y,z,x}={2+2*Cos[t],2+2*Sin[t],2z-2};
Integrate[Dot[{y^3,z^3,x^3},D[r,t]],{t,0,2*Pi}]
\end{Verbatim}
\end{tcolorbox}

\subsection{Superficie de un paraboloide}

Evaluar $\ds \iint {\rm rot\,}{\bf F}\cdot d{\bf S}$, si ${\bf F\,}(x,y,z)= \la y-xz,y^2,z \ra $ y $S$ es la superficie sobre el paraboloide $2z=x^2+y^2$ acotada por el plano $-x+y+z=1$.

\begin{figure}[H]
\centering
\caption{$2z\leq x^2+y^2$ \& $z \leq 1+x-y$}
\includegraphics[width=0.4\textwidth]{Img/72a.pdf} \hspace{5mm}
\includegraphics[width=0.4\textwidth]{Img/T_26a.pdf}

\includegraphics[width=0.4\textwidth]{Img/T_26b.pdf}
\fuentepropia
\end{figure}

Con las siguientes instrucciones se obtiene la figura en \verb"Mathematica":

\begin{tcolorbox}[title=\bf Instrucción \verb"Mathematica"]\begin{Verbatim}[formatcom=\letraverbatim]
a=SliceVectorPlot3D[{2*x,2*y,-2},{2*z==x^2+y^2},
{x,-1,3},{y,-3,1},{z,0,6},RegionFunction->Function[{x,
y,z},(x-1)^2+(y+1)^2<=4],VectorStyle->Black,
VectorScale->{0.1,Scaled[0.2],Automatic},
VectorPoints->15];
b=ParametricPlot3D[{1+2*Cos[t],-1+2*Sin[t],
3+2*Cos[t]-2*Sin[t]},{t,0,2Pi},PlotStyle->Red];
Show[a,b]
\end{Verbatim}
\end{tcolorbox}

Las superficies mencionadas se trabajaron en la página \pageref{curvas1}.
La curva de intersección se parametrizó por medio de la función

\[{\bf r}(t)= \la 1+ 2\cos(t), -1+ 2\sin(t), 3+2\cos(t)-2\sin(t) \ra, \, t \in [0,2 \pi].\]
Entonces, ${\bf F}({\bf r}(t))\cdot{\bf r}'(t)= 2(2\cos(t)+\cos^2(t)-\sin(t))(2\sin(t)-1)$,
por lo tanto,

\begin{equation}\label{stokes3}
\int_C {\bf F} \cdot d{\bf r}=\int_0^{2\pi}2(2\cos(t)+\cos^2(t)-\sin(t))(2\sin(t)-1) dt = -8\pi.
\end{equation}
Del campo ${\bf F}(x,y,z)$ se sigue que ${\rm rot\,}{\bf F}(x,y,z)= \la 0,-x,-1\ra$.

La región sobre el paraboloide acotada por la curva $C$, se puede parametrizar mediante las expresiones
$x=1+2u\cos(t)$,  $y =-1+2u\sin(t)$.
La expresión para $z$ se obtiene de la ecuación del paraboloide, es decir,

\[z = (x^2 + y^2)/2=1+2u^2+2u\cos(t)-2u\sin(t);\] en donde $u\in [0,1]$,
$t\in [0,2\pi]$.
Con la función

\[{\bf r}(u,t)= \la 1+2u\cos(t),-1+2u\sin(t), 1+2u^2+2u\cos(t)-2u\sin(t) \ra,\]
se sigue que

{\small
\begin{eqnarray*}
{\bf r}_u &=& \la 2\cos(t), 2\sin(t), 2(2u +\cos(t)-\sin(t))\ra \\
{\bf r}_t &=& \la -2 u \sin(t), 2u\cos(t),-2u(\cos(t) + \sin(t)) \ra\\
{\bf r}_u \times {\bf r}_t &=& \la -4u(1+2u\cos(t)),4u(1-2u\sin(t)), 4 u  \ra.
\end{eqnarray*}
}

De modo que
${\rm rot\,}{\bf F} \cdot ({\bf r}_u \times {\bf r}_t) = 8u(-1-u\cos(t)+u\sin(t)+u^2\sin(2t)) $, entonces

{\small
\begin{multline*}
\iint_S {\rm rot\,}{\bf F}\cdot d{\bf S}=
\int_0^1\int_0^{2\pi}8u(-1-u\cos(t)+u\sin(t)+u^2\sin(2t))dtdu= -8\pi.
\end{multline*}
}

Este cálculo se obtiene como sigue, y es el mismo valor obtenido en \ref{stokes3}.

\begin{tcolorbox}[title=\bf Instrucción \verb"Mathematica"]\begin{Verbatim}[formatcom=\letraverbatim]
G=Curl[{y-x*z,y^2,z};,{x,y,z}]
r={x,y,z}={1+2*u*Cos[t],-1+2*u*Sin[t],(x^2+y^2)/2};
Integrate[Dot[G,Cross[D[r,u],D[r,t]]],
{t,0,2*Pi},{u,0,1}]
\end{Verbatim}
\end{tcolorbox}

\subsection{Superficie sobre una esfera}

Sea $S$ la superficie sobre la esfera con ecuación $x^{2}+y^{2}+z^{2}=2x$, con
$x\leq \frac{3}{2}$.
%\clearpage

\begin{figure}[H]
\centering
\caption{Superficie $S$ sobre una esfera}
\includegraphics[width=0.4\textwidth]{Fig/183.pdf} \hspace{5mm}
\includegraphics[width=0.4\textwidth]{Fig/184.pdf}
\fuentepropia
\end{figure}

Consideremos ${\bf F}\left(x,y,z\right)=\left\langle -y,y-z,-x\right\rangle $. Evaluar
{\small
\[\displaystyle \iint_{S}{\rm rot}\,{\bf F}\cdot dS.\]}

La superficie $S$ es una parte de la esfera de radio $1$ con centro en el punto  $(1,0,0)$, por lo tanto, usaremos coordenadas esféricas en la forma $x=\rho \cos \phi $, $y=\rho \cos \theta \sin \phi $,
$z=\rho \sin \theta \sin \phi $ para construir una parame- \\ trización de $S$. Tendremos en cuenta que la ecuación $x^{2}+y^{2}+z^{2}=2x$ se reescribe como $\rho =2\cos \phi $,
de modo que esta superficie $S$ se parametriza como
$x=2\cos ^{2}\phi $, $y=2\cos \theta \cos \phi \sin \phi $, $z=2\cos \phi
\sin \theta \sin \phi $. O de manera más simple:

{\small
\begin{align*}
\mathbf{u}(\theta, \phi) &= \left \langle 2 \cos^2 \phi, 2 \cos \theta \cos \phi \sin \phi, 2 \cos \phi \sin \theta \sin \phi \right \rangle \\
&= \left \langle 1 + \cos (2\phi), \cos \theta \sin (2\phi), \sin \theta \sin (2\phi) \right \rangle,
\end{align*}
}
con $\theta \in \lbrack 0,2\pi ]$, $\phi \in \lbrack \pi /6,\pi /2]$.
De esta función se obtiene

{\small
\begin{align*}
\mathbf{u}_{\theta} &= \left \langle 0, -\sin 2\phi \sin \theta, \sin 2\phi \cos \theta \right \rangle \\
\mathbf{u}_{\phi} &= \left \langle -2 \sin 2\phi, 2 \cos 2\phi \cos \theta, 2 \cos 2\phi \sin \theta \right \rangle \\
\mathbf{u}_{\theta} \times \mathbf{u}_{\phi} &= \left \langle -2 \cos 2\phi \sin 2\phi, -2 \sin^2 (2\phi) \cos \theta, -2 \sin^2 (2\phi) \sin \theta \right \rangle \\
&= \left \langle -\sin (4\phi), -2 \sin^2 (2\phi) \cos \theta, -2 \sin^2 (2\phi) \sin \theta \right \rangle.
\end{align*}
}

Además,

${\rm rot}\,{\bf F}=\nabla \times \left\langle -y,y-z,-x\right\rangle =
\left\langle 1,1,1\right\rangle $
y
\begin{eqnarray*}
{\bf F}({\bf u}(\theta,\phi))\cdot ({\bf u}_{\theta }\times {\bf u}_{\phi }) &=& -\sin 4\phi -2\left( \cos \theta +\sin \theta \right) \sin^{2}\left( 2\phi \right)\!.
\end{eqnarray*}
Con lo anterior, podemos evaluar la integral

{\small
\begin{multline*}
\iint_S {\rm rot}\, {\bf F} \cdot d{\bf S} = -\int_{\pi /6}^{\pi /2}\int_{0}^{2\pi }\left( \sin 4\phi +2\left( \cos
\theta +\sin \theta \right) \sin ^{2}\left( 2\phi \right) \right) d\theta d\phi = \frac{3\pi}{4}.
\end{multline*}
}

En \Verb"Mathematica" este proceso se hace con las siguientes instrucciones:

\begin{tcolorbox}[title=\bf Instrucción \verb"Mathematica"]\begin{Verbatim}[formatcom=\letraverbatim]
r[\[Theta]_,\[Phi]_]:={1+Cos[2\[Phi]],
Cos[\[Theta]]*Sin[2\[Phi]],Sin[\[Theta]]*Sin[2\[Phi]]}
ParametricPlot3D[r[\[Theta],\[Phi]],{\[Theta],0,2*Pi},
{\[Phi],Pi/6,Pi/2}]
G=Curl[{-y,y-z,-x},{x,y,z}];
{x,y,z}={r[\[Theta],\[Phi]][[1]],r[\[Theta],
\[Phi]][[2]],r[\[Theta],\[Phi]][[3]]};
A=Simplify[Dot[G,Cross[D[r[\[Theta],\[Phi]],
\[Theta]],D[r[\[Theta],\[Phi]],\[Phi]]]]]
Integrate[A,{\[Theta],0,2*Pi},{\[Phi],Pi/6,Pi/2}]
\end{Verbatim}
\end{tcolorbox}

La figura de la mencionada superficie y el campo vectorial ${\rm rot}\, {\bf F}$ se muestran a continuación:

\begin{figure}[H]
\centering
\caption{Campo vectorial ${\bf F}$ sobre la superficie $S$}
\includegraphics[width=0.4\textwidth]{Img/T_18a.pdf} \hspace{5mm}
\includegraphics[width=0.4\textwidth]{Img/T_18b.pdf}
\fuentepropia
\end{figure}

Las figuras anteriores se obtienen ejecutando las siguientes instrucciones:

\begin{tcolorbox}[title=\bf Instrucción \verb"Mathematica"]\begin{Verbatim}[formatcom=\letraverbatim]
a=SliceVectorPlot3D[{1,1,1},{x^2+y^2+z^2==2*x},
{x,0,3/2},{y,-1,1},{z,-1,1},VectorPoints->12,
VectorStyle->Black,VectorScale->{0.12,Scaled[0.2],
Automatic},Ticks->{{0,1},{-1,0,1},{-1,0,1}},
PlotStyle->Opacity[0]];
b=ParametricPlot3D[{1+Sqrt[3]/2,Cos[\[Theta]]/2,
Sin[\[Theta]]*Sin[2*Pi/6]},{\[Theta],0,2Pi},
PlotStyle->Blue];
c=ParametricPlot3D[{1+Cos[2*\[Phi]],Cos[\[Theta]]*
Sin[2*\[Phi]],Sin[\[Theta]]*Sin[2*\[Phi]]},
{\[Theta],0,2Pi},{\[Phi],Pi/6,Pi/2},Mesh->False];
Show[a,b,c]
\end{Verbatim}
\end{tcolorbox}

La frontera de la superficie $S$ es la circunferencia $C$ determinada por la ecuación
$\vec{\alpha}(t) =\left\langle \frac{3}{2},\frac{\sqrt{3}}{2}\cos t,\frac{\sqrt{3}}{2}\sin t\right\rangle $,
$t\in \lbrack 0,2\pi ];$ entonces,
\begin{eqnarray*}
{\bf F}(\ve{\alpha}(t))\cdot d\ve{\alpha}
& =& \frac{3}{4}\left ( \sin ^{2}t-\cos t\sin t-\sqrt{3}\cos t \right );
\end{eqnarray*}
entonces
\[\int_{C}{\bf F}\cdot d\alpha=\int_{0}^{2\pi}\frac{3}{4}\left(\sin^{2}t-\cos t\sin t-\sqrt{3}\cos t\right)dt=\frac{3}{4}\pi. \]
El cálculo anterior se obtiene al ejecutar las siguientes instrucciones:

%

\begin{tcolorbox}
\begin{Verbatim}[formatcom=\letraverbatim]
{x,y,z}={3/2,Sqrt[3]/2*Cos[t],Sqrt[3]/2*Sin[t]};
F[x_,y_,z_]:={-y,y-z,-x}
Dot[F[x,y,z],D[{x,y,z},t]]
Integrate[Dot[F[x,y,z],D[{x,y,z},t]],{t,0,2*Pi}]
\end{Verbatim}
\end{tcolorbox}

\subsection{Superficie sobre un elipsoide} \index{Elipsoide}

Evaluar $\ds \iint_{S}{\rm rot}\,{\bf F}\cdot d{\bf S}$, en donde $S$ es la parte del elipsoide
$x^{2}+4y^{2}+z^{2}=16$ con $z \geq 0$ y ${\bf F}\left(x,y,z\right) =e^{x}\vi+\cos \left( yz\right)\vj+z\vk$.
%\clearpage

\begin{figure}[H]
\centering
\caption{Superficie sobre el elipsoide}
\includegraphics[width=0.8\textwidth]{Fig/185.pdf}
\fuentepropia
\end{figure}

Utilizaremos la parametrización
{\small
${\bf r}\left( u,v\right) =\left\langle 4\cos u\sin v,2\sin u\sin v,4\cos v\right\rangle $,} con
{\small $u\in \lbrack 0,2\pi ]$, $v\in \lbrack 0,\pi/2]$.}

Entonces,
\begin{eqnarray*}
{\bf r}_{u}\left( u,v\right) &=& \left\langle -4\sin u\sin v,2\cos u\sin v,0\right\rangle \\
{\bf r}_{v}\left( u,v\right) &=&\left\langle 4\cos u\cos v,2\cos v\sin u,-4\sin v\right\rangle \\
{\bf r}_{u}\times {\bf r}_{v}&=&\left\langle -8\cos u\sin ^{2}v,-16\sin u\sin ^{2}v,-8\cos v\sin v\right\rangle\!.
\end{eqnarray*}
Con los elementos anteriores es sencillo ver que ${\rm rot}\,{\bf F}= \left\langle y\sin yz,0,0\right\rangle$, además
${\rm rot}\,{\bf F}({\bf r}(u,v)) \cdot ({\bf r}_{u}\times {\bf F}_{v})= -8\sin(2u)\sin^3(v)\sin(4\sin(u)\sin(2v)), $
por lo tanto,
{\footnotesize
\begin{multline*}
\iint_{S}{\rm rot}\,{\bf F}\cdot d{\bf S}=
\int_{0}^{\pi /2}\int_{0}^{2\pi}\left( -8\sin(2u)\sin^3(v)\sin(4\sin(u)\sin(2v))\right)dudv= 0.
\end{multline*}
}

Con las siguientes instrucciones se obtienen, además de la figura, los cálculos anteriores.

\begin{tcolorbox}[title=\bf Instrucción \verb"Mathematica"]\begin{Verbatim}[formatcom=\letraverbatim]
Clear[x,y,z,F,G,r]
F[x_,y_,z_]:={Exp[x],Cos[y*z],z}
G=Curl[F[x,y,z],{x,y,z}]
r[u_,v_]:={4*Cos[u]*Sin[v],2*Sin[u]*Sin[v],4*Cos[v]}
ParametricPlot3D[r[u,v],{u,0,2*Pi},{v,0,Pi/2}]
{x,y,z}=r[u,v];
A=Simplify[Dot[G,Cross[D[r[u,v],u],D[r[u,v],v]]]]
NIntegrate[A,{u,0,2*Pi},{v,0,Pi/2}]
\end{Verbatim}
\end{tcolorbox}

La frontera de $S$ es la curva $C$ determinada por la función vectorial
$\ve{\alpha}(t) =\left\langle 4\cos t,2\sin t,0\right\rangle$, $t\in \lbrack 0,2\pi ]$.
El mismo cálculo se puede evaluar usando una integral de línea, en este caso
${\bf F}(\ve{\alpha}(t))\cdot \ve{\alpha}^{\prime}(t)=\la e^{4\cos t},\cos(0),0\ra\cdot\la -4\sin t,2\cos t,0\ra$,
por lo tanto,

{\small
\[\int_{C}{\bf F}\cdot d\ve{\alpha} = \int_{0}^{2\pi}{\bf F}(\ve{\alpha}(t))\cdot
\ve{\alpha}^{\prime }\left( t\right) dt = \int_{0}^{2\pi}\left( 2\cos t-4e^{4\cos t}\sin t\right) dt= 0.\]}

%

\begin{tcolorbox}[title=\bf Instrucción \verb"Mathematica"]\begin{Verbatim}[formatcom=\letraverbatim]
{x,y,z}={4*Cos[t],2*Sin[t],0};
F[x_,y_,z_]:={Exp[x],Cos[y*z],z}
Integrate[Dot[F[x,y,z],D[{x,y,z},t]],{t,0,2*Pi}]
\end{Verbatim}
\end{tcolorbox}

\subsection{Flujo sobre una superficie}

Calcular el flujo del campo vectorial ${\bf F}(x,y,z)= \la 2x,2y,z\ra$ sobre la superficie $S$
determinada por la ecuación $z=3-x^2-y^4+2y^2$.

\begin{figure}[H]
\centering
\caption{Campo vectorial ${\bf F}(x,y,z)=\la 2x,2y,z\ra$ sobre la superficie $S$}
\includegraphics[width=0.4\textwidth]{Img/T_4a.pdf} \hspace{5mm}
\includegraphics[width=0.4\textwidth]{Img/T_4b.pdf}
\fuentepropia
\end{figure}

Las siguientes instrucciones permiten obtener la figura que se muestra:

\begin{tcolorbox}[title=\bf Instrucción \verb"Mathematica"]\begin{Verbatim}[formatcom=\letraverbatim]
a=ContourPlot3D[3-x^2-y^4+2*y^2==z,{x,-2,2},
{y,-2,2},{z,0,4.2},Mesh->False,PlotPoints->60];
Show[a,SliceVectorPlot3D[{2*x,2*y,z},{3-x^2-y^4+
2*y^2==z},{x,-2,2},{y,-2,2},{z,0,4.4},
VectorStyle->Black,VectorScale->{0.12,Scaled[0.3],
Automatic},VectorPoints->12,PlotStyle->Opacity[0]]]
\end{Verbatim}
\end{tcolorbox}

La ecuación $z=3-x^{2}-y^{4}+2y^{2}$ determina la superficie $S$
que intersecta al plano $xy$, en una curva caracterizada por la expresión
$4=x^{2}+\left( y^{2}-1\right) ^{2}$. Esta región se muestra enseguida:

\begin{figure}[H]
\centering
\caption{$z=3-x^2-y^4+2y^2$ y su proyección en el plano $xy$}
\includegraphics[width=0.4\textwidth]{Fig/186.pdf} \hspace{5mm}
\includegraphics[width=0.4\textwidth]{Fig/187.pdf}
\fuentepropia
\end{figure}

Es claro que si $x=0$ entonces {\footnotesize $y=\pm \sqrt{3}$,} esta región se puede parametrizar formando segmentos rectilíneos entre dos puntos de la forma $A$ y $B$, que son de forma {\footnotesize $A=\left( -\sqrt{3-t^{4}+2t^{2}},t\right)$,} {\footnotesize $B=\left( \sqrt{3-t^{4}+2t^{2}},t\right)$.}
Una combinación convexa entre $A$ y $B$, permite parametrizar la proyección de la superficie sobre el plano $xy$,
esto es {\footnotesize $\left( 1-u\right)A+uB=\left(\left(2u-1\right)\sqrt{2t^{2}-t^{4}+3},t\right)$.}
La coordenada $z$ sobre la superficie, se obtiene al reemplazar $x=(2u-1)\sqrt{2t^{2}-t^{4}+3}$ y $y=t$, en
$z=3-x^{2}-y^{4}+2y^{2}$ obteniendo $z=4u(1-u)(3-t^{4}+2t^{2})$.
Por lo anterior, la superficie se describe por medio de la función
\[{\bf r}(t,u) =\la (2u-1) \sqrt{2t^{2}-t^{4}+3},t,4u(1-u)\left(3-t^{4}+2t^{2}\right)\ra,\]
con $t\in \lbrack -\sqrt{3},\sqrt{3}]$, $u\in \lbrack 0,1]$.
Con lo cual,
\[{\bf F}\left({\bf r}\left(t,u\right) \right) =\left\langle 2\left( 2u-1\right) \sqrt{2t^{2}-t^{4}+3},2t,4u\left( 1-u\right) \left( 3-t^{4}+2t^{2}\right)\right\rangle\!.\]
De la función $\bf  {r}$ obtenemos los siguientes vectores:
\begin{eqnarray*}
{\bf r}_{t}&=& \left\langle -\frac{1}{2}\left( 4t^{3}-4t\right) \frac{2u-1}{\sqrt{3-t^{4}+2t^{2}}},1,4u\left( 4t^{3}-4t\right) \left( u-1\right) \right\rangle \\
{\bf r}_{u}&=& \left\langle 2\sqrt{2t^{2}-t^{4}+3},0,4\left( t^{2}-3\right)
\left( t^{2}+1\right) \left( 2u-1\right)  \right\rangle\!.
\end{eqnarray*}
Para tener la orientación de la superficie hacia afuera se hace el producto
{\footnotesize
\[
{\bf r}_{u} \times {\bf r}_{t} =
\left\langle
4(1 - 2u)(t^{2} + 1)(t^{2} - 3),\;
\frac{-8t(t^6 - 3t^4 - t^2 + 3)}{\sqrt{3 - t^{4} + 2t^{2}}},\;
2\sqrt{3 - t^{4} + 2t^{2}}
\right\rangle\!.
\]
}

El término ${\bf F}\left({\bf r}(u,v)\right) \cdot \left({\bf r}_{u} \times {\bf r}_{v}\right)$ se puede simplificar de la siguiente mane\-ra:

{\small
\begin{multline*}
8\sqrt{2t^{2} - t^{4} + 3} \left(
3 + t^2(t^2u + t^2 - 6u) + {} \right.
\left. 3u^2(2t^{2} - t^{4} + 3) - 9u
\right),
\end{multline*}
}

de modo que el flujo de ${\bf F}$ sobre $S$ está dado por:

\[
\int_{-\sqrt{3}}^{\sqrt{3}} \int_{0}^{1}
{\\bf F}\left({\bf r}(u,v)\right) \cdot \left({\bf r}_{u} \times {\bf r}_{v}\right)
\, du\, dt \approx 136.205.
\]

Los cálculos anteriores se pueden obtener en \verb|Mathematica| con las siguientes instrucciones; con estos, además de hallar la expresión que define la proyección en el plano $xy$, se construye la parametrización mostrada antes.

\begin{tcolorbox}[title=\bf Instrucción \verb"Mathematica"]\begin{Verbatim}[formatcom=\letraverbatim]
Solve[3-x^2-y^4+2y^2==0,x]
{x1,x2,y}={-Sqrt[3+2t^2-t^4],Sqrt[3+2t^2-t^4],t};
x=(1-u)*x1+u*x2;
z=Factor[3-x^2-t^4+2t^2];
r={x,y,z};
A=Factor[Dot[{2*x,2*y,z},Cross[D[r,u],D[r,t]]]];
NIntegrate[A,{u,0,1},{t,-Sqrt[3],Sqrt[3]}]
\end{Verbatim}
\end{tcolorbox}

\subsection{Una región acotada por dos cilindros parabólicos}

Evalúe $\ds \iint_S \rot{F} \cdot d{\bf S}$ en donde ${\bf F}(x,y,z)=\left\langle y^{2},yz-x,xz-y\right\rangle $ y %% $
$S$ es la superficie acotada por las figuras de las expresiones $x=z^2$ y $x=4y-y^2$.
Esta superficie se presenta en la página \pageref{curvas3}.

En la figura~\ref{curvas3} se muestra el sólido y la representación de los campos vectoriales $\la -1,0,-2z\ra$ y  $\la 1,-4+2y,0\ra$, que determinan la orientación hacia afuera de la superficie.

\begin{figure}[H]
\centering
\caption{$x = z^2$ con $(y-2)^2+z^2 \leq 4$}
\label{curvas3}
\includegraphics[width=0.4\textwidth]{Img/53e.pdf} \hspace{5mm}
\includegraphics[width=0.4\textwidth]{Img/53h.pdf}

\includegraphics[width=0.4\textwidth]{Img/53g.pdf}
\fuentepropia
\end{figure}

Estas figuras se consiguen ejecutando las siguientes instrucciones:

\begin{tcolorbox}[title=\bf Instrucción \verb"Mathematica"]\begin{Verbatim}[formatcom=\letraverbatim]
a=ParametricPlot3D[{4*Sin[t]^2,2+2*Cos[t],2*Sin[t]},
{t,0,2*Pi},PlotStyle->{Black,Thickness[0.01]}];
b=SliceVectorPlot3D[{-1,0,2*z},{x==z^2},{x,0,4.5},
{y,0,4},{z,-2.5,2.5},RegionFunction->Function[{x,y,z},
(y-2)^2+z^2<=4],VectorScale->{0.14,Scaled[0.2],
Automatic},VectorPoints->18,VectorStyle->Black];
Show[b,a]
c=ParametricPlot3D[{4*Sin[t]^2,2+2*Cos[t],2*Sin[t]},
{t,0,2*Pi},PlotStyle->{Black,Thickness[0.01]}];
d=SliceVectorPlot3D[{1,-4+2*y,0},{x==4*y-y^2},
{x,-0.5,4},{y,-0.5,5},{z,-2,2},
RegionFunction->Function[{x,y,z},(y-2)^2+z^2<=4],
VectorScale->{0.2,Scaled[0.2],Automatic},
VectorStyle->Black,VectorPoints->20];
Show[d,c]
\end{Verbatim}
\end{tcolorbox}

Las ecuaciones $x=z^{2}$ y $x=4y-y^{2}$ representan dos cilindros parabólicos, su intersección proyectada en el plano $yz$ está dada por la ecuación $\left( y-2\right) ^{2}+z^{2}=4$, la cual se parametriza en la forma $y=2+2u\cos v$, $z=2u\sin v$, con $u\in \lbrack 0,1]$, $v\in \lbrack 0,2\pi ]$. Haremos el cálculo sobre cada región por separado.

\subsubsection{Integración sobre el cilindro \texorpdfstring{$x=z^{2}$}{x=z^2}}

Del campo vectorial se  sigue que $\rot{F}= \left\langle -y-1,-z,-2y-1\right\rangle $.
Además, la superficie en mención se parametriza por medio de la función
\[{\bf r}\left( u,v\right) =\left\langle 4u^{2}\sin ^{2}v,2+2u\cos v,2u\sin v\right\rangle\!,\, u\in \lbrack 0,1],\,v\in \lbrack 0,2\pi ].\] Por lo tanto,
$\rot{F}\left({\bf r}\left(u,v\right)\right) = \left\langle -2u\cos v-3,-2u\sin v,-4u\cos v-5\right\rangle $.
De la función ${\bf r}$ se obtienen los siguientes vectores

\begin{align*}
{\bf r}_{u} &= \left\langle 8u\sin ^{2}v,2\cos v,2\sin v\right\rangle \\
{\bf r}_{v} &=\left\langle 8u^{2}\cos v\sin v,-2u\sin v,2u\cos v\right\rangle \\
{\bf r}_{u}\times {\bf r}_{v} &=\left\langle 4u,0,-16u^{2}\sin v\right\rangle\!.
\end{align*}
Además, $\rot{F}\left({\bf r}(u,v)\right)\cdot\left({\bf r}_{u}\times{\bf r}_{v}\right)=-4 u(4u^2\cos(2v)-4u^2+4 u \cos (v)+5)$, entonces,
\begin{equation}\label{stokes6}
\int_{0}^{2\pi }\int_{0}^{1}\left( -4 u(4u^2\cos(2v)-4u^2+4 u \cos (v)+5)\right) dudv= -12\pi
\end{equation}
Las instrucciones requeridas en \verb"Mathematica" se presentan a continuación:

\begin{tcolorbox}[title=\bf Instrucción \verb"Mathematica"]\begin{Verbatim}[formatcom=\letraverbatim]
G=Curl[{y^2,y*z-x,x*z-y},{x,y,z}]
r={y,z,x}={2+2*u*Cos[v],2*u*Sin[v],z^2}
Integrate[Dot[G,Cross[D[r,u],D[r,v]]],
{u,0,1},{v,0,2Pi}]
\end{Verbatim}
\end{tcolorbox}

\subsubsection{Integración sobre el cilindro \texorpdfstring{$x=4y-y^{2}$}{x=4y-y^2}}

Sobre esta región se satisface que
$y= (2+2u\cos v)$ y $z=2u\sin v$ por lo tanto $x=4-2u^{2}\left( 1+\cos 2v\right);$
por lo tanto, la superficie se parametriza utilizando la función
\[{\bf r}\left( u,v\right) =\left\langle 4-2u^{2}\left( 1+\cos 2v\right) ,2+2u\cos v,2u\sin v\right\rangle\!,\]
con $u\in[0,1],\, v\in [0,2\pi]$.
Por lo tanto, $\rot{F}\left({\bf r}\left( u,v\right)\right) =\big \langle -2u\cos v-3,-2u\sin v,-4u\cos v-5\big \rangle$.
De la función ${\bf r}$ se sigue que
\begin{eqnarray*}
{\bf r}_{u} &=& \left\langle -4u\left( \cos 2v+1\right) ,2\cos v,2\sin v\right\rangle \\
{\bf r}_{v}&=& \left\langle 4u^{2}\sin 2v,-2u\sin v,2u\cos v\right\rangle \\
{\bf r}_{u}\times {\bf r}_{v}&=& \left\langle 4u,16u^{2}\cos v,0\right\rangle;
\end{eqnarray*}
además \[\rot{F}\left({\bf r}(u,v) \right)\cdot\left({\bf r}_{u}\times {\bf r}_{v}\right)= -4u\left( 4u^{2}\sin 2v+2u\cos v+3\right),\]
entonces,
\begin{equation}\label{stokes7}
\int_{0}^{2\pi }\int_{0}^{1}-4u\left( 4u^{2}\sin 2v+2u\cos v+3\right)dudv= -12\pi
\end{equation}
Se presentan las instrucciones que se requieren en \verb"Mathematica".

\begin{tcolorbox}[title=\bf Instrucción \verb"Mathematica"]\begin{Verbatim}[formatcom=\letraverbatim]
F={y^2,y*z-x,x*z-y};
G=Curl[{y^2,y*z-x,x*z-y},{x,y,z}];
{y,z,x}={2+2*u*Cos[v],2*u*Sin[v],4*y-y^2}
Integrate[Dot[G,Cross[D[{x,y,z},u],D[{x,y,z},v]]],
{u,0,1},{v,0,2Pi}]
\end{Verbatim}
\end{tcolorbox}

\subsubsection{Integración sobre la curva de intersección}

En la página \pageref{curvas3}, la curva de intersección se parametrizó  por medio de la función ${\bf r}(t)=\left\langle 4\sin ^{2}t,2+2\cos t,2\sin t\right\rangle$, $t \in [0,2\pi]$.
De esta expresión se sigue que
${\bf r}^{\prime}(t)  = \left\langle 8\cos t\sin t,-2\sin t,2\cos t\right\rangle$. Además,
{\small
\[{\bf F}\left({\bf r}(t) \right) =  \big \langle 4(\cos t+1) ^{2},4(\cos t+1)(\cos t+\sin t-1), 2(3\sin t-\cos t-\sin 3t-1) \big \rangle.\]
}

Por lo tanto,
{\small
\[{\bf F}\left({\bf r}\left(t\right) \right) \cdot {\bf r}^{\prime}(t) = -8\cos ^{2}\left( \frac{t}{2}\right) (1+\cos(t)-\cos(2t)-3\sin(t)-5\sin(2t)-\sin(3t));\]
}
entonces,
\begin{equation}\label{stokes8}
\int_{0}^{2\pi}{\bf F}\left({\bf r}\left(t\right) \right) \cdot {\bf r}^{\prime}(t)dt
= -12\pi
\end{equation}
De acuerdo con el teorema de Stokes, se buscará la igualdad de los planteamientos que conducen a los resultados obtenidos en \ref{stokes6}, \ref{stokes7} y \ref{stokes8}.
En \verb"Mathematica" esta integral se evalúa usando las siguientes instrucciones:

\begin{tcolorbox}[title=\bf Instrucción \verb"Mathematica"]\begin{Verbatim}[formatcom=\letraverbatim]
F={y^2,y*z-x,x*z-y};
{y,z,x}={2+2*Cos[t],2*Sin[t],z^2};
Integrate[Simplify[Dot[F,D[{x,y,z},t]]],{t,0,2*Pi}]
\end{Verbatim}
\end{tcolorbox}

\index{Flujo}
En las siguientes secciones se calcula el flujo de un campo vectorial sobre ciertas superficies.
Este valor se halla por diversos métodos.

\subsection{Una superficie de tipo exponencial}

Consideremos el campo vectorial ${\bf F}(x,y,z) = \la x-y,x+y,2z \ra$, comprobar el teorema de Stokes sobre la superficie con ecuación $z=(x+y^2)e^{1-x^2-y^2}$ para $x^2+y^2\leq 4$.
%\clearpage
%\vspace*{-1.6em}

\begin{figure}[H]
\centering
\caption{$z=(x+y^2)e^{1-x^2-y^2}$}
\includegraphics[width=0.4\textwidth]{Img/T_33a.pdf} \hspace{5mm}
\includegraphics[width=0.4\textwidth]{Img/T_33c.pdf}

\includegraphics[width=0.4\textwidth]{Img/T_33b.pdf}
\fuentepropia
\end{figure}

Con las siguientes orientaciones se consiguen estas figuras; en ellos se hace uso de la instrucción:
\verb"ListSliceVectorPlot3D". \index{ListSliceVectorPlot3D}

\begin{tcolorbox}[title=\bf Instrucción \verb"Mathematica"]\begin{Verbatim}[formatcom=\letraverbatim]
datos=Table[{x-y,x+y,2*z},{z,-2,2},{y,-2,2},{x,-2,2}];
a=ContourPlot3D[z==(x+y^2)*Exp[1-x^2-y^2],{x,-2,2},
{y,-2,2},{z,-4/3,4/3},MeshStyle->Gray,
BoxRatios->{1,1,0.75}];
b=ListSliceVectorPlot3D[datos,ImplicitRegion[z==(x+
y^2)*Exp[1-x^2-y^2],{x,y,z}],DataRange->{{-2,2},
{-2,2},{-4/3,3/2}},VectorPoints->15,VectorStyle
->{Thickness->0.005,Black,Thickness->0.004},
PlotPoints->60,PlotStyle->{Opacity[0]}];
c=ParametricPlot3D[{2*Cos[t],2*Sin[t],2*Exp[-3]
*(Cos[t]+2Sin[t]^2)},{t,0,2Pi},PlotStyle->Red];
Show[a,b,PlotRange->{{-2,2},{-2,2},{-4/3,3/2}}]
\end{Verbatim}
\end{tcolorbox}

En el plano $xy$, la región determinada por la desigualdad $x^2+y^2\leq 4$ se describe por medio de las expresiones
$x=2u\cos(t)$, $y=2u\sin(t)$, con lo cual \[z=(x+y^2)e^{1-x^2-y^2}=2ue^{1-4u^2}(\cos(t)+2u\sin^2(t)).\]
Por lo tanto, la función
\[{\bf r}(t,u) = \la 2u\cos(t),2u\sin(t) ,2ue^{1-4u^2}(\cos(t)+2u\sin^2(t)))\ra,\]
con $t\in [0,2\pi]$, $u\in [0,1]$,  parametriza la superficie en mención.
Puede verse que
\begin{multline*}
{\bf r}_u \times {\bf r}_t = \bigg \langle 4 e^{1-4 u^2} u \left(4 u^3 \cos (t)-4 u^3
\cos (3 t)+4 u^2 \cos (2 t)+4 u^2-1\right),\\
16e^{1-4 u^2} u^2 \sin (t) \left(2 u^2 \cos (2t)-2 u \cos (t)-2 u^2+1\right),4 u\bigg \rangle;
\end{multline*}
además,
$\rot{F}\cdot ({\bf r}_u \times {\bf r}_t)= 8u$, entonces
\begin{equation} \label{stokes9}
\iint_S \rot{F}\cdot {\bf n} d \sigma= \int_0^1\int_0^{2\pi} 8u dtdu =8\pi.
\end{equation}
Este cálculo se hace en \verb"Mathematica" con las siguientes instrucciones:

\begin{tcolorbox}[title=\bf Instrucción \verb"Mathematica"]\begin{Verbatim}[formatcom=\letraverbatim]
G=Curl[{x-y,x+y,2*z},{x,y,z}]
r={x,y,z}={2*u*Cos[t],2*u*Sin[t],
(x+y^2)*Exp[1-x^2-y^2]};
Integrate[Dot[G,Cross[D[r,u],D[r,t]]],{u,0,1},
{t,0,2Pi}]
\end{Verbatim}
\end{tcolorbox}

Para evaluar la integral de línea sobre la frontera de la superficie en mención, se puede usar la función
\[{\bf r}(t) = \la 2\cos(t),2\sin(t) ,2e^{-3}(\cos(t)+2\sin^2(t)))\ra,\]
con $t\in [0,2\pi]$ como parametrización de dicha curva. Es claro que
\[{\bf F}({\bf r}(t))= \la 2\cos(t)-2\sin(t),2\sin (t)+2\cos(t),4e^{-3}\left(2\sin ^2(t)+\cos(t)\right)\ra;\]
además,
${\bf F}({\bf r}(t)) \cdot {\bf r}^\prime (t)= 4 + 4e^{-6}(3\sin(2t)-\sin(t)+3\sin(3t)-2\sin(4t))$.

Por lo tanto,
\begin{equation} \label{stokes10}
\scalebox{0.85}{%
$\displaystyle \int_C {\bf F} \cdot d{\bf r} = \int_0^{2\pi}(4 + 4e^{-6}(3\sin(2t)-\sin(t)+3\sin(3t)-2\sin(4t))) dt = 8\pi$
}
\end{equation}

Es evidente la igualdad de los resultados obtenidos en \ref{stokes9} y \ref{stokes10}.
Las siguientes instrucciones permiten realizar este procedimiento:

\begin{tcolorbox}[title=\bf Instrucción \verb"Mathematica"]\begin{Verbatim}[formatcom=\letraverbatim]
{x,y}={2*Cos[t],2*Sin[t]};
{x,y,z}={2*Cos[t],2*Sin[t],(x+y^2)*Exp[1-x^2-y^2]};
Integrate[Dot[{x-y,x+y,2*z},D[{x,y,z},t]],{t,0,2Pi}]
\end{Verbatim}
\end{tcolorbox}

\subsection{Flujo sobre la superficie \texorpdfstring{$\bm{\mathsf{z=\ln(x^2+y^2)}}$}{z=ln(x²+y²)}}

Calcular el flujo del campo vectorial ${\bf F}(x,y,z)= \la yz,-xz,-xy \ra$ sobre la superficie determinada por las expresiones.
$z=\ln(x^2+y^2)$, con $\frac{1}{9} \leq x^2+y^2 \leq 9$.

\begin{figure}[H]
\centering
\caption{Campo vectorial ${\bf F}$ sobre la superficie $z=\ln (x^2+y^2)$}
\includegraphics[width=0.4\textwidth]{Img/T_5a.pdf} \hspace{5mm}
\includegraphics[width=0.4\textwidth]{Img/T_5b.pdf}
\fuentepropia
\end{figure}

La figura se consigue ejecutando las siguientes instrucciones:

\begin{tcolorbox}[title=\bf Instrucción \verb"Mathematica"]\begin{Verbatim}[formatcom=\letraverbatim]
SliceVectorPlot3D[{y*z,-x*z,-x*y},{z==Log[x^2+y^2]},
{x,-3,3},{y,-3,3},{z,-1,2},VectorPoints->14,
RegionFunction->Function[{x,y,z},0.1<x^2+y^2<9],
VectorStyle->Black,VectorScale->{0.16,Scaled[0.22],
Automatic}]
\end{Verbatim}
\end{tcolorbox}

La superficie en mención se parametriza mediante la función
\[{\bf r}\left( u,v\right) =\left\langle u\cos v,u\sin v,\ln u^{2}\right\rangle\!,\,\, v\in [0,2\pi], \,\, u\in \left[1/3,3\right ].\] Con esta función se hallan los siguientes vectores
{\small
\begin{align*}
{\bf r}_{u} &= \frac{\partial}{\partial u}\left \langle u\cos v, u\sin v, \ln u^2 \right \rangle = \left \langle \cos v, \sin v, \frac{2}{u} \right \rangle \\
{\bf r}_{v} &= \frac{\partial}{\partial v}\left \langle u\cos v, u\sin v, \ln u^2 \right \rangle = \left \langle -u\sin v, u\cos v, 0 \right \rangle \\
{\bf r}_{u} \times {\bf r}_{v} &= \left \langle \cos v, \sin v, \frac{2}{u} \right \rangle \times \left \langle -u\sin v, u\cos v, 0 \right \rangle = \left \langle -2\cos v, -2\sin v, u \right \rangle.
\end{align*}
}
Además, del campo vectorial considerado se obtiene
%${\bf F}\left( x,y,z\right) =\left\langle yz,-xz,-xy\right\rangle $
\begin{eqnarray*}
{\bf F}({\bf r}(u,v)) %&=&\left\langle (u\sin v) \ln u^{2},-(u\cos v) \ln u^{2},-(u\cos v)(u\sin v)\right\rangle \\
&=& \left\langle u\sin v\ln u^{2},-u\cos v\ln u^{2},-u^{2}\cos v\sin v \right\rangle
\end{eqnarray*}
y también ${\bf F}\left({\bf r}\left( u,v\right) \right) \cdot \left({\bf r}_{u}\times {\bf r}_{v}\right) = -u^{3}\cos v\sin v$,
por lo tanto, el flujo de ${\bf F}$ sobre la superficie en mención es
\[\int_{0}^{2\pi }\int_{1/3}^{3}-u^{3}\cos v\sin vdudv= 0.\]

\subsection{Flujo sobre la superficie \texorpdfstring{$z=xe^{1-x^2-y^2}$}{z=xe a la potencia de 1-x2-y2}}

Calcular el flujo del campo ${\bf F}(x,y,z) =\la x^{2},-y,1+z\ra $ sobre la superficie $S$ determinada por la expresión $z=xe^{1-x^2-y^2}$, $|x|\leq 2$, $|y|\leq 2$.

\begin{figure}[H]
\centering
\caption{Campo vectorial ${\bf F}$ sobre la superficie $z=xe^{1-x^2-y^2}$}
\includegraphics[width=0.4\textwidth]{Img/T_6b.pdf} \hspace{5mm}
\includegraphics[width=0.4\textwidth]{Img/T_6a.pdf}
\fuentepropia
\end{figure}

Con las siguientes instrucciones se obtiene la figura de la superficie en mención:

\begin{tcolorbox}[title=\bf Instrucción \verb"Mathematica"]\begin{Verbatim}[formatcom=\letraverbatim]
a=ContourPlot3D[z==x*Exp[1-x^2-y^2],{x,-2,2},{y,-2,2},
{z,-1.2,1.2},Mesh->False,ContourStyle->Orange];
Show[a,SliceVectorPlot3D[{x^2,-y,1+z},
{x*Exp[1-x^2-y^2]==z},{x,-2,2},{y,-2,2},
{z,-1.4,1.4},PlotStyle->{Opacity[0]},
VectorPoints->20,VectorScale->{0.15,Scaled[0.25],
Automatic},VectorStyle->Black],BoxRatios->{1,1,0.5}]
\end{Verbatim}
\end{tcolorbox}

La superficie que se considera se puede parametrizar de la siguiente mane\-ra
${\bf u}\left(x,y\right) =\la x,y,xe^{1-x^{2}-y^{2}}\ra $, con $x\in[-2,2]$, $y\in[-2,2]$.
De esta función se obtienen los siguientes vectores:
\begin{eqnarray*}
{\bf u}_{x} &=& = \la 1,0,\left( 1-2x^{2}\right) e^{-x^{2}-y^{2}+1}\ra \\
{\bf u}_{y} &=&  \la 0,1,-2xye^{-x^{2}-y^{2}+1}\ra \\
{\bf u}_{x}\times {\bf u}_{y} &=& \la e^{-x^{2}-y^{2}+1}(2x^{2}-1),2xye^{-x^{2}-y^{2}+1},1 \ra.
\end{eqnarray*}
También se puede verificar que
${\bf F}\left({\bf u}\left( x,y\right) \right) =\la x^{2},-y,1+xe^{1-x^{2}-y^{2}}\ra$
y
\[{\bf F}\left({\bf u}\left( x,y\right) \right) \cdot \left({\bf u}_{x}\times {\bf u}_{y}\right)
= x\left( 2x^{3}-x-2y^{2}+1\right)e^{-x^{2}-y^{2}+1}+1.\]
Entonces, el flujo de ${\bf F}$ a través de $S$ está dado por
{\small
\[\int_{-2}^{2}\int_{-2}^{2}\left( x\left( 2x^{3}-x-2y^{2}+1\right)e^{-x^{2}-y^{2}+1}+1\right) dydx \approx 22.706.\]
}

El cálculo anterior es posible realizarlo en \verb"Mathamatica" ejecutando las siguientes instrucciones:

\begin{tcolorbox}[title=\bf Instrucción \verb"Mathematica"]\begin{Verbatim}[formatcom=\letraverbatim]
u={x,y,x*Exp[1-x^2-y^2]};
NIntegrate[Dot[{x^2,-y,1+x*Exp[1-x^2-y^2]},
Cross[D[u,x],D[u,y]]],{x,-2,2},{y,-2,2}]
\end{Verbatim}
\end{tcolorbox}

\section{Cálculo de flujo utilizando el teorema de la divergencia}

Presentamos ejemplos en los que se utiliza el teorema de la divergencia y se compara con la integración sobre la superficie.

\subsection{Flujo sobre la superficie de un elipsoide} \index{Elipsoide}
Hallar el flujo del campo vectorial ${\bf F}(x,y,z)= \la xz,y,x \ra$ sobre la superficie del elipsoide con ecuación $x^2+4y^2+4z^2=4$. La siguiente figura ilustra esta situación.
\vskip2cm

\begin{figure}[H]
\centering
\caption{Campo vectorial ${\bf F}(x,y,z)= \la xz,y,x \ra$ sobre el elipsoide}
\includegraphics[width=0.4\textwidth]{Img/T_1a.pdf} \hspace{5mm}
\includegraphics[width=0.4\textwidth]{Img/T_1b.pdf}
\fuentepropia
\end{figure}

Para obtener esta figura, es necesario compilar las siguientes instrucciones:

\begin{tcolorbox}[title=\bf Instrucción \verb"Mathematica"]\begin{Verbatim}[formatcom=\letraverbatim]
SliceVectorPlot3D[{x*z,y,x},{x^2+4*y^2+4*z^2==4},
{x,-2,2},{y,-1,1},{z,-1,1},VectorPoints->12,
VectorStyle->Black,VectorScale->{0.15,
Scaled[0.2],Automatic},BoxRatios->{1,1,1/2},
Ticks->{{-2,0,2},{-2,0,2},{-1,0,1}}]
\end{Verbatim}
\end{tcolorbox}

La ecuación del elipsoide se reescribe en la forma $\frac{x^{2}}{4}+y^{2}+z^{2}=1$,
y se puede parametrizar multiplicando por $2$ la componente $x$ de la expresión correspondiente a las ecuaciones en coordenadas esféricas, que describen la esfera unitaria. Es por lo anterior que utilizaremos la función
\[{\bf r}\left( u,v\right) =\left\langle 2\cos u\sin v,\sin u\sin v,\cos
v\right\rangle\!,\,u\in \lbrack 0,2\pi ],\, v\in \lbrack 0,\pi ].\]
Al derivar parcialmente, obtenemos los siguientes vectores

\begin{eqnarray*}
{\bf r}_{u} &=& \left\langle -2\sin u\sin v,\cos u\sin v,0\right\rangle  \\
{\bf r}_{v} &=& \left\langle 2\cos u\cos v,\cos v\sin u,-\sin v\right\rangle \\
{\bf r}_{u}\times {\bf r}_{v}&=&  \left\langle \cos u\sin ^{2}v,2\sin u\sin^{2}v,2\sin v\cos v\right\rangle\!.
\end{eqnarray*}

Y con ellos obtenemos las siguientes expresiones:
{\small
\begin{align*}
{\bf F}\left( {\bf r}(u,v) \right) &=
\left\langle
2\cos u \sin v \cos v,\;
\sin u \sin v,\;
2\cos u \sin v
\right\rangle\!, \\
{\bf F}({\bf r}(u,v)) \cdot ({\bf r}_{u} \times {\bf r}_{v}) &=
2\left(
\sin^{2} u \sin v + (\cos u \cos v)(2 + \cos u \sin v)
\right) \sin^{2} v.
\end{align*}
}

Entonces, el flujo de ${\bf F}$ a través de $S$ está dado por:
\[
\int_{0}^{\pi} \int_{0}^{2\pi}
{\bf F}\left( {\bf r}(u,v) \right) \cdot \left({\bf r}_{u} \times {\bf r}_{v}\right)
\, du\, dv = \frac{8}{3}\pi.
\]

Los cálculos anteriores se pueden hacer en \verb|Mathematica| ejecutando las siguientes sentencias:

\begin{tcolorbox}[title=\bf Instrucción \verb"Mathematica"]\begin{Verbatim}[formatcom=\letraverbatim]
r={x,y,z}={2*Cos[u]*Sin[v],Sin[u]*Sin[v],Cos[v]};
A=Simplify[Dot[{x*z,y,x},Cross[D[r,v],D[r,u]]]]
Integrate[A,{u,0,2*Pi},{v,0,Pi}]
\end{Verbatim}
\end{tcolorbox}

Otra posibilidad de hacer este cálculo es mediante el uso del teorema de la divergencia.
Para este propósito se necesita calcular la divergencia del campo vectorial ${\bf F}$, esto es
\[{\rm div}{\bf F}=\nabla \cdot \langle xz,y,x\rangle = z+1.\]
En este caso, el flujo está dado por
{\small
\[\bigintss_{-2}^{2}\bigintss_{-\sqrt{1-\frac{x^{2}}{4}}}^{\sqrt{1-\frac{x^{2}}{4}}}\bigintss_{-\sqrt{1-\frac{x^{2}}{4}-y^{2}}}^{\sqrt{1-\frac{x^{2}}{4}-y^{2}}}\left( z+1\right)dzdydx= \frac{8}{3}\pi.\]
}

Este cálculo se evaluó en \verb"Mathematica" con la siguiente instrucción:

\begin{tcolorbox}[title=\bf Instrucción \verb"Mathematica"]\begin{Verbatim}[formatcom=\letraverbatim]
Integrate[z+1,{x,-2,2},{y,-Sqrt[1-x^2/4],
Sqrt[1-x^2/4]},{z,-Sqrt[1-x^2/4-y^2],
Sqrt[1-x^2/4-y^2]}]
\end{Verbatim}
\end{tcolorbox}

\subsection{Flujo sobre la superficie de un paraboloide}

Calcular el flujo del campo vectorial ${\bf F}\left( x,y,z\right) =\left\langle xz,x(y+1),z(y+1)\right\rangle$,
sobre la superficie del paraboloide $z=4-x^{2}-y^{2}$ y el disco sobre el plano $z=0$, que resulta de proyectar esta región sobre dicho plano.

\begin{figure}[H]
\centering
\caption{Campo vectorial ${\bf F}(x,y,z)= \la xz,x(y+1),z(y+1) \ra$ sobre el paraboloide}
\includegraphics[width=0.4\textwidth]{Img/T_2a.pdf} \hspace{5mm}
\includegraphics[width=0.4\textwidth]{Img/T_2b.pdf}
\fuentepropia
\end{figure}

La figura anterior se consigue con la ayuda de las siguientes instrucciones:

\begin{tcolorbox}[title=\bf Instrucción \verb"Mathematica"]\begin{Verbatim}[formatcom=\letraverbatim]
SliceVectorPlot3D[{x*z,x*(y+1),z*(y+1)},{4-x^2-
y^2==z},{x,-2,2},{y,-2,2},{z,0,4},VectorPoints->12,
VectorStyle->Black,VectorScale->{0.16,Scaled[0.22],
Automatic},Ticks->{{-2,0,2},{-2,0,2},{0,2,4}}]
\end{Verbatim}
\end{tcolorbox}

La proyección en el plano $xy$ del paraboloide corresponde a un círculo centrado en el origen del radio $2$.
Para calcular el flujo sobre el disco que forma la base del sólido, se parametriza esta superficie en la forma
\[{\bf r}\left(u,v\right) = \left\langle u\cos v,u\sin v,0\right\rangle ,\, u\in \lbrack 0,2],\,v\in \lbrack 0,2\pi ].\]
Es sencillo ver que
${\bf F}\left({\bf r}\left(u,v\right)\right)=\left\langle 0,u\left(u\sin v+1\right)\cos v,0\right\rangle;$
con la función ${\bf r}$ se obtienen los vectores:
\[{\bf r}_{u} = \left\langle \cos v, \sin v, 0 \right\rangle\] y
\[{\bf r}_{v} = \left\langle -u\sin v, u\cos v, 0 \right\rangle\!.\]
Para conseguir que la orientación sea hacia afuera, se evalúa:
\[
{\bf r}_{v} \times {\bf r}_{u} =
\left\langle -u\sin v, u\cos v, 0 \right\rangle \times
\left\langle \cos v, \sin v, 0 \right\rangle =
\left\langle 0, 0, -u \right\rangle\!.
\]

Es claro que
\[
{\\ F}\left({\bf r}(u,v)\right) \cdot \left({\bf r}_{u} \times {\bf r}_{v}\right) = 0;
\]
por lo tanto, el flujo sobre la superficie en mención es
\begin{equation} \label{stokes4}
\int_{0}^{2\pi }\int_{0}^{2}{\bf F}\left({\bf r}\left(u,v\right)\right) \cdot \left({\bf r}_{u}\times {\bf r}_{v}\right) dudv =0
\end{equation}

La superficie sobre el paraboloide se puede parametrizar en la forma
${\bf r}\left(u,v\right) =\left\langle u\cos v,u\sin v,4-u^{2}\right\rangle\!,\, u\in \lbrack 0,2],\, v\in \lbrack 0,2\pi ]$.
Con esta función hallamos los siguientes vectores:
\begin{eqnarray*}
{\bf r}_{u}&=& \left\langle \cos v,\sin v,-2u\right\rangle \\
{\bf r}_{v} &=& \left\langle -u\sin v,u\cos v,0\right\rangle \\
{\bf r}_{u}\times {\bf r}_{v} &=& \left\langle 2u^{2}\cos v,2u^{2}\sin v,u\right\rangle\!.
\end{eqnarray*}
Sustituyéndose adecuadamente, el término ${\bf F}\left({\bf r}\left(u,v\right)\right) \cdot \left({\bf r}_{u}\times {\bf r}_{v}\right)$ se simplifica como
$\big(4u-u^{3}+u^{2}((4-u^{2}) \sin v+(8u-2u^{3})\cos^{2}v +u(u\sin v+1) \sin (2v))\big)$.
Integrando esta expresión obtendremos el flujo de ${\bf F}$ a través de $S$,
\begin{equation} \label{stokes5}
\int_{0}^{2\pi }\int_{0}^{2}{\bf F}\left({\bf r}\left(u,v\right)\right) \cdot \left({\bf r}_{u}\times {\bf r}_{v}\right) dudv= \frac{56}{3}\pi
\end{equation}
Los procedimientos que conducen a la integral (\ref{stokes5}) se logran con la ejecución de las siguientes instrucciones:

\begin{tcolorbox}[title=\bf Instrucción \verb"Mathematica"]\begin{Verbatim}[formatcom=\letraverbatim]
r={x,y,z}={u*Cos[v],u*Sin[v],4-u^2};
Integrate[Dot[{x*z,x*(y+1),z*(y+1)},Cross[D[r,v],
D[r,u]]],{u,0,2*Pi},{v,0,Pi}]
\end{Verbatim}
\end{tcolorbox}

La otra alternativa para calcular el flujo es usar el teorema de la divergencia, como ${\rm div}{\bf F}=\nabla \cdot \left\langle xz,x\left( y+1\right) ,z\left( y+1\right) \right\rangle = x+y+z+1$.
Entonces, el flujo de ${\bf F}$ sobre la superficie de $S$ está dado por
\[\int_{-2}^{2}\int_{-\sqrt{4-x^{2}}}^{\sqrt{4-x^{2}}}\int_{0}^{4-x^{2}-y^{2}}\left( x+y+z+1\right) dzdydx= \frac{56}{3}\pi.\]
Este valor es el producto de sumar los resultados obtenidos en \ref{stokes4} y \ref{stokes5}.
Usando \verb"Mathematica" las siguientes sentencias realizan este cálculo.

\begin{tcolorbox}[title=\bf Instrucción \verb"Mathematica"]\begin{Verbatim}[formatcom=\letraverbatim]
F={x*z,x*(y+1),z*(y+1)};
Integrate[Div[F,{x,y,z}],{x,-2,2},{y,-Sqrt[4-x^2],
Sqrt[4-x^2]},{z,0,4-x^2-y^2}]
\end{Verbatim}
\end{tcolorbox}

\subsection{Flujo sobre la superficie de un paraboloide y un plano}

Consideremos el sólido $E$ determinado por el paraboloide con ecuación $9y=3-9x^{2}-4z^{2}+4z$ y el plano $4z-9y=9$. Calcular el flujo neto del campo vectorial
{\small ${\bf F}\left( x,y,z\right) =\la xz,x(y+1),z(y+1) \ra$,} a través de la superficie de $E$. La proyección de esta región en el plano $xz$ está dada por la ecuación  $9x^{2}+4z^{2}=12$. Esta región se puede parametrizar en la forma
$x=\frac{2}{\sqrt{3}}u\cos v$, $z=\sqrt{3}u\sin v$, en donde  $u \in [0,1]$, $v \in [0,2\pi]$.
Las siguientes figuras muestran el sólido.
%\clearpage
%

\begin{figure}[H]
\centering
\caption{Superficies $9y=3-9x^{2}-4z^{2}+4z$ y $4z-9y=9$}
\includegraphics[width=0.4\textwidth]{Img/T_T1.pdf} \hspace{5mm}
\includegraphics[width=0.4\textwidth]{Img/T_T2.pdf}
\fuentepropia
\end{figure}

Las siguientes son las instrucciones para obtener las figuras.

\begin{tcolorbox}[title=\bf Instrucción \verb"Mathematica"]\begin{Verbatim}[formatcom=\letraverbatim]
ContourPlot3D[{9*y==3-9*x^2-4*z^2+4*z,4z-9y==9},
{x,-2,2},{y,-2,1},{z,-2.5,3},Mesh->False,
ContourStyle->{Cyan,Orange}]
{x,z}={2/Sqrt[3]*u*Cos[t],Sqrt[3]*u*Sin[t]};
a=ParametricPlot3D[{x,(3-9*x^2-4*z^2+4*z)/9,z},
{u,0,1},{t,0,2Pi},PlotStyle->Orange,PlotPoints->80];
b=ParametricPlot3D[{x,(4*z-9)/9,z},{u,0,1},{t,0,2Pi},
PlotStyle->Cyan,PlotPoints->80];
Show[a,b,BoxRatios->{1,1,1}]
\end{Verbatim}
\end{tcolorbox}

La siguiente es una representación del campo vectorial que actúa sobre $E$.

\begin{figure}[H]
\centering
\caption{Campo vectorial ${\bf F}(x,y,z)= \la xz,x(y+1),z(y+1) \ra$ sobre la superficie $E$}
\includegraphics[width=0.4\textwidth]{Img/T_21a.pdf} \hspace{5mm}
\includegraphics[width=0.4\textwidth]{Img/T_21b.pdf}

\includegraphics[width=0.4\textwidth]{Img/T_21c.pdf}
\fuentepropia
\end{figure}

Esta figura se consigue con la ayuda de las siguientes instrucciones:

\begin{tcolorbox}[title=\bf Instrucción \verb"Mathematica"]\begin{Verbatim}[formatcom=\letraverbatim]
a=RegionPlot3D[(4*z-9)/9<=y<=(3-9*x^2-4*z^2+4*z)/9
&&9*x^2+4*z^2<=12,{x,-2,2},{y,-2,1},{z,-2,2},
PlotPoints->90,Mesh->False];
b=SliceVectorPlot3D[{x*z,x*(y+1),z(y+1)},
{4*z-9*y==9,9*y==3-9*x^2-4*z^2+4*z},{x,-2,2},
{y,-2,1},{z,-2,2},RegionFunction->Function[{x,y,z},
9*x^2+4*z^2<12],VectorPoints->20,VectorStyle->Black,
PlotPoints->60,VectorScale->{0.3,Scaled[0.4],
Automatic},PlotStyle->Opacity[0]]
Show[a,b]
\end{Verbatim}
\end{tcolorbox}

\subsubsection{Integral sobre la porción del plano}

Sobre el plano la coordenada $y$ está dada por
$y=\frac{4}{9}z-1=\frac{4}{9}\sqrt{3}u\sin v-1$,
de modo que una parametrización de la región sobre la porción del plano acotada por el paraboloide es la siguiente:

${\bf r}\left( u,v\right) =\left\langle \frac{2}{\sqrt{3}}u\cos v,\frac{4}{9}\sqrt{3}u\sin v-1,\sqrt{3}u\sin v\right\rangle$,
$u\in \lbrack 0,1]$, $v\in\lbrack 0,2\pi ];$

con esta función hallamos los siguientes vectores:

{\small
\begin{eqnarray*}
{\bf r}_{u} &=&\left\langle \frac{2}{3}\sqrt{3}\cos v,\frac{4}{9}\sqrt{3}\sin v,\sqrt{3}\sin v\right\rangle \\
{\bf r}_{v}&=&\left\langle -\frac{2}{3}\sqrt{3}u\sin v,\frac{4}{9}\sqrt{3}u\cos v,\sqrt{3}u\cos v\right\rangle \\
{\bf r}_{u}\times {\bf r}_{v}&=& % \left\langle \frac{2}{3}\sqrt{3}\cos v,\frac{4}{9}\sqrt{3}\sin v,\sqrt{3}\sin v\right\rangle \times \left\langle -\frac{2}{3}\sqrt{3}u\sin v,\frac{4}{9}\sqrt{3}u\cos v,\sqrt{3}u\cos v\right\rangle \\
\left\langle 0,-2u,\frac{8}{9}u\right\rangle\!.
\end{eqnarray*}
}

También pueden comprobarse las siguientes igualdades,
{\small
\begin{eqnarray*}
{\bf F}\left( {\bf r}\left( u,v\right) \right) &=& \left\langle 2u^{2}\cos v\sin v,\frac{8}{9}u^{2}\cos v\sin v,\frac{4}{3}u^{2}\sin ^{2}v \right\rangle  \\
{\bf F}\left( {\bf r}\left( u,v\right) \right) \cdot \left( {\bf r}_{u}\times {\bf r}_{v}\right)
&=& %\left\langle 2u^{2}\cos v\sin v,\frac{8}{9}u^{2}\cos v\sin v,\frac{4}{3} u^{2}\sin ^{2}v \right\rangle\cdot \left\langle 0,-2u,\frac{8}{9}u\right\rangle =
\frac{32}{27}u^{3}\sin ^{2}v-\frac{16}{9}u^{3}\sin v\cos v
\end{eqnarray*}
}

De tal modo que el flujo sobre la porción del plano es
\begin{equation} \label{T1}
\scalebox{0.88}{%
$\displaystyle \iint_{S_1} {\bf F} \cdot d\boldsymbol{\sigma} = \int_{0}^{2\pi }\int_{0}^{1}\left(\frac{32}{27}u^{3}\sin ^{2}v-\frac{16}{9}u^{3}\sin v\cos v \right) dudv = \frac{8}{27}\pi$
}
\end{equation}

El procedimiento mostrado antes se puede efectuar en \verb"Mathematica"
ejecutando las siguientes sentencias:

\begin{tcolorbox}[title=\bf Instrucción \verb"Mathematica"]\begin{Verbatim}[formatcom=\letraverbatim]
{x,y}={2/Sqrt[3]*u*Cos[v],4/9*Sqrt[3]*u*Sin[v]-1}
z=Sqrt[3]*u*Sin[v];
F={x*z,x*(y+1),z*(y+1)};
A=Dot[F,Simplify[Cross[D[{x,y,z},u],D[{x,y,z},v]]]];
Integrate[A,{u,0,1},{v,0,2Pi}]
\end{Verbatim}
\end{tcolorbox}

\subsubsection{Integral sobre la porción del paraboloide}

Con los elementos anteriores, la coordenada $y$ sobre el paraboloide está dada por
$y = \frac{1}{9}\left( 3-9x^{2}-4z^{2}+4z\right)= \frac{4}{9}\sqrt{3}u\sin v+\frac{1}{3}-\frac{4}{3}u^{2}$,
de modo que una parametrización de la región sobre el paraboloide acotada por la intersección con el plano, es la siguiente

{\footnotesize
\[{\bf r}\left( u,v\right) =\left\langle \frac{2}{\sqrt{3}}u\cos v,\frac{4}{9}\sqrt{3}u\sin v+\frac{1}{3}-\frac{4}{3}u^{2},\sqrt{3}u\sin v\right\rangle\!,u\in \lbrack 0,1],\, v\in \lbrack 0,2\pi ].\]
}

Esta función nos permite hallar los siguientes vectores:
\begin{eqnarray*}
{\bf r}_{u} &=& \left\langle \frac{2}{3}\sqrt{3}\cos v,\frac{4}{9}\sqrt{3}\sin v-\frac{8}{3}u,\sqrt{3}\sin v\right\rangle \\
{\bf r}_{v} &=& \left\langle -\frac{2}{3}\sqrt{3}u\sin v,\frac{4}{9}\sqrt{3}u\cos v,\sqrt{3}u\cos v\right\rangle \\
{\bf r}_{v}\times {\bf r}_{u} %&=&\left\langle \frac{2}{3}\sqrt{3}\cos v,\frac{4}{9}\sqrt{3} \sin v-\frac{8}{3}u,\sqrt{3}\sin v\right\rangle \times \left\langle -\frac{2}{3}\sqrt{3}u\sin v,\frac{4}{9}\sqrt{3}u\cos v,\sqrt{3}u\cos v\right\rangle \\
&=&  \left\langle \frac{8}{3}\sqrt{3}u^{2}\cos v,2u,\frac{8}{9}u+\frac{16}{9}\sqrt{3}u^{2}\sin v\right\rangle\!.
\end{eqnarray*}
Además,
\begin{multline*}
{\bf F}\left( {\bf r}\left( u,v\right) \right) = \bigg\langle 2u^{2}\sin v\cos v,
\frac{2}{3}\sqrt{3}u\left( \frac{4}{9}\sqrt{3}u\sin v+\frac{4}{3}- \frac{4}{3}u^{2}\right) \cos v, \\
\sqrt{3}u\left( \frac{4}{9}\sqrt{3}u\sin v+\frac{4}{3}-\frac{4}{3}u^{2}\right) \sin v\bigg\rangle
\end{multline*}

{\small
\begin{multline*}
{\bf F}\left({\bf r}(u,v) \right) \cdot \left( {\bf r}_{v}\times {\bf r}_{u}\right) =
\frac{16}{3}\sqrt{3}u^{4}\sin v\cos ^{2}v+ \\
\left( \frac{16}{27}\sqrt{3}u^{2}\cos v - \frac{4}{9}\sqrt{3}u\sin v \left( \frac{8}{9}u-\frac{16}{9}\sqrt{3}u^{2}\sin v\right) \right) \left( \sqrt{3}u\sin v+3-3u^{2}\right)\!.
\end{multline*}
}

De lo anterior se sigue que
\begin{equation} \label{T2}
\iint_{S_2}{\bf F}\left( {\bf r}\left( u,v\right) \right) \cdot \left( {\bf r}_{u}\times {\bf r}_{v}\right)dudv=\frac{8}{27}\pi
\end{equation}
El procedimiento anterior se puede hacer en \verb"Mathematica" con ayuda de las siguientes instrucciones:

\begin{tcolorbox}[title=\bf Instrucción \verb"Mathematica"]\begin{Verbatim}[formatcom=\letraverbatim]
{x,z}={2/Sqrt[3]*u*Cos[v],Sqrt[3]*u*Sin[v]};
y=Simplify[(3-9x^2-4*z^2+4z)/9]
{r,F}={{x,y,z},{x*z,x*(y+1),z*(y+1)}};
A=Dot[F,Simplify[Cross[D[r,v],D[r,u]]]];
Integrate[A,{u,0,1},{v,0,2Pi}]
\end{Verbatim}
\end{tcolorbox}

\subsubsection{Usando el teorema de la divergencia}

Es claro que ${\rm div}{\bf F}=z+x+y+1$, por lo tanto, el flujo de ${\bf F}$ a través de $S$ está dado por
\[\bigintsss_{-\frac{2}{\sqrt{3}}}^{\frac{2}{\sqrt{3}}}\bigintsss_{-\frac{1}{2}\sqrt{12-9x^{2}}}^{\frac{1}{2}\sqrt{12-9x^{2}}}\bigintsss_{\frac{4}{9}z-1}^{\frac{4}{9}z+ \frac{1}{3}-x^{2}-\frac{4}{9}z^{2}}\left( z+x+y+1\right) dydzdx=\frac{16}{27}\pi,\]
que coincide con la suma de las expresiones \ref{T1} y \ref{T2}.
Este cálculo se consigue ejecutando las siguientes instrucciones:

\begin{tcolorbox}[title=\bf Instrucción \verb"Mathematica"]\begin{Verbatim}[formatcom=\letraverbatim]
F={x*z,x*(y+1),z*(y+1)}
Integrate[Div[F,{x,y,z}],{x,-2/Sqrt[3],2/Sqrt[3]},
{z,-Sqrt[12-9*x^2]/2,Sqrt[12-9*x^2]/2},{y,4z/9-1,
(3+4z-9*x^2-4z^2)/9}]
\end{Verbatim}
\end{tcolorbox}

\subsection{Flujo sobre una esfera}

Se calcula el flujo del campo vectorial ${\bf F}(x,y,z)= \la y^2,-xy,z \ra$ sobre la esfera con ecuación $x^2+y^2+z^2=4$.
El gráfico se obtiene con las siguientes instrucciones:

\begin{tcolorbox}[title=\bf Instrucción \verb"Mathematica"]\begin{Verbatim}[formatcom=\letraverbatim]
SliceVectorPlot3D[{y^2,-x*y,z},{4==x^2+y^2+z^2},
{x,-2,2},{y,-2,2},{z,-2,2},VectorPoints->10,
VectorPoints->12,VectorStyle->Black,
VectorScale->{0.2,Scaled[0.2],Automatic}]
\end{Verbatim}
\end{tcolorbox}

\begin{figure}[H]
\centering
\caption{Campo vectorial ${\bf F}(x,y,z)= \la y^2,-xy,z \ra$ sobre una esfera}
\includegraphics[width=0.4\textwidth]{Img/T_3a.pdf} \hspace{5mm}
\includegraphics[width=0.4\textwidth]{Img/T_3b.pdf}
\fuentepropia
\end{figure}

La esfera puede describirse utilizando la función
\[{\bf r}\left( u,v\right)=\left\langle 2\cos u\sin v,2\sin u\sin v,2\cos v\right\rangle ,\, u\in
\lbrack 0,2\pi ],\,v\in \lbrack 0,2\pi ].\]
Con esta función hallamos los siguientes vectores:
\begin{eqnarray*}
{\bf r}_{u}&=&\left\langle -2\sin u\sin v,2\cos u\sin v,0\right\rangle \\
{\bf r}_{v}&=& \left\langle 2\cos u\cos v,2\cos v\sin u,-2\sin v\right\rangle \\
{\bf r}_{v}\times {\bf r}_{u}&=&\la 4\cos u\sin ^{2}v,4\sin u\sin ^{2}v,4\cos v\sin v \ra.
\end{eqnarray*}
Con estos elementos se puede establecer la siguiente relación:
\[
\mathbf{F}\left( \mathbf{r}\left( u,v\right) \right) = \left\langle 4\sin^{2}u\sin^{2}v,\ -4\cos u\sin u\sin^{2}v,\ 2\cos v \right\rangle
\]

Entonces,
\[
\mathbf{F}\left( \mathbf{r}\left( u,v\right) \right) \cdot \left( \mathbf{r}_{u} \times \mathbf{r}_{v} \right) = -8\cos^{2}v\sin v.
\]

Por lo tanto, el flujo de \(\mathbf{F}\) a través de la superficie de la esfera está dado por:
\[
\int_{0}^{\pi} \int_{0}^{2\pi} 8\cos^{2}v\sin v\ dudv = \frac{32}{3}\pi.
\]

El proceso anterior se puede hacer en \verb"Mathematica" utilizando las siguien\-tes instrucciones:

\begin{tcolorbox}[title=\bf Instrucción \verb"Mathematica"]\begin{Verbatim}[formatcom=\letraverbatim]
F={y^2,-x*y,z};
r={x,y,z}={2*Cos[u]*Sin[v],2*Sin[u]*Sin[v],2*Cos[v]};
Integrate[Dot[Cross[D[r,v],D[r,u]],F],{u,0,2*Pi},
{v,0,Pi}]
\end{Verbatim}
\end{tcolorbox}

El flujo también se puede hallar usando el teorema de la divergencia; en ese caso, el cálculo es el siguiente
{\small
\[\iiint_E {\rm div} \,{\bf F} dV= \int _0^{\pi}\int _0^{2\pi}\int _0^{2} (1-\rho \sin (\phi))\rho^2\sin(\phi)d\rho d\theta d\phi = \frac{32}{3} \pi.\]}
Esta integral se evaluó en \verb"Mathematica" usando la siguiente instrucción:

\begin{tcolorbox}[title=\bf Instrucción \verb"Mathematica"]\begin{Verbatim}[formatcom=\letraverbatim]
Integrate[(1-\[Rho]*Sin[\[Phi]]*Cos[\[Theta]])*
\[Rho]^2*Sin[\[Phi]],{\[Rho],0,2},{\[Phi],0,Pi},
{\[Theta],0,2*Pi}]
\end{Verbatim}
\end{tcolorbox}

\section{Cálculo de flujo sobre superficies cerradas}

En esta sección se presentan algunos ejemplos en que se calcula el flujo sobre superficies cerradas y se compara el valor obtenido usando el teorema de la divergencia, cuando es posible hacerlo.
\begin{parrafonumerado}
\item Hallar el flujo neto del campo vectorial.

${\bf F}(x,y,z) = \la y^2,x^2,z \ra $ sobre la superficie cerrada del sólido acotado por los paraboloides con ecuaciones:

%$z+3=x^2+y^2$ y $2z=6-2x^2-y^2.$

\begin{gather*}
z + 3 = x^2 + y^2 \quad \text{y} \quad 2z = 6 - 2x^2 - y^2
\end{gather*}

\begin{figure}[H]
\centering
\caption{Campo vectorial ${\bf F}(x,y,z)=\la y^2,x^2,z\ra $ sobre una superficie}
\includegraphics[width=0.4\textwidth]{Img/T_8b.pdf} \hspace{5mm}
\includegraphics[width=0.4\textwidth]{Img/T_8d.pdf}
\fuentepropia
\end{figure}

Esta representación del campo vectorial se obtiene ejecutando las siguientes instrucciones:

\begin{tcolorbox}[title=\bf Instrucción \verb"Mathematica"]\begin{Verbatim}[formatcom=\letraverbatim]
a=RegionPlot3D[z<(6-2*x^2-y^2)/2&&z>x^2+y^2-3,
{x,-2,2},{y,-2,2},{z,-3,3},PlotPoints->90,
Mesh->False];
b=SliceVectorPlot3D[{y^2,x^2,z},{2*z==6-2*x^2-y^2,
z==x^2+y^2-3},{x,-2,2},{y,-2,2},{z,-3,3},
VectorPoints->20,VectorScale->{0.15,Scaled[0.25],
Automatic},RegionFunction->Function[{x,y,z},4*x^2
+3*y^2<=12],VectorStyle->Black,PlotStyle->Opacity[0]];
Show[a,b]
\end{Verbatim}
\end{tcolorbox}

De acuerdo con el ejercicio \ref{ejercicio1} de la página \pageref{ejercicio1}, la curva de intersección corresponde a la ecuación $4x^{2}+3y^{2}=12$, que representa una elipse.
Esta región y su interior se pueden representar en el plano $xy$ por medio de las ecuaciones
$x=\sqrt{3}u\cos v$, $y=2u\sin v$, con $u\in \lbrack 0,1]$,
$v\in \lbrack 0,2\pi ]$.
Calcularemos el flujo sobre las dos superficies que limitan el sólido.

\subsubsection{Flujo sobre el paraboloide \texorpdfstring{$\bm{\mathsf{2z=6-2x^2-y^2}}$}{2z=6-2x^2-y^2}}

En este caso, una función que parametriza esta región es:
{\small
\[
{\\\alpha}(u,v) =
\left\langle
\sqrt{3}u\cos v,\;
2u\sin v,\;
3 - \frac{1}{2}u^{2}\left(5 + \cos 2v\right)
\right\rangle\!,
u \in [0,1],\; v \in [0, 2\pi].
\]
}

Con esta función se obtienen los siguientes vectores:
\begin{eqnarray*}
{\bm\alpha}_{u}&=&\left\langle \sqrt{3}\cos v,2\sin v,-u\left( \cos 2v+5\right)
\right\rangle \\
{\bm\alpha}_{v}&=&\left\langle -\sqrt{3}u\sin v,2u\cos v,u^{2}\sin 2v\right\rangle \\
{\bm\alpha}_{u}\times {\bm\alpha}_{v}&=&\left\langle 12u^{2}\cos v,4\sqrt{3}u^{2}\sin v,2\sqrt{3}u\right\rangle
\end{eqnarray*}
${\bf F}\left({\bm \alpha}(u,v)\right)=\left\langle 4u^2\sin^2 v,3u^2\cos^2 v,3-\frac{1}{2}u^{2}\left(5+\cos 2v\right)\right\rangle $, entonces \\
${\bf F}\left({\bm\alpha}(u,v)\right) \cdot \left({\bm\alpha}_{u}\times {\bm\alpha} _{v}\right)
=48u^4\sin^2v\cos v+ \sqrt{3}u(6-u^2(\cos 2v+5)+12u^3\cos^2 v \sin v)$.
Con lo cual, el flujo sobre la superficie está dado por
\begin{equation}\label{flujo1}
\int_{0}^{1}\int_{0}^{2\pi}{\bf F}\left({\bm\alpha}(u,v)\right) \cdot \left({\bm\alpha}_{u}\times {\bm\alpha} _{v}\right)dvdu=\frac{7}{2}\sqrt{3}\pi
\end{equation}
El proceso descrito antes se obtiene con las siguientes instrucciones:

\begin{tcolorbox}[title=\bf Instrucción \verb"Mathematica"]\begin{Verbatim}[formatcom=\letraverbatim]
r={x,y,z}={Sqrt[3]*u*Cos[v],2*u*Sin[v],
(6-2x^2-y^2)/2};
Integrate[Dot[Cross[D[r,u],D[r,v]],{y^2,x^2,z}],
{u,0,1},{v,0,2*Pi}]
\end{Verbatim}
\end{tcolorbox}

\subsubsection{Flujo sobre el paraboloide \texorpdfstring{$\bm{\mathsf{z+3=x^2+y^2}}$}{z+3=x^2+y^2}}

Esta superficie se puede parametrizar mediante la siguiente función:
{\small
\[{\bm \beta} (u,v)=\left\langle \sqrt{3}u\cos v,2u\sin v,\frac{1}{2}u^{2}\left(7-\cos 2v\right)-3\right\rangle\!,
u\in \lbrack 0,1],\, v\in \lbrack 0,2\pi].\]
}

En este sentido,
\[{\bf F}({\bm \beta}(u,v))=\big \langle 4u^2\sin^2v,3u^2\cos^2v, \linebreak
\frac{1}{2}u^{2}(7-\cos 2v)-3\big \rangle .\]
Derivando parcialmente se hallan los siguientes vectores:
{\small
\begin{eqnarray*}
{\bm \beta}_{u}&=&\left\langle \sqrt{3}\cos v,2\sin v,u\left( 7-\cos 2v\right)\right\rangle \\
{\bm \beta}_{v}&=&\left\langle -\sqrt{3}u\sin v,2u\cos v,u^{2}\sin 2v\right\rangle \\
{\bm \beta}_{v}\times {\bm \beta}_{u}&=&\left\langle 12u^{2}\cos v,8\sqrt{3}u^{2}\sin v,-2\sqrt{3}u\right\rangle\!.
\end{eqnarray*}
}

Además,
{\footnotesize
\begin{multline*}
{\bf F}\left({\bm \beta}(u,v)\right) \cdot \left({\bm \beta}_{v}\times {\bm \beta}_{u}\right)
= 12u^{4}(2\sin v+\sqrt{3}\cos v)\sin(2v)+\sqrt{3}u(u^{2}(\cos(2v)-7) +6).
\end{multline*}
}
Con lo cual, el flujo sobre esta región está dado por
{\small
\begin{equation}\label{flujo2}
\int_{0}^{1}\int_{0}^{2\pi}{\bf F}\left({\bm \beta}(u,v)\right) \cdot \left({\bm \beta}_{v}\times {\bm \beta}_{u}\right)dvdu=\frac{5}{2}\sqrt{3}\pi
\end{equation}
}

El flujo sobre la superficie del sólido se determina sumando las expresiones \ref{flujo1} y \ref{flujo2}.
Dado que la ${\rm div}{\bf F}=1$, aplicando el teorema de la divergencia, el valor obtenido coincide con el volumen del sólido en la fórmula \ref{ejercicio1} de la página \pageref{ejercicio1}.

Con las siguientes instrucciones se obtiene el cálculo anterior:

\begin{tcolorbox}[title=\bf Instrucción \verb"Mathematica"]\begin{Verbatim}[formatcom=\letraverbatim]
r={x,y,z}={Sqrt[3]*u*Cos[v],2*u*Sin[v],x^2+y^2-3};
Integrate[Dot[Cross[D[r,v],D[r,u]],{y^2,x^2,z}],
{u,0,1},{v,0,2*Pi}]
\end{Verbatim}
\end{tcolorbox}

\item Hallar el flujo neto sobre el campo vectorial
${\bf F}(x,y,z) = \la y,-x,z \ra $ sobre la superficie frontera del sólido acotado por el cilindro $x^2+y^2=4$ y la esfera $x^2+y^2+z^2=9$.

\begin{figure}[H]
\centering
\caption{${\bf F}(x,y,z) = \la y,-x,z \ra$ sobre la superficie $x^2+y^2\leq 4$ \& $x^2+y^2+z^2\leq 9$}
\includegraphics[width=0.4\textwidth]{Img/T_9b.pdf} \hspace{5mm}
\includegraphics[width=0.4\textwidth]{Img/T_9c.pdf}
\fuentepropia
\end{figure}

La figura presentada se obtiene ejecutando las siguientes instrucciones:

\begin{tcolorbox}[title=\bf Instrucción \verb"Mathematica"]\begin{Verbatim}[formatcom=\letraverbatim]
a=RegionPlot3D[x^2+y^2<4&&x^2+y^2+z^2<=9,{x,-2,2},
{y,-2,2},{z,-3.5,3.5},PlotPoints->90,Mesh->False]
b=SliceVectorPlot3D[{y,-x,z},{x^2+y^2==4,x^2+y^2
+z^2==9},{x,-2,2},{y,-2,2},{z,-3.4,3.5},
VectorPoints->16,RegionFunction->Function[{x,y,z},
x^2+y^2<=4&&x^2+y^2+z^2<=9],VectorStyle->Black,
VectorScale->{0.15,Scaled[0.2],Automatic},
PlotStyle->Opacity[0]];
Show[a,b]
\end{Verbatim}
\end{tcolorbox}

La intersección de las superficies satisface la relación $z=\pm \sqrt{5}$.
El interior del cilindro se parametriza mediante las expresiones $x=2u\cos v$, $y=2u\sin v$,
para $u\in \lbrack 0,1]$, $v\in \lbrack 0,2\pi ]$. Para este sólido se consideran los siguientes dos casos.

\subsubsection{Flujo sobre la superficie de la esfera}
Utilizaremos la función
${\bf r}\left( u,v\right) =\langle 2u\cos v,2u\sin v,\sqrt{9-4u^{2}}\rangle $,
como parametrización de la esfera, con lo cual ${\bf F}({\bf r}(u,v))=\langle 2u\sin v,-2u\cos v, \linebreak \sqrt{9-4u^{2}}\rangle $.
Derivando parcialmente se obtienen los siguientes vectores:

\begin{eqnarray*}
{\bf r}_{u} &=& \left\langle 2\cos v,2\sin v,\frac{-4u}{\sqrt{9-4u^{2}}}\right\rangle \\
{\bf r}_{v}&=&\left\langle -2u\sin v,2u\cos v,0\right\rangle \\
{\bf r}_{u}\times {\bf r}_{v}&=&\left\langle \frac{8u^{2}\cos v}{\sqrt{9-4u^{2}}},\frac{8u^{2}\sin v}{\sqrt{9-4u^{2}}},4u\right\rangle\!.
\end{eqnarray*}

Con lo cual, {\small ${\bf F}\left({\bf r}\left(u,v\right)\right) \cdot \left({\bf r}_{u}\times {\bf r}_{v}\right)= 4u\sqrt{9-4u^{2}}$.}
Al ser dos regiones iguales (sobre la esfera) multiplicamos por 2 y haremos un solo cálculo, es decir, el flujo sobre la porción de la esfera que limita el sólido está dado por

\begin{equation}\label{flujo3}
2\int_{0}^{1}\int_{0}^{2\pi }4u\sqrt{9-4u^{2}}dvdu= \frac{4}{3}\pi \left( 27-5\sqrt{5}\right)
\end{equation}

El cálculo anterior se obtiene en \verb"Mathematica" con las siguientes instrucciones:

\begin{tcolorbox}[title=\bf Instrucción \verb"Mathematica"]\begin{Verbatim}[formatcom=\letraverbatim]
r={x,y,z}={2*u*Cos[v],2*u*Sin[v],Sqrt[9-x^2-y^2]};
Integrate[Dot[Cross[D[r,u],D[r,v]],{y,-x,z}],
{u,0,1},{v,0,2*Pi}]
\end{Verbatim}
\end{tcolorbox}

\subsubsection{Flujo sobre la superficie del cilindro}

En este caso la parametrización escogida es
${\bf w}\left( u,v\right) =\left\langle 2\cos u,2\sin u,v\right\rangle $, para $ u\in \lbrack 0,2\pi ]$, $v\in \lbrack -\sqrt{5},\sqrt{5}]$.
Con esta función hallamos los siguientes vectores:
\begin{eqnarray*}
{\bf w}_{u}&=&\left\langle -2\sin u,2\cos u,0\right\rangle \\
{\bf w}_{v}&=&\left\langle 0,0,1\right\rangle \\
{\bf w}_{u}\times {\bf w}_{v}&=&\left\langle 2\cos u,2\sin u,0\right\rangle\!.
\end{eqnarray*}
Además, ${\bf F}\left( {\bf w}\left( u,v\right) \right) =\left\langle 2\sin u,-2\cos u,v\right\rangle $
y ${\bf F}\left({\bf w}\left( u,v\right) \right) \cdot \left({\bf w}_{u}\times {\bf w}_{v}\right)= 0$,

con lo cual el flujo sobre la superficie del cilindro esta dado por
\begin{equation}\label{flujo4}
\int_S {\bf F}\cdot d{\bf S}=0
\end{equation}

El flujo neto es la suma de los resultados \ref{flujo3} y \ref{flujo4}. Como ${\rm div}\,{\bf F}=1$ en este caso, el valor hallado coincide con el volumen del sólido encontrado en la expresión \ref{ejercicio4} de la página \pageref{ejercicio4}.

\item Hallar el flujo neto sobre el campo vectorial ${\bf F}(x,y,z) = \la z,y,x \ra $
sobre la superficie frontera del sólido acotada por el hiperboloide de un manto con ecuaciones
$2x^2+3y^2=z^2+7$ y el elipsoide $2x^2+y^2+z^2=9$.

\label{elipsoidehiperboloide}

\begin{figure}[H]
\centering
\caption{Campo vectorial ${\bf F}(x,y,z)=\la z,y,x\ra$}
\includegraphics[width=0.4\textwidth]{Img/T_11c.pdf} \hspace{5mm}
\includegraphics[width=0.4\textwidth]{Img/T_11d.pdf}
\fuentepropia
\end{figure}

Con las siguientes instrucciones se consigue esta figura.

\begin{tcolorbox}[title=\bf Instrucción \verb"Mathematica"]\begin{Verbatim}[formatcom=\letraverbatim]
a=RegionPlot3D[2*x^2+3*y^2<z^2+7&&2*x^2+y^2+z^2<9,
{x,-2,2},{y,-2,2},{z,-3,3},PlotPoints->90,Mesh->False];
b=SliceVectorPlot3D[{z,y,x},{2*x^2+3*y^2==z^2+7,
2*x^2+y^2+z^2==9},{x,-Sqrt[15/4],Sqrt[15/4]},
{y,-Sqrt[15/2],Sqrt[15/2]},{z,-4,3},VectorPoints->20,
VectorStyle->Black,RegionFunction->Function[{x,y,z},
x^2+y^2<=4],VectorScale->{0.07,0.2,None},
PlotStyle->Opacity[0]];
Show[a,b]
\end{Verbatim}
\end{tcolorbox}

Calcularemos el flujo sobre la superficie que acota el sólido, es decir, la porción del elipsoide y, posteriormente, la parte del hiperboloide.

\subsubsection{Flujo sobre la porción del elipsoide}

Con las ecuaciones que definen el elipsoide y el hiperboloide de un manto se elimina $z$ y se obtiene la expresión
$x^2+y^2=4;$ la región que esta ecuación representa se parametriza mediante el uso de las expresiones
$x=2u\cos v$ y $y=2u \sin v$, con $u \in [0,2\pi]$ y $u\in [0,1]$.
El término que define a $z$ está dado por
\[z =\sqrt{9-2x^2-y^2} = \sqrt{9-6u^2-2u^2 \cos(2v)}.\]
Entonces; la función
\[{\bf r}(u,v) = \left \langle 2u\cos v,2u \sin v ,\sqrt{9-6u^2-2u^2 \cos(2v)} \right\rangle\]
parametriza la región de interés sobre el elipsoide. Con esta función se obtienen los siguientes vectores:
{\small
\begin{eqnarray*}
{\bf r}_u &=& \left \langle 2\cos v,2 \sin v , -\frac{2u(3+2\cos v)}{\sqrt{9-6u^2-2u^2 \cos(2v)}} \right\rangle \\
{\bf r}_v &=& \left \langle -2u\sin v,2u \cos v , \frac{2u^2 \sin v}{\sqrt{9-6u^2-2u^2 \cos(2v)}} \right\rangle \\
{\bf r}_u \times {\bf r}_v &=& \left \langle \frac{16u^2 \cos v}{\sqrt{9-6u^2-2u^2 \cos(2v)}}, \frac{8u^2 \sin v}{\sqrt{9-6u^2-2u^2 \cos(2v)}}, 4u \right\rangle\!.
\end{eqnarray*}
}

Además
${\bf F}({\bf r}(u,v) ) = \left \langle \sqrt{9-6u^2-2u^2 \cos(2v)},2u \sin v , 2u\cos v \right\rangle \!$,
con lo cual
{\small
\[{\bf F}({\bf r}(u,v) ) \cdot ({\bf r}_u \times {\bf r}_v)=\frac{ 8u \left (u-u \cos (2v) +3 \cos v \, \sqrt{9-2u^2(3+\cos(2v))}\right )}{\sqrt{9-6u^2-2u^2 \cos(2v)}};\]
}

por lo tanto, el flujo sobre la porción del elipsoide está dado por

\begin{equation} \label{ejercicio6e}
\int_0^1 \int_0^{2\pi} {\bf F}({\bf r}(u,v) ) \cdot ({\bf r}_u \times {\bf r}_v) dvdu \approx 10.844
\end{equation}
Este cálculo se hizo en \verb"Mathematica" ejecutando las siguientes instrucciones:

\begin{tcolorbox}[title=\bf Instrucción \verb"Mathematica"]\begin{Verbatim}[formatcom=\letraverbatim]
r={x,y,z}={2*u*Cos[v],2*u*Sin[v],Sqrt[9-2*x^2-y^2]};
NIntegrate[2*Dot[{z,y,x},Cross[D[r,u],D[r,v]]],
{v,0,2*Pi},{u,0,1}]
\end{Verbatim}
\end{tcolorbox}

Antes de presentar el siguiente procedimiento, nótese que la función $f(u)= (b-a)u+a$ satisface que $f(0)=a$ y $f(1)=b$.
De esta manera $f$ aplica el intervalo $[0,1]$ en el intervalo $[a,b]$. A su vez la función inversa
$g(v)= \frac{v-a}{b-a}$ aplica el intervalo $[a,b]$ en el intervalo $[0,1]$.
Usaremos estas ideas para parametrizar la región acotada por la porción de la circunferencia y la hipérbola que se des\-cribe a continuación.

\subsubsection{Parametrizando la región sobre el hiperboloide}

Al eliminar $y$ de las ecuaciones originales se obtiene $x^2+z^2=5$,
esta expresión representa, en el plano $xz$, una circunferencia de radio $\sqrt{5}$ que se parametriza como
$x=\sqrt{5}\cos t$, $z=\sqrt{5}\sin t$.
El hiperboloide proyectado en el plano $xz$ corresponde a un una hipérbola con ecuación $2x^2-z^2=7$.
Esta región se puede describir definiendo $x=\sqrt{\frac{7}{2}}\cosh (t),\,\, z=\sqrt{7}\sinh (t)$.
Las anteriores dos curvas se intersectan en los puntos $(\pm 2, \pm 1)$.

Entonces la región sobre la circunferencia para $x\in [-2,2]$ se describe por medio de la función
\[\overrightarrow{\alpha}(t)=\langle \sqrt{5}\cos t, \sqrt{5}\sin t \rangle, \,\,
t \in \left[ \arccos \left( \frac{2}{\sqrt{5}}\right),\pi -\arccos \left( \frac{2}{\sqrt{5}}\right) \right].\]

La región sobre la hipérbola desde el punto $(2,-1)$ al punto $(2,1)$ se describe por medio de la función
\[\overrightarrow{\beta}(t) = \left \langle \sqrt{\frac{7}{2}}\cosh (t), \sqrt{7}\sinh (t) \right \rangle, \,t \in \left[ -t_0 ,t_0 \right].\]
Se usó la abreviatura $t_0={\rm arccosh}\left( \sqrt{\frac{2}{7}}\right)$.
Estas curvas se muestran a continuación:

\begin{figure}[H]
\centering
\caption{Parametrización de la proyección del sólido en plano $xz$}
\begin{minipage}[c]{0.39\linewidth}
\centering

\includegraphics[width=\linewidth]{Fig/188.pdf}
\end{minipage}
\hspace{1em}
\begin{minipage}[c]{0.32\linewidth}
\centering

\includegraphics[width=\linewidth]{Fig/189.pdf}
\end{minipage}

\fuentepropia
\end{figure}

Para abreviar la escritura, usaremos las siguientes constantes:
{\small
\[a= \arccos \left( \frac{2}{\sqrt{5}}\right),\,\, b=\pi -\arccos \left( \frac{2}{\sqrt{5}}\right)\]
\[c= -{\rm arccosh}\left( \frac{2\sqrt{2}}{\sqrt{7}}\right) ,\, d={\rm arccosh}\left( \frac{2\sqrt{2}}{\sqrt{7}}\right)\!.\]}

Al redefinir las funciones que describen las curvas en mención con el sentido deseado y unificar el intervalo de definición a $t \in [0,1]$,

{\small
\begin{eqnarray*}
\overrightarrow{\alpha}(t) &=& \left \langle \sqrt{7/2}\cosh((d-c)t+c),\sqrt{7}\sinh((d-c)t+c)\right \rangle,\\
\overrightarrow{\beta}(t) &=&\left \langle \sqrt{5}\cos((b-a)(1-t)+a),\sqrt{5}\sin((b-a)(1-t)+a) \right \rangle.
\end{eqnarray*}
}

Formando una combinación convexa de estas expresiones se parametriza la región sombreada en el gráfico.
Con las siguientes instrucciones se dibujan las curvas mencionadas y se consigue una animación en \verb"Mathematica" que permite ver el recorrido de la región que se considera.

\begin{tcolorbox}[title=\bf Instrucción \verb"Mathematica"]\begin{Verbatim}[formatcom=\letraverbatim]
{a,b}={ArcCos[2/Sqrt[5]],Pi-a};
{c,d}={-ArcCosh[2*Sqrt[2/7]],c};
R[w_]:={Sqrt[5]*Cos[(b-a)*(1-w)+a],
Sqrt[5]*Sin[(b-a)*(1-w)+a]};
r[w_]:={Sqrt[7/2]*Cosh[(d-c)*w+c],
Sqrt[7]*Sinh[(d-c)*w+c]};
Animate[ParametricPlot[{R[u],(1-u)*r[t]+u*R[t],r[u]},
{u,0,1},PlotRange->{{-2,3},{-2,3}},PlotStyle->{Red,
Black,Red}],{t,0,1}]
\end{Verbatim}
\end{tcolorbox}

La función obtenida al hacer esta construcción es la siguiente:

\begin{multline}
{\bf r}(t,u)=(1-u)\left(\sqrt{7/2}\cosh((d-c)t+c),\sqrt{7}\sinh((d-c)t+c)\right )+ \\
u\left (\sqrt{5}\cos((b-a)(1-t)+a),\sqrt{5}\sin((b-a)(1-t)+a)\right) \label{ejercicio6ec}
\end{multline}
En donde $t \in [0,1],\,\, u \in [0,1]$. La expresión correspondiente a $y$ se consigue al sustituir $x$ y $z$
de la función anterior en la expresión que define el hiperboloide de un manto. La región obtenida corresponde a la cuarta parte de la superficie mencionada.

\begin{figure}[H]
\centering
\caption{Parametrización de la superficie del sólido}
\includegraphics[width=0.4\textwidth]{Img/T_11e.pdf} \hspace{5mm}
\includegraphics[width=0.4\textwidth]{Img/T_11f.pdf}
\fuentepropia
\end{figure}

\begin{tcolorbox}[title=\bf Instrucción \verb"Mathematica"]\begin{Verbatim}[formatcom=\letraverbatim]
r={x,y,z}={2*u*Cos[v],2*u*Sin[v],Sqrt[9-2*x^2-y^2]}
A1=ParametricPlot3D[{{x,y,z},{x,y,-z}},{u,0,1},
{v,0,2*Pi},PlotStyle->{RGBColor[0.8,0.2,0.2]},
MeshStyle->Black,PlotPoints->90];
{a,b}={ArcCos[2/Sqrt[5]],Pi-ArcCos[2/Sqrt[5]]};
{c,d}={-ArcCosh[2*Sqrt[2/7]],ArcCosh[2*Sqrt[2/7]]};
R=(1-u)*{Sqrt[7/2]*Cosh[(d-c)*t+c],Sqrt[7]*Sinh[(d-
c)*t+c]}+u*{Sqrt[5]*Cos[(b-a)*(1-t)+a],
Sqrt[5]*Sin[(b-a)*(1-t)+a]};
{x,z,y}={R[[1]],R[[2]],Sqrt[(z^2+7-2*x^2)/3]};
A2=ParametricPlot3D[{{x,y,z},{x,-y,z},-{x,y,z},
-{x,-y,z}},{u,0,1},{t,0,1},PlotStyle->{Cyan,Orange,
Cyan,Orange},PlotPoints->60];
Show[A2,A1,BoxRatios->{1,1,1}]
\end{Verbatim}
\end{tcolorbox}

Utilizando la expresión \ref{ejercicio6ec}, se establece que el flujo sobre la superficie del hiperboloide $H$,
está dado por la expresión
\begin{equation} \label{ejercicio6h}
4 \iint_H {\bf F}({\bf r}(t,u) ) \cdot ({\bf r}_u \times {\bf r}_v) dt dt \approx 41.0176
\end{equation}
Este valor se obtuvo con \verb"Mathematica" con estas instrucciones:

\begin{tcolorbox}[title=\bf Instrucción \verb"Mathematica"]\begin{Verbatim}[formatcom=\letraverbatim]
{a,b,d,c}={ArcCos[2/Sqrt[5]],Pi-a,
ArcCosh[2*Sqrt[2/7]],-d};
R=(1-u)*{Sqrt[7/2]*Cosh[(d-c)*t+c],
Sqrt[7]*Sinh[(d-c)*t+c]}+u*{Sqrt[5]*Cos[(b-a)*(1-t)
+a],Sqrt[5]*Sin[(b-a)*(1-t)+a]};
r={x,y,z}={R[[1]],Sqrt[(z^2+7-2*x^2)/3],R[[2]]};
NIntegrate[4*Dot[{z,y,x},Cross[D[r,u],D[r,t]]],
{u,0,1},{t,0,1}]
\end{Verbatim}
\end{tcolorbox}

El flujo sobre la superficie del sólido es la suma de las expresiones
\ref{ejercicio6h} y \ref{ejercicio6e}.
Como ${\rm div}\, {\rm F} =1$, el valor del flujo neto de ${\bf F}$ coincide con el volumen del sólido que se halló en la página \pageref{ejercicio6}, con la expresión \ref{ejercicio6}.

\item Hallar el flujo neto sobre el campo vectorial
${\bf F}(x,y,z) = \la 3x,y,3z \ra $ sobre la superficie frontera del sólido acotado por el paraboloide $x^2+y^2=3-z$ y el paraboloide hiperbólico $x^2-y^2=3z+6$. \index{Paraboloide!hiperbólico}

\begin{figure}[H]
\centering
\caption{Campo vectorial ${\bf F}(x,y,z)=\la 3x,y,3z\ra$}
\includegraphics[width=0.4\textwidth]{Img/T_10c.pdf} \hspace{5mm}
\includegraphics[width=0.4\textwidth]{Img/T_10d.pdf}
\fuentepropia
\end{figure}

Esta figura es el producto de ejecutar las siguientes instrucciones:

\begin{tcolorbox}[title=\bf Instrucción \verb"Mathematica"]\begin{Verbatim}[formatcom=\letraverbatim]
a=RegionPlot3D[x^2+y^2<3-z&&x^2-y^2<3*z+6,{x,-2,2},
{y,-3,3},{z,-5,3},PlotPoints->90,Mesh->False]
b=SliceVectorPlot3D[{3*x,y,3*z},{x^2+y^2==3-z,x^2-
y^2==3*z+6},{x,-Sqrt[15/4],Sqrt[15/4]},{y,-Sqrt[15/2],
Sqrt[15/2]},{z,-5,3},RegionFunction->Function[{x,y,z},
4*x^2+2*y^2<=15],VectorPoints->18,VectorStyle->Black,
VectorScale->{0.14,Scaled[0.25],Automatic},
PlotStyle->Opacity[0]];
Show[a,b]
\end{Verbatim}
\end{tcolorbox}

En la página \pageref{ejercicio4a} se trata esta región.
La expresión $4x^{2}+2y^{2}=15$ describe la curva de intersección en el plano $xy$. Esta región puede representarse mediante las expresiones
$x=\frac{\sqrt{15}}{2}u\cos v$, $y=\sqrt{\frac{15}{2}}u\sin v$, con $u\in \lbrack 0,1]$,
$v\in \lbrack 0,2\pi ]$.
Calcularemos el flujo neto sobre la superficie del sólido en dos partes: primero sobre la superficie del paraboloide y luego sobre el paraboloide hiperbólico.

\subsubsection{Flujo sobre el paraboloide}

Al reemplazar las expresiones anteriores en la ecuación $x^2+y^2=3-z$, se obtiene que
$z=\frac{15}{8}u^{2}\cos 2v-\frac{45}{8}u^{2}+3;$
por lo tanto, una parametrización de la superficie en mención es
\[{\bf r}\left( u,v\right) =\left\langle \frac{\sqrt{15}}{2}u\cos v,\sqrt{\frac{15}{2}}u\sin v,
\frac{15}{8}u^{2}\cos 2v-\frac{45}{8}u^{2}+3\right\rangle\!,\]
en donde $u\in \lbrack 0,1]$, $v\in \lbrack 0,2\pi ]$.
Con la función ${\bf r}$, al evaluarla en el campo vectorial se sigue que
\[{\bf F}\left({\bf r}\left( u,v\right) \right) =\left\langle \frac{3\sqrt{15}}{2}u\cos v,\sqrt{\frac{15}{2}}u\sin v,\frac{45}{8}u^{2}\cos 2v-\frac{135}{8}u^{2}+9\right\rangle\!.\]
% ${\bf F}\left({\bf r}\left( u,v\right) \right) =\left\langle \frac{3\sqrt{15}}{2}u\cos v,\sqrt{\frac{15}{2}}u\sin v,\frac{45}{8}u^{2}\cos 2v-\frac{135}{8}u^{2}+9\right\rangle ,$
Derivando parcialmente se obtienen los siguientes vectores;
\begin{eqnarray*}
{\bf r}_{u} &=&\left\langle \frac{1}{2}\sqrt{15}\cos v,\frac{1}{2}\sqrt{2}\sqrt{15}%
\sin v,\frac{15}{4}u\cos 2v-\frac{45}{4}u \right\rangle \\
{\bf r}_{v}&=& \left\langle -\frac{1}{2}\sqrt{15}u\sin v,\frac{1}{2}\sqrt{2}\sqrt{15}u\cos v,-\frac{15}{4}u^{2}\sin 2v \right\rangle \\
{\bf r}_{u}\times {\bf r}_{v}&=&\left\langle \frac{15}{4}\sqrt{2}\sqrt{15}u^{2}\cos v,\frac{15}{2}\sqrt{15}u^{2}\sin v,\frac{15}{4}\sqrt{2}u\right\rangle;
\end{eqnarray*}
entonces,
\[{\bf F}\left({\bf r}\left(u,v\right) \right) \cdot \left({\bf r}_{u}\times {\bf r}_{v}\right) =\frac{225\sqrt{2}}{32}u\left( u^{2}+\frac{24}{5}+5u^{2}\cos 2v\right)\!.\]
Por lo expuesto antes, el flujo sobre esta región está dado por
\begin{equation}\label{flujo5}
\int_{0}^{1}\int_{0}^{2\pi }\left( \frac{225\sqrt{2}}{32}u\left( u^{2}+ \frac{24}{5}+5u^{2}\cos 2v\right) \right) dvdu= \frac{2385}{64}\sqrt{2}\pi
\end{equation}
Este valor también se obtiene como producto de ejecutar las siguientes instrucciones:

\begin{tcolorbox}[title=\bf Instrucción \verb"Mathematica"]\begin{Verbatim}[formatcom=\letraverbatim]
r={x,y,z}={Sqrt[15]/2*u*Cos[v],Sqrt[15/2]*u*Sin[v],
3-x^2-y^2};
Integrate[Dot[Cross[D[r,u],D[r,v]],{3x,y,3z}],
{u,0,1},{v,0,2*Pi}]
\end{Verbatim}
\end{tcolorbox}

\subsubsection{Flujo sobre el paraboloide hiperbólico}

Con las mismas ecuaciones que describen el conjunto de puntos que satisfacen la condición $4x^{2}+2y^{2} \leq 15$,
la coordenada $z$ sobre esta superficie satisface que
\[z=\frac{x^2-y^2-6}{6}=\frac{15}{8}u^{2}\cos 2v-\frac{5}{8}u^{2}-2,\]
por lo tanto, utilizaremos la siguiente parametrización, para describir la porción del paraboloide hiperbólico:

{\small
\[{\bf r}\left( u,v\right) =\left\langle \frac{\sqrt{15}}{2}u\cos v,\sqrt{\frac{15}{2}}u\sin v,\frac{15}{8}u^{2}\cos 2v-\frac{5}{8}u^{2}-2\right\rangle\!,\]}

en donde $u\in \lbrack 0,1]$, $v\in \lbrack 0,2\pi ]$.

Con ella se consiguen los siguientes vectores:

{\small
\begin{align*}
{\bf r}_{u}&=\left\langle \frac{1}{2}\sqrt{15}\cos v,\frac{1}{2}\sqrt{2}\sqrt{15}\sin v,\frac{15}{4}u\cos 2v-\frac{5}{4}u\right\rangle \\
{\bf r}_{v}&= \left\langle -\frac{1}{2}\sqrt{15}u\sin v,\frac{1}{2}
\sqrt{2}\sqrt{15}u\cos v,-\frac{15}{4}u^{2}\sin 2v \right\rangle \\
{\bf r}_{v}\times {\bf r}_{u}&=\left\langle \frac{5}{4}\sqrt{2}\sqrt{15}u^{2}\cos v,-\frac{5}{2}\sqrt{3}\sqrt{5}u^{2}\sin v,-\frac{15}{4}\sqrt{2}u\right\rangle\!.
\end{align*}
}

Es de anotar que para conservar la orientación de la superficie se calculó ${\bf r}_{v}\times {\bf r}_{u}$. Al evaluar la parametrización en el campo vectorial, se sigue que

{\small
\[{\bf F}\left( {\bf r}\left( u,v\right) \right) =\left\langle \frac{3\sqrt{15}}{2}u\cos
v,\sqrt{\frac{15}{2}}u\sin v,\frac{45}{8}u^{2}\cos 2v-\frac{15}{8}u^{2}-6\right\rangle\!,\]}

entonces,

{\small
\[{\bf F}\left({\bf r}\left( u,v\right) \right) \cdot \left({\bf r}_{u}\times {\bf r}_{v}\right) =
\frac{75\sqrt{2}}{32}u\left( 5u^{2}+\frac{48}{5}+u^{2}\cos 2v\right)\!.\]}

Por lo señalado antes, el flujo sobre esta superficie está dado por
{\small
\begin{equation}\label{flujo6}
\int_{0}^{1}\int_{0}^{2\pi }\left( \frac{75\sqrt{2}}{32}u\left( 5u^{2} + \frac{48}{5}+u^{2}\cos 2v\right) \right) dvdu= \frac{1815}{64}\sqrt{2}\pi
\end{equation}
}

Con las siguientes instrucciones se hace el cálculo anterior:

\begin{tcolorbox}[title=\bf Instrucción \verb"Mathematica"]\begin{Verbatim}[formatcom=\letraverbatim]
r={x,y,z}={Sqrt[15]/2*u*Cos[v],
Sqrt[15/2]*u*Sin[v],(x^2-y^2-6)/6};
Integrate[Dot[Cross[D[r,v],D[r,u]],{3x,y,3z}],
{u,0,1},{v,0,2*Pi}]
\end{Verbatim}
\end{tcolorbox}

El flujo neto sobre toda la superficie del sólido es la suma de los resultados hallados en \ref{flujo5} y
\ref{flujo6}, esto es \[\frac{2385}{64}\sqrt{2}\pi +\frac{1815}{64}\sqrt{2}\pi = \frac{525}{8}\sqrt{2}\pi.\]
Dado que ${\rm div}{\bf F}=7$, evidentemente, el valor obtenido es $7$ veces el volumen del sólido que se obtuvo en la expresión \ref{ejercicio3} de la página \pageref{ejercicio3}

\item Hallar el flujo neto del campo vectorial ${\bf F}(x,y,z) = \la y,-x,z \ra $
sobre la superficie del sólido interior a la esfera con ecuacion $x^2+y^2+z^2=9 $ y al cilindro cardioidico
$r=1+\cos \theta$.

\begin{figure}[H]
\centering
\caption{Campo vectorial ${\bf F}(x,y,z)=\la y,-x,z\ra$, sobre $x^2+y^2+z^2=9 $ \& $r=1+\cos \theta$}
\label{cvFxyzsobrex2y2z2}
\includegraphics[width=0.4\textwidth]{Img/T_12c.pdf} \hspace{5mm}
\includegraphics[width=0.4\textwidth]{Img/T_12d.pdf}
\fuentepropia
\end{figure}

Las instrucciones que permiten obtener la figura~\ref{cvFxyzsobrex2y2z2} son las siguientes.

\begin{tcolorbox}[title=\bf Instrucción \verb"Mathematica"]\begin{Verbatim}[formatcom=\letraverbatim]
a=RegionPlot3D[(x^2+y^2-x)^2<(x^2+y^2)&&
x^2+y^2+z^2<9,{x,-1,2.2},{y,-2,2},{z,-3,3.2},
PlotPoints->90,Mesh->False,BoxRatios->{1,1,1}];
b=SliceVectorPlot3D[{y,-x,z},{(x^2+y^2-x)^2==x^2+y^2,
x^2+y^2+z^2==9},{x,-1,2.2},{y,-2,2},{z,-3,3.2},
RegionFunction->Function[{x,y,z},(x^2+y^2-x)^2<(x^2
+y^2)&&x^2+y^2+z^2<9],VectorPoints->17,
VectorStyle->Black,VectorScale->{0.12,Scaled[0.25],
Automatic},PlotStyle->Opacity[0]];
Show[a,b]
\end{Verbatim}
\end{tcolorbox}

A partir de la ecuación $r = 1 + \cos \theta$, se puede construir en el plano $xy$ la parametrización de la región que rellena esta cardioide. Usaremos las ecuaciones:
\[
x = u(1 + \cos \theta) \cos \theta, \quad
y = u(1 + \cos \theta) \sin \theta, \quad
u \in [0,1],\; \theta \in [0,2\pi].
\]

La región sobre la esfera se parametriza completamente hallando $z$, es decir:
\[
z = \sqrt{9 - x^2 - y^2} = \sqrt{9 - u^2(1 + \cos \theta)^2}.
\]

Por lo tanto, la siguiente función describe la porción sobre la esfera interior al cilindro cardioidico:
\[
{\bf r}(u,\theta) =
\left\langle
u(1 + \cos \theta) \cos \theta,\;
u(1 + \cos \theta) \sin \theta,\;
\sqrt{9 - u^2(1 + \cos \theta)^2}
\right\rangle\!,
\]
en donde $u \in [0,1],\; \theta \in [0,2\pi]$. Es sencillo ver que:
{\small
\[
{\bf F}({\bf r}(u,\theta)) =
\left\langle
u(1 + \cos \theta) \sin \theta,\;
- u(1 + \cos \theta) \cos \theta,\;
\sqrt{9 - u^2(1 + \cos \theta)^2}
\right\rangle\!.
\]
}

Con la parametrización elegida hallamos los siguientes vectores:
{\footnotesize
\begin{eqnarray*}
{\bf r}_{u} &=& \left \langle (\cos \theta +1) \cos \theta ,\left(\cos \theta +1\right) \sin \theta ,\frac{-4u\cos^2\left( \frac{\theta}{2} \right)}{\sqrt{9-u^2(1+\cos(\theta))^2}}\right\rangle \\
{\bf r}_{\theta}&=&\left\langle -u( 1+2\cos \theta)\sin \theta ,u(\cos \theta +\cos (2\theta)),\frac{u^{2}(\sin 2\theta+2\sin \theta)}{\sqrt{2}\sqrt{9-u^2(1+\cos(\theta))^2}} \right\rangle ,
\end{eqnarray*}
}
cuyo producto cruz es
{\small
\begin{multline*}
{\bf r}_{u}\times {\bf r}_{\theta}= \bigg \langle \frac{8u^2\cos^6(\frac{\theta}{2})\cos(\theta)}{\sqrt{9-u^2(1+\cos(\theta))^2}}, \frac{6u^2\cos^7(\frac{\theta}{2})\sin(\frac{\theta}{2})}{\sqrt{9-u^2(1+\cos(\theta))^2}},4u\cos^4\left(\frac{\theta}{2}\right )\bigg \rangle.
\end{multline*}
}
Además,
\[{\bf F}\left({\bf r}\left( u,\theta \right) \right) \cdot \left({\bf r}_{u}\times {\bf r}_{\theta}\right) =
4u \cos^4\left(\frac{\theta}{2}\right)\sqrt{9-u^2(1+\cos(\theta))^2}.\]
Por lo tanto, el flujo sobre la superficie de la esfera e interior al cilindro cardioidico, esta dado por
\begin{equation}\label{flujo7}
2\int_{0}^{1}\int_{0}^{2\pi}{\bf F}\left({\bf r}\left( u,\theta \right) \right) \cdot \left({\bf r}_{u}\times {\bf r}_{\theta}\right) d\theta du\approx 25.8179
\end{equation}
Se ha multiplicado por $2$ debido a que la región considerada es similar al hemisferio opuesto de la esfera.
El cálculo anterior se hizo en \verb"Mathematica" utilizando las siguientes instrucciones:

\begin{tcolorbox}[title=\bf Instrucción \verb"Mathematica"]\begin{Verbatim}[formatcom=\letraverbatim]
{x,y}={u*(1+Cos[t])*Cos[t],u*(1+Cos[t])*Sin[t]};
r={x,y,Sqrt[9-x^2-y^2]};
A=Dot[{y,-x,z},TrigExpand[Cross[D[r,u],D[r,t]]]];
2*NIntegrate[A,{u,0,1},{t,0,2*Pi}]
\end{Verbatim}
\end{tcolorbox}

Para calcular el flujo sobre el cilindro cardióidico, se debe parametrizar dicha región.
Al proyectar el cilindro sobre el plano $xy$, la curva resultante puede parametrizarse en la forma
$x=\left( 1+\cos \theta \right) \cos \theta $, $y=\left( 1+\cos \theta\right) \sin \theta $, $\theta \in [0,2\pi]$.
Sobre la esfera, la coordenada $z$ puede ser de dos formas
\[y = \pm \frac{\sqrt{15-4\cos(\theta)-\cos(2\theta)}}{\sqrt{2}}.\]
Con dos puntos de la forma;
{\small
\[P=\left(\left( 1+\cos \theta \right) \cos \theta, \left( 1+\cos \theta \right) \sin \theta,\frac{1}{\sqrt{2}}\sqrt{15-4\cos(\theta)-\cos(2\theta)} \right )\]
\[Q=\left(\left( 1+\cos \theta \right)\cos\theta,\left(1+\cos\theta\right)\sin\theta,-\frac{1}{\sqrt{2}}\sqrt{15-4\cos(\theta)-\cos(2\theta)} \right)\]
}

se puede hacer una combinación convexa que determine una función ${\bf r}$ que parametrice la región en mención. Esto es

{\footnotesize
\[{\bf r}(u,\theta ) =\bigg \langle \left( \cos \theta +1\right ) \cos \theta,\linebreak
(1+\cos\theta) \sin \theta, \frac{(1-2u)\sqrt{15-4\cos(\theta) -\cos(2\theta)}}{\sqrt{2}}\bigg\rangle,\]
}

entonces,
{\footnotesize
\[{\bf F}\left({\bf r}\left(u,\theta \right) \right)  \linebreak
=\big \langle \left( 1+\cos \theta \right) \sin \theta , -\left( \cos \theta +1\right) \cos \theta ,\frac{(1-2u)\sqrt{15-4\cos(\theta)-\cos(2\theta)}}{\sqrt{2}} \big \rangle.\]
}

Con la función ${\bf r}$ se obtienen los siguientes vectores:
{\footnotesize
\begin{eqnarray*}
{\bf r}_{u} &=& \left \langle 0,0,-\sqrt{2}\sqrt{15-4\cos(\theta)-\cos(2\theta)}\right\rangle \\
{\bf r}_{\theta}&=& \left\langle -\sin(\theta)-\sin(2\theta),\cos(\theta)+\cos(2\theta),\frac{(1-2u)(2\sin(\theta)+\sin(2\theta))}{\sqrt{2}\sqrt{15-4\cos(\theta)-\cos(2\theta)}} \right \rangle;
\end{eqnarray*}
}

además, estos vectores permiten calcular
\begin{multline*}
{\bf r}_{u}\times {\bf r}_{\theta } = \bigg\langle 2\sqrt{2}\cos^2(\theta/2)(-1+2\cos(\theta))\sqrt{15-4\cos(\theta)-\cos(2\theta)}, \\
\sqrt{2}(1+2\cos(\theta))\sqrt{15-4\cos(\theta)-\cos(2\theta)}\sin(\theta),0 \bigg\rangle .
\end{multline*}
La densidad de flujo está dada por
\[{\bf F}\left({\bf r}\left( u,\theta \right) \right) \cdot \left({\bf r}_{u}\times {\bf r}_{\theta}\right) =\frac{1}{\sqrt{2}}\sqrt{15 - 4 \cos(\theta) - \cos(2\theta)},\]
cuya integral proporciona el flujo de ${\bf F}$ sobre la superficie del cilindro cardióidico, esto es
\begin{equation}\label{flujo8}
\int_{0}^{1}\int_{0}^{2\pi} \frac{1}{\sqrt{2}}\sqrt{15 - 4 \cos(\theta) - \cos(2\theta)}\,
d\theta du=0.
\end{equation}
Estos cálculos se obtienen ejecutando las siguientes instrucciones:

\begin{tcolorbox}[title=\bf Instrucción \verb"Mathematica"]\begin{Verbatim}[formatcom=\letraverbatim]
{x,y,z}={(1+Cos[t])*Cos[t],(1+Cos[t])*Sin[t],
Sqrt[9-x^2-y^2]};
r=Simplify[(1-u)*{x,y,z}+u*{x,y,-z}];
A=Simplify[Dot[{y,-x,z},Cross[D[r,u],D[r,t]]]];
Integrate[A,{u,0,1},{t,0,2*Pi}]
\end{Verbatim}
\end{tcolorbox}

El flujo neto sobre la superficie del sólido considerado es la suma de las expresiones \ref{flujo7}
y \ref{flujo8}.
Dado que ${\rm div} {\bf F}=1$, este valor coincide con el volumen del sólido que se halló con la expresión \ref{ejercicio3}, página \pageref{ejercicio3}.

\item Hallar el flujo neto sobre el campo vectorial ${\bf F}(x,y,z) = \la x,1,y+z \ra $
sobre la superficie sólido acotado por el semicono $y=\sqrt{x^2+z^2}$ y la esfera $x^2+y^2+z^2=9$.

La siguiente figura ilustra esta situación.

\begin{figure}[H]
\centering
\caption{Campo vectorial ${\bf F}(x,y,z) = \la x,1,y+z \ra $ }
\includegraphics[width=0.4\textwidth]{Img/T_14c.pdf} \hspace{5mm}
\includegraphics[width=0.4\textwidth]{Img/T_14d.pdf}
\fuentepropia
\end{figure}

La figura se consigue ejecutando las siguientes instrucciones:

\begin{tcolorbox}[title=\bf Instrucción \verb"Mathematica"]\begin{Verbatim}[formatcom=\letraverbatim]
a=RegionPlot3D[y>Sqrt[x^2+z^2]&&x^2+y^2+z^2<9,
{x,-3,3},{y,0,3},{z,-3,3},PlotPoints->90,Mesh->False];
b=SliceVectorPlot3D[{x,1,y+z},{x^2+y^2+z^2==9,
y==Sqrt[x^2+z^2]},{x,-3,3},{y,0,3},{z,-3,3},
VectorPoints->18,VectorScale->{0.16,Scaled[0.25],
Automatic},VectorStyle->Black,RegionFunction->
Function[{x,y,z},x^2+z^2<=9/2],PlotStyle->Opacity[0]];
Show[a,b]
\end{Verbatim}
\end{tcolorbox}

Las superficies mencionadas se trataron en la página \pageref{ejercicio7a}.
Su intersección se caracteriza por la ecuación
$x^{2}+z^{2}=\frac{9}{2}$. Por lo tanto, la porción del cono se parametriza como
${\bf r}\left( u,v\right) =\left\langle u\cos v,u,u\sin v\right\rangle $
y la porción de la esfera se parametriza en la forma
${\bf r}\left(u,v\right) =\left\langle u\cos v,\sqrt{9-u^{2}},u\sin v\right\rangle $, en donde $u\in \lbrack 0,\frac{3}{\sqrt{2}}]$, $v\in \lbrack 0,2\pi ]$.
Calcularemos el flujo en cada región por separado.

\subsubsection{Flujo sobre la superficie del cono}

En este caso ${\bf F}\left( {\bf r}\left( u,v\right) \right) =\left\langle u\cos v,1,u+u\sin v\right\rangle $
y también
\begin{align*}
{\bf r}_{u} &= \left\langle \cos v, 1, \sin v \right\rangle \\
{\bf r}_{v} &=\left\langle -u\sin v,0,u\cos v\right\rangle \\
{\bf r}_{u}\times {\bf r}_{v}&= \left\langle u\cos v,-u,u\sin v\right\rangle\!,
\end{align*}
con lo cual,
${\bf F}\left({\bf r}\left(u,v\right)\right)\cdot \left({\bf r}_{u}\times {\bf r}_{v}\right) = u^{2}\sin v-u+u^{2}$.
Entonces, el flujo de ${\bf F}$ sobre la superficie del cono está dado por

\begin{equation}
\label{flujo9}
\int_{0}^{2\pi}\int_0^{\frac{3}{\sqrt{2}}}\left(u^{2}\sin v-u+u^{2}\right)dudv= \frac{9}{2}\pi \left( \sqrt{2}-1\right)
\end{equation}

\subsubsection{Flujo sobre la superficie de la esfera}

En este sentido,
${\bf F}\left( {\bf r}\left( u,v\right) \right) =\left\langle u\cos v,1,\sqrt{9-u^{2}}+u\sin v\right\rangle \!$, y con la parametrización escogida se obtienen los siguientes vectores:
\begin{align*}
{\bf r}_{u} &= \left\langle \cos v,-\frac{u}{\sqrt{9-u^{2}}},\sin v\right\rangle \\
{\bf r}_{v} &= \left\langle -u\sin v,0,u\cos v\right\rangle \\
{\bf r}_{v}\times {\bf r}_{u}&= \left\langle\frac{u^{2}\cos v}{\sqrt{9-u^{2}}},u,\frac{u^{2}\sin v}{\sqrt{9 -u^{2}}}\right\rangle\!.
\end{align*}
De lo anterior, se obtiene
${\bf F}\left( {\bf r}\left( u,v\right) \right) \cdot \left( {\bf r}_{u}\times {\bf r}_{v}\right) =\frac{u^{3}}{\sqrt{9-u^{2}}}+u+u^{2}\sin v$,
cuya integración con los límites establecidos proporciona el flujo de ${\bf F}$ sobre la porción de la esfera.
\begin{equation}\label{flujo10}
\int_{0}^{2\pi }\int_{0}^{\frac{3}{\sqrt{2}}}\left( \frac{u^{3}}{\sqrt{9-u^{2}}}+u+u^{2}\sin v\right) dudv= \frac{9}{2}\pi \left( 9-5\sqrt{2}\right)
\end{equation}

El flujo sobre toda la superficie es la suma de los valores obtenidos en \ref{flujo9}
y \ref{flujo10}. Como ${\rm div\,}{\bf F}=2$, es claro que este valor corresponde a dos veces el volumen del sólido el cual se obtuvo con la expresión \ref{ejercicio7}.

\item Hallar el flujo neto del campo vectorial ${\bf F}(x,y,z) = \la x,y,2z \ra $
sobre la superficie del sólido determinado por la expresión $\rho =1+\cos \phi$.

\begin{figure}[H]
\centering
\caption{Campo vectorial ${\bf F}(x,y,z)=\la x,y,2z \ra$}
\includegraphics[width=0.4\textwidth]{Img/T_13b.pdf} \hspace{5mm}
\includegraphics[width=0.4\textwidth]{Img/T_13c.pdf}
\fuentepropia
\end{figure}

La figura anterior se obtiene al ejecutar las siguientes instrucciones:

\begin{tcolorbox}[title=\bf Instrucción \verb"Mathematica"]\begin{Verbatim}[formatcom=\letraverbatim]
a=RegionPlot3D[(x^2+y^2+z^2-x)^2<x^2+y^2+z^2,
{x,-0.5,2.2},{y,-1.5,1.5},{z,-1.5,1.5},PlotPoints->90,
Mesh->False,BoxRatios->{1,1,1}];
b=SliceVectorPlot3D[{x,y,2*z},{(x^2+y^2+z^2-x)^2==
x^2+y^2+z^2},{x,-0.5,2.2},{y,-1.5,1.5},{z,-1.5,1.5},
VectorStyle->Black,VectorScale->{0.12,Scaled[0.25],
Automatic},VectorPoints->15,PlotStyle->Opacity[0]];
Show[a,b]
\end{Verbatim}
\end{tcolorbox}

Con la ecuación \(\rho = 1 + \cos \phi\) y las expresiones que definen las coordenadas esféricas, se propone la siguiente función:
\[
\mathbf{r}(\theta,\phi) =
\left\langle
(1 + \cos \phi)\sin \phi \cos \theta,\
(1 + \cos \phi)\sin \phi \sin \theta,\
(1 + \cos \phi)\cos \phi
\right\rangle
\]

la cual parametriza la superficie del sólido con $\phi \in[0,\pi], \theta\in[0,2\pi]$.
De esta función se obtiene
{\small
\begin{eqnarray*}
{\bf r}_{\theta} &=& \left\langle -\left(\cos\phi +1\right)\sin\theta\sin\phi,(\cos \phi +1)\cos\theta\sin\phi,0\right\rangle \\
{\bf r}_{\theta } &=& \big \langle -\left(\cos\phi +1\right)\sin\theta\sin\phi,(\cos \phi +1)\cos\theta\sin\phi, 0 \big\rangle \\
{\bf r}_{\phi } &=& \big\langle \left( \cos \phi +\cos (2\phi) \right)\cos \theta , \left( \cos \phi +\cos (2\phi) \right)\sin \theta,
\left( 2\cos \phi +1\right)\sin \phi\big \rangle, \\
\end{eqnarray*}
}

Además,
{\small
\[{\bf F}({\bf r}(\theta,\phi))=(1+\cos\phi) \left\langle\sin\phi\cos\theta,\sin\phi\sin\theta,2\cos\phi\right\rangle\!,\]

\begin{multline*}
{\bf r}_{\phi}\times {\bf r}_{\theta}=\bigg\langle\left(\cos\phi+1\right)\left(\sin2\phi+\sin\phi\right)\cos\theta\sin \phi,\\
\left( \cos \phi +1\right) \left( \sin 2\phi +\sin \phi \right) \sin \theta \sin \phi ,\left( \cos \phi +1\right)
^{2}\left( 2\cos \phi -1\right) \sin \phi \bigg\rangle .
\end{multline*}
}

Por lo tanto,

\[
\mathbf{F}(\mathbf{r}(\phi,\theta)) \cdot (\mathbf{r}_{u} \times \mathbf{r}_{\phi}) =
\frac{1}{2}(\cos \phi + 1)^{3} \left( \sin 3\phi - \sin 2\phi + 3\sin \phi \right)\!.
\]

Entonces, el flujo de \(\mathbf{F}\) sobre la superficie en mención es:
{\small
\[
\int_{0}^{\pi} \int_{0}^{2\pi}
\frac{1}{2}(\cos \phi + 1)^{3} \left( \sin 3\phi - \sin 2\phi + 3\sin \phi \right)
\, d\theta\, d\phi = \frac{32}{3}\pi.
\]
}

Como ${\rm div\,}{\bf F} =4$, es claro que el flujo corresponde a cuatro veces el volumen de la región, el cual se obtuvo en la fórmula \ref{ejercicio5}.

%\clearpage

\item Hallar el flujo neto del campo vectorial ${\bf F}(x,y,z) = \la x,y,z \ra $
sobre la superficie del sólido acotado por el paraboloide $4y=x^2+z^2$ y la esfera $x^2+y^2+z^2=12$.

Este sólido se consideró en la página \pageref{ejercicio8a}, la intersección de las superficies mencionadas y proyectada en el plano $xz$ se determina por la ecuación $x^2+z^2=8$.

La figura se obtiene mediante la ejecución de las siguientes instrucciones:

\begin{tcolorbox}[title=\bf Instrucción \verb"Mathematica"]\begin{Verbatim}[formatcom=\letraverbatim]
a=RegionPlot3D[4*y>x^2+z^2&&x^2+y^2+z^2<12,{x,-3,3},
{y,0,4},{z,-3,3},PlotPoints->90,Mesh->False];
b=SliceVectorPlot3D[{x,y,z},{x^2+y^2+z^2==12,
4*y==x^2+z^2},{x,-3,3},{y,0,3.7},{z,-3,3},
RegionFunction->Function[{x,y,z},x^2+z^2<=8],
VectorPoints->16,PlotStyle->Opacity[0],
VectorStyle->Black,VectorScale->{0.15,Scaled[0.25],
Automatic}];
Show[a,b]
\end{Verbatim}
\end{tcolorbox}

\begin{figure}[H]
\centering
\caption{Campo vectorial ${\bf F}(x,y,z)=\la x,y,z\ra$ sobre $4y=x^2+z^2$ \& $x^2+y^2+z^2=12$}
\includegraphics[width=0.4\textwidth]{Img/T_15c.pdf} \hspace{5mm}
\includegraphics[width=0.4\textwidth]{Img/T_15d.pdf}
\fuentepropia
\end{figure}

Con el propósito de parametrizar la superficie mencionada antes, es posible asignar
$x=\sqrt{8}u\cos v$, $z=\sqrt{8}u\sin v$, en donde $u\in [0,1]$ y $v \in [0,2\pi]$.
Es por esta razón que al sustituir en las ecuaciones respectivas utilizaremos las funciones
${\bf r}(u,v) = \langle\sqrt{8}u \cos v ,2u^2 ,\sqrt{8}u\sin v \rangle $, como parametrización del paraboloide y
${\bf r}(u,v) = \langle  \sqrt{8}u \cos v ,\sqrt{12-8u^2} ,\sqrt{8}u\sin v \rangle$, en el caso de la esfera.
Los cálculos se harán por separado.

\subsubsection{Flujo sobre la superficie de la esfera}

En este caso, con la parametrización escogida se obtiene
\begin{eqnarray*}
{\bf r}_u &=& \left \langle 2 \sqrt{2}\cos(v), -\frac{4u}{\sqrt{3-2u^2}},2\sqrt{2}\sin(v) \right \rangle \\
{\bf r}_v &=&\langle -2\sqrt{2}u\sin(v),0,2\sqrt{2}u\cos(v) \rangle \\
{\bf r}_v \times {\bf r}_u &=& \left \langle \frac{8\sqrt{2}u^2 \cos(v)}{\sqrt{3-2u^2}},8u,\frac{8\sqrt{2}u^2 \sin(v)}{\sqrt{3-2u^2}} \right\rangle\!.
\end{eqnarray*}

Por lo tanto, ${\bf F}({\bf r}(u,v)) = \langle \sqrt{8}u \cos v ,\sqrt{12-8u^2} ,\sqrt{8}u\sin v \rangle$ y
${\bf F}({\bf r}(u,v))\cdot({\bf r}_v \times {\bf r}_u)=\dfrac{48u}{\sqrt{3-2u^2}}$.
El flujo de ${\bf F}$ sobre la esfera está dado por
\begin{equation}\label{flujo12}
\int_0^1\int_0^{2\pi}\frac{48u}{\sqrt{3-2u^2}} dvdu= 48\pi (\sqrt{3}-1)
\end{equation}

Este resultado también se obtiene ejecutando  las siguientes instrucciones:

\begin{tcolorbox}[title=\bf Instrucción \verb"Mathematica"]\begin{Verbatim}[formatcom=\letraverbatim]
r={x,y,z}={Sqrt[8]*u*Cos[v],Sqrt[12-8*u^2],
Sqrt[8]*u*Sin[v]};
A=Simplify[Dot[{x,y,z},Cross[D[r,v],D[r,u]]]]
Integrate[A,{u,0,1},{v,0,2*Pi}]
\end{Verbatim}
\end{tcolorbox}

\subsubsection{Flujo sobre la superficie del paraboloide}

Con la función que parametriza el paraboloide, se obtiene

\begin{eqnarray*}
{\bf r}_u &=& \left \langle 2 \sqrt{2}\cos(v), 4u,2\sqrt{2}\sin(v) \right \rangle \\
{\bf r}_v &=&\langle -2\sqrt{2}u\sin(v),0,2\sqrt{2}u\cos(v) \rangle \\
{\bf r}_u \times {\bf r}_v &=& \left \langle 2\sqrt{2}u^2 \cos(v),-8u,2\sqrt{2}u^2 \sin(v) \right\rangle
\end{eqnarray*}

Luego, ${\bf F}({\bf r}(u,v)) = \langle \sqrt{8}u \cos v ,2u^2 ,\sqrt{8}u\sin v \rangle$ y
${\bf F}({\bf r}(u,v))\cdot({\bf r}_u \times {\bf r}_v)=16u^3 $.
El flujo de ${\bf F}$ sobre el paraboloide está determinado por
\begin{equation}\label{flujo11}
\int_0^1\int_0^{2\pi } 16 u^3 dvdu= 8\pi.
\end{equation}
El cálculo anterior se obtiene ejecutando las siguientes instrucciones:

\begin{tcolorbox}[title=\bf Instrucción \verb"Mathematica"]\begin{Verbatim}[formatcom=\letraverbatim]
r={x,y,z}={Sqrt[8]*u*Cos[v],2*u^2,Sqrt[8]*u*Sin[v]};
Integrate[Dot[{x,y,z},Cross[D[r,u],D[r,v]]],{u,0,1},
{v,0,2*Pi}]
\end{Verbatim}
\end{tcolorbox}

El flujo neto sobre la superficie del sólido es la suma de los valores obtenidos en \ref{flujo11}
y \ref{flujo12}. En este caso ${\rm div\,}{\bf F} =3$, el valor obtenido es tres veces el volumen obtenido en la expresión \ref{ejercicio8}.

\item Hallar el flujo neto del campo vectorial ${\bf F}(x,y,z)=\la x-2,z,-y \ra $
sobre la superficie del sólido acotado por el semicono
$x=\sqrt{y^2+z^2}$ y la esfera con ecuación $y^2+z^2+4x=12$.

\begin{figure}[H]
\centering
\caption{Campo vectorial ${\bf F}(x,y,z)=\la x-2,z,-y \ra$ sobre una superficie}
\includegraphics[width=0.4\textwidth]{Img/T_17d.pdf} \hspace{5mm}
\includegraphics[width=0.4\textwidth]{Img/T_17e.pdf}

\includegraphics[width=0.4\textwidth]{Img/T_17c.pdf}
\fuentepropia
\end{figure}

Las figuras se obtienen al ejecutar las siguientes instrucciones:

\begin{tcolorbox}[title=\bf Instrucción \verb"Mathematica"]\begin{Verbatim}[formatcom=\letraverbatim]
a=RegionPlot3D[x>Sqrt[y^2+z^2]&&y^2+z^2+4*x<12,
{x,0,4},{y,-3,3},{z,-3,3},PlotPoints->90,
Mesh->False];
b=SliceVectorPlot3D[{x-2,z,-y},{x==Sqrt[y^2+z^2],
y^2+z^2+4*x==12},{x,0,3.5},{y,-2,2},{z,-2,2},
VectorPoints->18,RegionFunction->Function[{x,y,z},
y^2+z^2<=4],VectorStyle->Black,VectorScale->{0.16,
Scaled[0.24],Automatic},PlotStyle->Opacity[0]];
Show[a,b]
\end{Verbatim}
\end{tcolorbox}

La intersección de estas dos superficies se consigue eliminando $y^{2}+z^{2}$ de las dos ecuaciones, lo que se reduce a $x^{2}+4x=12$ e implica que $x=2$.
La proyección en el plano $yz$ se determina por la ecuación $y^{2}+z^{2}=4$. Esta región con su interior se describe asignando $y=u\cos v$, $z=u\sin v$,
en donde $u\in \lbrack 0,2]$, $v\in \lbrack 0,2\pi ]$.
Sobre el cono se satis\-face que $x=\sqrt{y^{2}+z^{2}}=u$ y sobre el paraboloide se cumple que
$x=\frac{1}{4}\left( 12-y^{2}-z^{2}\right) =\frac{1}{4}\left(12-u^{2}\right) = 3-\frac{1}{4}u^{2}$,

Por lo descrito anteriormente, resulta apropiado escoger como parametrizaciones de la porción del paraboloide y del cono a las siguientes funciones, respectivamente.

\begin{eqnarray*}
{\bf r}\left( u,v\right) &=&\left\langle 3-\frac{1}{4}u^{2},\cos v,u\sin v\right\rangle\\
{\bf r}( u,v) &=&\left\langle u,\cos v,u\sin v\right\rangle\!.
\end{eqnarray*}
con $u\in [0,2]$, $v\in [0,2\pi ]$.
Calcularemos el flujo sobre cada una de las regiones componentes.
\subsubsection{Flujo sobre la porción del cono}

Es claro que \[{\bf F}({\bf r}(u,v))=\left\langle 1-\frac{1}{4}u^{2},u\sin v, -\cos v\right\rangle\] y con la parametrización escogida se sigue que
\begin{eqnarray*}
{\bf r}_{u}&=&\left\langle 1,0,\sin v\right\rangle \\
{\bf r}_{v}&=&\left\langle 0,-\sin v,u\cos v\right\rangle \\
{\bf r}_{u}\times {\bf r}_{v}&=& \left\langle \sin ^{2}v,-u\cos v,-\sin v\right\rangle\!.
\end{eqnarray*}
Entonces, ${\bf F}({\bf r}(u,v)) \cdot ({\bf r}_{u}\times {\bf r}_{v}) = u(2-u)$.
Con lo cual, el flujo está dado por
\begin{equation}\label{flujo15}
\int_0^2\int_0^{2\pi} u(2-u) du dv = \frac{8\pi}{3}
\end{equation}

\subsubsection{Flujo sobre la porción del paraboloide}

Con la parametrización escogida se sigue que
\[{\bf F}({\bf r}(u,v))=\left\langle u-2,u\sin v,-\cos v \right\rangle\] y de la función ${\bf r}$,
se obtienen los siguientes vectores.
\begin{eqnarray*}
{\bf r}_{u}&=&\left\langle -\frac{1}{2}u,0,\sin v\right\rangle \\
{\bf r}_{v}&=&\left\langle 0,-\sin v,u\cos v\right\rangle \\
{\bf r}_{u}\times {\bf r}_{v}&=&\left\langle \sin ^{2}v,\frac{1}{2}u^{2}\cos v,\frac{1}{2}u\sin v\right\rangle\!.
\end{eqnarray*}
Entonces, ${\bf F}({\bf r}(u,v)) \cdot ({\bf r}_{u}\times {\bf r}_{v}) = \frac{1}{4}u(u^2-4)$, por lo tanto, el flujo está dado por
\begin{equation}\label{flujo16}
\int_0^2\int_0^{2\pi}\frac{1}{4}\left (u-\frac{u^3}{4} \right ) du dv = 2\pi
\end{equation}
Sumando los resultados obtenidos en \ref{flujo15} y \ref{flujo16}, se obtiene el flujo neto sobre el sólido que se considera. Este valor coincide con el volumen hallado en la página \pageref{ejercicio10}, pues en este caso ${\rm div}\,{\bf F}=1$.

\item Hallar el flujo neto del campo vectorial ${\bf F}(x,y,z) = \left\langle yz,y-xz,xy\right\rangle $
sobre la superficie del sólido acotado por el paraboloide $x+3=y^2+z^2$
y el paraboloide hiperbólico $y^2-z^2=9x-3$.

La figura de esta región se produce con las siguientes instrucciones:

\begin{tcolorbox}[title=\bf Instrucción \verb"Mathematica"]\begin{Verbatim}[formatcom=\letraverbatim]
a=RegionPlot3D[x+3>y^2+z^2&&y^2-z^2>9*x-3,{x,-3,2},
{y,-2,2},{z,-2,2},PlotPoints->90,Mesh->False,
BoxRatios->{1,1,1}];
b=SliceVectorPlot3D[{y*z,y-x*z,x*y},{y^2-z^2==9*x-3,
x+3==y^2+z^2},{x,-3,1.2},{y,-2,2},{z,-2,2},
VectorPoints->20,VectorStyle->Black,
RegionFunction->Function[{x,y,z},4*y^2+5*z^2<=15],
VectorScale->{0.18,Scaled[0.3],Automatic},
PlotStyle->Opacity[0]];
Show[a,b]
\end{Verbatim}
\end{tcolorbox}

\begin{figure}[H]
\centering
\caption{${\bf F}(x,y,z)=\la yz,y-xz,xy\ra $ sobre $x+3=y^2+z^2$ \& $y^2-z^2=9x-3$}
\includegraphics[width=0.4\textwidth]{Img/T_16d.pdf} \hspace{5mm}
\includegraphics[width=0.4\textwidth]{Img/T_16e.pdf}

\includegraphics[width=0.4\textwidth]{Img/T_16f.pdf}
\fuentepropia
\end{figure}

El paraboloide y el paraboloide hiperbólico se intersecan en una curva cuya proyección en el plano $yz$, está dada por
$4y^{2}+5z^{2}=15$, esta ecuación representa una elipse.
Haciendo $y=\frac{\sqrt{15}}{2}u\cos v$, $z=\sqrt{3}u\sin v$ se des\-cribe esta región. También se cumplen las siguientes relaciones para el paraboloide.

\begin{eqnarray*}
x&=&\left( \frac{\sqrt{15}}{2}u\cos v\right)^{2}+\left(\sqrt{3}u\sin v\right)^{2}-3 =\frac{3}{8}u^{2}\cos 2v
+\frac{27}{8}u^{2}-3
\end{eqnarray*}
Asimismo el paraboloide hiperbólico
{\small
\begin{eqnarray*}
x&=&\frac{1}{9}\left( \left( \frac{\sqrt{15}}{2}u\cos v\right) ^{2}-\left( \sqrt{3}u\sin v\right) ^{2}+3\right)
= \frac{3}{8}u^{2}\cos 2v+\frac{1}{24}u^{2}+\frac{1}{3}.
\end{eqnarray*}
}

El flujo neto se hará por separado sobre las dos regiones que acotan el sólido.

\subsubsection{Flujo sobre la porción del paraboloide}
Por lo mencionado antes, resulta natural escoger la siguiente parametri\-zación para el paraboloide.

{\small
\[{\bf r}\left( u,v\right) = \langle \frac{3}{8}u^{2}\cos 2v+\frac{27}{8}u^{2}-3, \sqrt{\frac{15}{4}}u\cos v,\sqrt{3}u\sin v \rangle,\]}

con $u \in [0,1]$, $v \in [0,2\pi]$. De esta función se obtienen los siguientes vectores:

{\small
\begin{eqnarray*}
{\bf r}_{u}&=&\left\langle \frac{27}{4}u+\frac{3}{4}u\cos 2v,\frac{1}{2}\sqrt{15}\cos v,\sqrt{3}\sin v\right\rangle \\
{\bf r}_{v}&=&\left\langle -\frac{3}{4}u^{2}\sin 2v,-\frac{1}{2}\sqrt{15}u\sin v,\sqrt{3}u\cos v\right\rangle \\
{\bf r}_{u}\times {\bf r}_{v}&=&\left\langle -\frac{3}{2}\sqrt{5}u,\frac{15}{2}\sqrt{3}u^{2}\cos v, 3\sqrt{3}\sqrt{5}u^{2}\sin v\right\rangle
\end{eqnarray*}
}

Asimismo
\begin{multline*}
{\bf F}({\bf r}(u,v))=\bigg \langle \frac{3}{4}\sqrt{5}u^2\sin(2v),
\frac{1}{16}\sqrt{3}u(8\sqrt{5}\cos(v)+(48-51u^2)\sin(v) \\
3u^2\sin(3v)), \frac{3}{16}\sqrt{15}u\cos(v)(u^2(9+\cos(2v))-8)\bigg \rangle,
\end{multline*}
con lo cual
${\bf F}({\bf r}(u,v)) \cdot \left({\bf r}_{u}\times {\bf r}_{v}\right)=\frac{45}{4} u^3\cos(v)(\sqrt{5}\cos(v)-\sin(v))$.
Por lo tanto, el flujo sobre la porción del paraboloide que se considera es:
\begin{equation} \label{flujo13}
\int_{0}^{1}\int_{0}^{2\pi}\frac{45}{4} u^3\cos(v)(\sqrt{5}\cos(v)-\sin(v)) dvdu = \frac{45 \sqrt{5}}{16} \pi
\end{equation}
El cálculo anterior se puede hacer en \verb"Mathematica" ejecutando las siguien\-tes instrucciones:

\begin{tcolorbox}[title=\bf Instrucción \verb"Mathematica"]\begin{Verbatim}[formatcom=\letraverbatim]
{y,z,x}={Sqrt[15]/2*u*Cos[v],Sqrt[3]*u*Sin[v],
y^2+z^2-3};
F=TrigFactor[{y*z,y-x*z,x*y}]
Integrate[Dot[F,Cross[D[{x,y,z},v],D[{x,y,z},u]]],
{u,0,1},{v,0,2*Pi}]
\end{Verbatim}
\end{tcolorbox}

\subsubsection{Flujo sobre el paraboloide hiperbólico}

En este caso, la parametrización escogida es la siguiente:
\[
{\bf r}(u,v) =
\left\langle
\frac{3}{8}u^{2}\cos 2v + \frac{1}{24}u^{2} + \frac{1}{3},\;
\frac{\sqrt{15}}{2}u\cos v,\;
\sqrt{3}u\sin v
\right\rangle\!,
\]
en donde \( u \in [0,1],\; v \in [0,2\pi] \). También se puede ver que
\[
\mathbf{F}(\mathbf{r}(u,v)) =
\left\langle
\frac{3}{2}\sqrt{5}\,u,\
-\frac{5u^2\cos(v)}{2\sqrt{3}},\
\sqrt{\frac{5}{3}}\, u^2 \sin(v)
\right\rangle\!.
\]
Con la función ${\bf r}$ obtenemos los siguientes vectores:
\begin{eqnarray*}
{\bf r}_{u} &=&\left\langle \frac{1}{12}u+\frac{3}{4}u\cos 2v,\frac{1}{2}\sqrt{15}\cos v,\sqrt{3}\sin v\right\rangle\\
{\bf r}_{v}&=&\left\langle -\frac{3}{4}u^{2}\sin 2v,-\frac{1}{2}\sqrt{15}u\sin v,\sqrt{3}u\cos v\right\rangle \\
{\bf r}_{u}\times {\bf r}_{v}&=& \la \frac{3}{2}\sqrt{5}u,-\frac{5}{6}\sqrt{3} u^{2}\cos v,\frac{1}{3}\sqrt{15}u^{2}\sin v\ra,
\end{eqnarray*}
entonces,
{\small
\[{\bf F}({\bf r}(u,v)) \cdot \left({\bf r}_{u}\times {\bf r}_{v}\right)=\frac{5u^3}{24}\cos(v)(6\sqrt{5}\cos(v)-(62+u^2(1+9\cos(2v)))\sin(v)).\]
}
Por lo tanto, sobre la porción del paraboloide hiperbólico el flujo de ${\bf F}$ es
\begin{equation} \label{flujo14}
\int_0^1\int_0^{2\pi}{\bf F}({\bf r}(u,v)) \cdot \left({\bf r}_{u}\times{\bf r}_{v}\right)dvdu=-\frac{5 \sqrt{5}}{16}\pi
\end{equation}
Con la ayuda de las siguientes instrucciones, se puede obtener el cálculo anterior.

\begin{tcolorbox}[title=\bf Instrucción \verb"Mathematica"]\begin{Verbatim}[formatcom=\letraverbatim]
{y,z,x}={Sqrt[15]/2*u*Cos[v],Sqrt[3]*u*Sin[v],
(y^2-z^2+3)/9};
Integrate[Dot[{y*z,y-x*z,x*y},Cross[D[{x,y,z},u],
D[{x,y,z},v]]],{u,0,1},{v,0,2*Pi}]
\end{Verbatim}
\end{tcolorbox}

El flujo neto sobre la superficie del sólido es la suma de los valores obtenidos en \ref{flujo13} y
\ref{flujo14}, es decir $\frac{5 \sqrt{5}}{16} \pi$ . Como ${\rm div\,}{\bf F}=1$,
este valor coincide con el volumen del sólido hallado con la expresión \ref{ejercicio9}.
\end{parrafonumerado}

\begin{ejer}
Resuelva los siguientes problemas.
\begin{parrafonumerado}

\item Calcule el flujo del campo vectorial ${\bf F}(x,y,z)= \la x-2y,3x+y,z+1\ra $ sobre la superficie del hiperboloide de un manto $x^2 + y^2 - z^2 = 1$, con $-2 \leq z \leq 2$.

\begin{figure}[H]
\centering
\caption{${\bf F}(x,y,z)=\la x-2y,3x+y,z+1\ra$ sobre un hiperboloide}
\includegraphics[width=0.4\textwidth]{Img/T_35a.pdf} \hspace{5mm}
\includegraphics[width=0.4\textwidth]{Img/T_35b.pdf}
\end{figure}

\item Considere la superficie sobre el paraboloide  $z=x^2-2x+y^2$ interior al cilindro cardioidico $r=1+\cos t$.

\begin{figure}[H]
\centering
\caption{Superficie sobre un paraboloide}
\includegraphics[width=0.4\textwidth]{Img/T_36a.pdf} \hspace{5mm}
\includegraphics[width=0.4\textwidth]{Img/T_36b.pdf}
\end{figure}

Calcule el flujo del campo vectorial ${\bf F}(x,y,z)= \la x-2,y,2z\ra $ sobre la superficie del paraboloide.
$z=x^2-2x+y^2$, interior al cilindro cardioidico $r=1+\cos t$.

%

\begin{figure}[H]
\centering
\caption{Campo vectorial ${\bf F}(x,y,z)= \la x-2,y,2z\ra $ sobre un paraboloide}
\includegraphics[width=0.4\textwidth]{Img/T_36c.pdf} \hspace{5mm}
\includegraphics[width=0.4\textwidth]{Img/T_36d.pdf}
\end{figure}

Evalúe $\ds \int_S\rot{F}\cdot d{\bf S}$ y $\ds\int_C{\bf F}\cdot d{\bf r}$; en donde $S$
es la superficie en mención y $C$ es su frontera.

\begin{comment}
x=Cos[t]*(1+Cos[t]);
y=Sin[t]*(1+Cos[t]);
z=x^2-2*x+y^2;
a=ParametricPlot3D[{x,y,z},{t,0,2*Pi},PlotStyle->Red];

Show[SliceVectorPlot3D[{x-2,y,2z},{z==x^2-2*x+y^2},{x,-1/4,2},{y,-3/2,3/2},
{z,-1,1.1},VectorStyle->Black,RegionFunction->Function[{x,y,z},
(x^2+y^2-x)^2<=x^2+y^2],VectorScale->{0.1,Scaled[0.2],Automatic},VectorPoints->17],a]
\end{comment}

\item Calcule el flujo del campo vectorial ${\bf F}(x,y,z)= \la z-y,y,z^2 \ra $ sobre la superficie del paraboloide hiperbólico $9z=x^2-y^2$, interior al cilindro $x^2+y^2 \leq 16$ y exterior al cilindro cardioidico $r=1+\cos t$.

\begin{figure}[H]
\centering
\caption{Campo vectorial ${\bf F}(x,y,z)= \la z-y,y,z^2\ra $ sobre una silla de montar}
\includegraphics[width=0.4\textwidth]{Img/T_37a.pdf} \hspace{5mm}
\includegraphics[width=0.4\textwidth]{Img/T_37b.pdf}
\end{figure}

\begin{comment}
x=Cos[t]*(1+Cos[t]);
y=Sin[t]*(1+Cos[t]);
z=(x^2-y^2)/9;
a=ParametricPlot3D[{x,y,z},{t,0,2*Pi},
PlotStyle->{Thickness->0.005,Blend[{Blue,Black},1/3]}];
b=ParametricPlot3D[{4*Cos[t],4*Sin[t],16*Cos[2*t]/9},{t,0,
2*Pi},PlotStyle->{Thickness->0.005,
Blend[{Blue,Black},1/3]}];
datos=Table[{z-y,y,z^2},{x,-4,4},{y,-4,4},{z,-2,2}];
Show[ListSliceVectorPlot3D[datos,
ImplicitRegion[
9*z==x^2-y^2&&(x^2+y^2-x)^2>=x^2+y^2&&
x^2+y^2<=16,{x,y,z}],
DataRange->{{-4,4},{-4,4},{-2,2}},VectorPoints->15,
VectorStyle->{Thin,Black,Thick},
BoxRatios->{1,1,2/3}],a,b,PlotPoints->60]
\end{comment}

\end{parrafonumerado}
\end{ejer}